% ============================================================
%  R. Nielsen, "Den speculative Methodes Anvendelse paa den hellige Historie"
%  [The Application of the Speculative Method to Sacred History]
%  Paa Dansk udgivet af B. C. Bøggild. Kjøbenhavn: H. C. Klein, 1842. 162 pp.
%  TRANSLATION (English)
%  Source of truth: transcription.tex (Danish), COMPLETE; sampled clean
%  2026-09-26 (12 pp., 0.08 errors/page).
%  Mirrors transcription.tex 1:1; every printed page carries \opage{N} at the
%  corresponding point in the English, with its "% --- p. N ---" line.
%  Translated 26 Sep 2026 in ten batches of about 16 printed pages
%  (TRANSLATION-PLAYBOOK.md). German, Latin and Greek quotations are kept as
%  printed, German misprints included; printer's misprints in the Danish are
%  translated by their intended sense; the book's Rettelser, applied in the
%  Danish, are applied here too. Footnote marks restart on every printed page,
%  as in the print. The printed Contents and the section heads disagree in six
%  places (§§ 3, 7, 10, 13, 15, 21); each is translated as printed.
%  Terms: Empirie = empirical knowledge; Begreb = concept; Forestilling =
%  representation; Erkjendelse = cognition; Tilværelse = existence;
%  Virkelighed = actuality; Anskuen = intuiting; Forsoning = atonement
%  (theological) / reconciliation (of faith and knowledge).
% ============================================================
\documentclass[12pt]{book}
\usepackage[utf8]{inputenc}
\usepackage[T1]{fontenc}
\usepackage[english]{babel}
\usepackage[protrusion=true,expansion=true]{microtype}
\usepackage[margin=1.4in]{geometry}
\usepackage{setspace}
\usepackage{amsmath}
\usepackage{enumitem}
\usepackage{graphicx}    % for \reflectbox (the old Danish »ɔ:« id-est mark)
\DeclareUnicodeCharacter{0254}{\reflectbox{c}}  % ɔ (U+0254) = reversed c
\usepackage{libertinus}
\usepackage{libertinust1math}
\usepackage{textalpha}   % polytonic Greek typed directly
\usepackage{fancyhdr}
\usepackage[colorlinks=true,linkcolor=black,urlcolor=black,
  unicode=true,plainpages=false,pdfpagelabels,hyperfootnotes=false]{hyperref}
\setlength{\headheight}{15pt}
\pagestyle{fancy}
\fancyhf{}
\fancyhead[LE,RO]{\thepage}
\fancyhead[RE]{\textit{The Application of the Speculative Method}}
\fancyhead[LO]{\textit{\leftmark}}
\setlength{\parindent}{1.5em}
\setlength{\parskip}{0pt}
\setstretch{1.3}

\renewcommand{\thefootnote}{\arabic{footnote})}

\newcommand{\deel}[2]{%
  \par\vspace{2.2\baselineskip}
  \begin{center}
    {\Large\bfseries #1}\\[0.7\baselineskip]
    {\large\bfseries\emph{#2}}
  \end{center}
  \vspace{0.9\baselineskip}
  \markboth{#1}{#1}%
  \phantomsection
  \addcontentsline{toc}{chapter}{\texorpdfstring{#1\quad #2}{#1 #2}}%
  \par\nopagebreak}

\newcommand{\parag}[2]{%
  \par\vspace{1.6\baselineskip}
  \begin{center}
    {\bfseries §~#1.}\\[0.45\baselineskip]
    {\emph{#2}}
  \end{center}
  \vspace{0.5\baselineskip}
  \markright{§ #1. #2}%
  \phantomsection
  \addcontentsline{toc}{section}{\texorpdfstring{§~#1.\quad #2}{§ #1. #2}}%
  \par\nopagebreak}

\newcommand{\division}[1]{%
  \par\vspace{2.0\baselineskip}
  \begin{center}{\large\bfseries #1}\end{center}
  \vspace{0.7\baselineskip}
  \markboth{#1}{#1}%
  \phantomsection
  \addcontentsline{toc}{chapter}{#1}%
  \par\nopagebreak}

\newcommand{\dsi}{\reflectbox{c}:}

\title{\textbf{The Application of the Speculative Method\\ to Sacred History}}
\author{R.\ Nielsen\\[4pt]
  \normalsize published in Danish by B.\ C.\ Bøggild\\[4pt]
  \small translated from the Danish}
\date{Copenhagen: H.\ C.\ Klein, 1842}

\usepackage{marginnote}
\newcommand{\opage}[1]{\marginnote{\footnotesize\textit{[#1]}}}
\newcommand{\apage}[1]{\marginnote{\footnotesize\textit{[#1]}}}

\begin{document}
\hypersetup{
  pdftitle={The Application of the Speculative Method to Sacred History},
  pdfauthor={R. Nielsen}}
\frontmatter
\maketitle
\thispagestyle{empty}

\tableofcontents

\cleardoublepage
\thispagestyle{empty}
\phantomsection
\addcontentsline{toc}{chapter}{Title Page}
\begin{center}
  {\Large The Application}\\[1.3\baselineskip]
  {\large of the Speculative Method to}\\[1.3\baselineskip]
  {\huge\bfseries Sacred History.}\\[1.1\baselineskip]
  \rule{4em}{0.4pt}\\[1.8\baselineskip]
  By\\[0.9\baselineskip]
  {\Large R.\ Nielsen,}\\[0.4\baselineskip]
  {\small Licentiate in Theology, Professor of Philosophy\\
   at the University of Copenhagen.}\\[1.3\baselineskip]
  \rule{9em}{0.4pt}\\[1.5\baselineskip]
  Published in Danish\\[0.5\baselineskip]
  by\\[0.5\baselineskip]
  {\Large B.\ C.\ Bøggild,}\\[0.35\baselineskip]
  {\small\textit{Cand. theol.}}\\[2.6\baselineskip]
  \rule{13em}{0.4pt}\\[1.6\baselineskip]
  {\large Copenhagen.}\\[0.5\baselineskip]
  {\small Published by H.\ C.\ Klein, Vimmelskaftet No.\ \textit{23}.\\
   Printed at the Seidelin Press, by Louis Klein.\\[0.3\baselineskip]
   \textit{1842}.}
\end{center}

\chapter*{Translator's Preface}
\phantomsection
\addcontentsline{toc}{chapter}{Translator's Preface}

\textit{[The beginning of the Translator's Preface is missing in the copy
used: a whole leaf --- two pages --- is not included in the scan, and the
Preface therefore begins in the middle of a sentence.]}

\noindent meet us in this work, will surely lead the reader to overlook that
deficiency. If the application of the speculative method to theology, in
general as in a single sphere, is to be defended in truth, this must first
and foremost be done through a thorough opposition, conducted from within,
against the middle-way system of critical, moderate rationalism; and this is
precisely a chief merit of this work. And while \textit{Dr.} Martensen, in his
splendid work ``The Autonomy of Human Self-Consciousness'', has from another
side carried through a thorough opposition against subjective rationalism
and, by proving that Schleiermacher too was a subjective rationalist, has also
extended his opposition to Schleiermacher's system, but, owing to the
particular economy of his treatise, only briefly touched on objective
rationalism, Prof.\ Nielsen in his second Part, ``Sacred History'', has
carried the dialectical polemic on into the territory of objective
rationalism (of panlogism, of Straußianism), and, in the same way as with
subjective rationalism, has endeavored from within to show the logical
insufficiency of this system. He has thus opened the struggle between
``speculative orthodoxy'' and panlogism, which necessarily had to arise among
us, just as in Germany, and has carried it onto the only proper ground for
such a struggle, namely science.

I take the liberty of referring those who may wish to have a complete survey
of the whole train of thought of the work to the review of it that is found in
``Fædrelandet'' No.\ \textit{329}, \textit{1840}.

As regards the translation itself, I feel full well how much, particularly with
respect to purity of style and to uniformity of spelling, I stand in need of
the lenient judgment of the kind reader; might I only hope, by a true rendering
of the author's thoughts and of his light, lively style, to reconcile the
reader to those deficiencies!

\bigskip
\begin{center}
Copenhagen, the \textit{24}th of January \textit{1842}.
\end{center}

\nopagebreak
\begin{flushright}
{\large B.\ C.\ Bøggild.}\hspace*{3em}
\end{flushright}

\chapter*{Contents}
\phantomsection
\addcontentsline{toc}{chapter}{Contents}

\begingroup
\parindent=0pt
\parfillskip=0pt

\begin{center}
  {\large\bfseries First Part.}\\[0.5\baselineskip]
  \emph{The Relation between Philosophy and History, Considered in
  General.}
\end{center}
\medskip

{\raggedleft Page.\par}
\smallskip

§~\textit{1.}\quad The Relation between the A Posteriori and the A Priori\dotfill\textit{1.}\par
§~\textit{2.}\quad Subjective Faith and Speculative Logic\dotfill\textit{5.}\par
§~\textit{3.}\quad Vulgar Empirical Knowledge\dotfill\textit{11.}\par
§~\textit{4.}\quad Transition to Speculative Empirical Knowledge\dotfill\textit{18.}\par
§~\textit{5.}\quad The Relation between Vulgar and Speculative Empirical Knowledge\dotfill\textit{25.}\par
§~\textit{6.}\quad The Concept of Science\dotfill\textit{30.}\par
§~\textit{7.}\quad The Speculative Method\dotfill\textit{36.}\par
§~\textit{8.}\quad Conceiving and Intuiting\dotfill\textit{42.}\par
§~\textit{9.}\quad Our Knowledge Is Relative\dotfill\textit{47.}\par
§~\textit{10.}\quad The Relation between Dogmatics and Sacred History\dotfill\textit{51.}\par

\bigskip
\begin{center}
  {\large\bfseries Second Part.}\\[0.5\baselineskip]
  {\large Sacred History.}
\end{center}
\medskip

§~\textit{11.}\quad Catholicism\dotfill\textit{53.}\par
§~\textit{12.}\quad The Fundamental Contradiction of Catholicism\dotfill\textit{56.}\par
§~\textit{13.}\quad Transition from Catholicism to Protestantism\dotfill\textit{60.}\par
§~\textit{14.}\quad Faith and Holy Scripture\dotfill\textit{63.}\par
§~\textit{15.}\quad The Symbols. Lutheran Scholasticism\dotfill\textit{71.}\par
§~\textit{16.}\quad Continuation. Transition to Criticism\dotfill\textit{82.}\par
\hangindent=3.5em\hangafter=1
§~\textit{17.}\quad Subjective Rationalism. The Profanation of Sacred
  History\dotfill\textit{87.}\par
§~\textit{18.}\quad Faith and Knowledge in Critical Theology\dotfill\textit{96.}\par
§~\textit{19.}\quad Examination of Critical Theology\dotfill\textit{104.}\par
§~\textit{20.}\quad Continuation. Transition to Panlogism\dotfill\textit{119.}\par
§~\textit{21.}\quad The Metaphysical Testing of Panlogism\dotfill\textit{135.}\par
§~\textit{22.}\quad Return to the Church. Speculative Theology\dotfill\textit{143.}\par
§~\textit{23.}\quad Miracle and Sacred History\dotfill\textit{150.}\par
\hspace*{3.2em}Conclusion\dotfill\textit{161.}\par

\endgroup

\chapter*{Errata}
\phantomsection
\addcontentsline{toc}{chapter}{Errata}

\begin{center}
{\large\bfseries Errata,}\\[0.4\baselineskip]
which the reader is asked to correct before reading the book.
\end{center}
\medskip

\begin{center}
\begin{tabular}{@{}l@{~}r@{~}l@{~}l@{}}
Page & \textit{4}   & line  & \textit{17}: naar den read: naar det\\
---  & \textit{4}   & ---   & \textit{26}: de maa read de maae\\
---  & \textit{7}   & ---   & \textit{11}: Subfaction read Subtraction\\
---  & \textit{10}  & ---   & \textit{3} from bottom: almindeligt, read almindelige\\
---  & \textit{12}  & ---   & \textit{1} from bottom: gaae ind read slaae ind\\
---  & \textit{15}  & ---   & \textit{4} from bottom: under et nyt read under Skikkelse af et nyt\\
---  & \textit{41}  & ---   & \textit{1}: næsten read næste\\
---  & \textit{42}  & ---   & \textit{13}: fast read faste\\
---  & \textit{48}  & ---   & \textit{6}: Virken read Viden\\
---  & \textit{71}  & ---   & \textit{12}: \emph{Scholastisme} read \emph{Scholasticisme}\\
---  & \textit{89}  & ---   & \textit{5}: κεκαυτερισμένυος read κεκαυτερισμένους\\
---  & \textit{132} & Note. & \textit{1}: das Mystische read das Mythische
\end{tabular}
\end{center}

\mainmatter

% --- p. 1 ---
\opage{1} %
\deel{First Part.}{The Relation between Philosophy and History, Considered in General.}

\parag{1}{The Relation between the A Posteriori and the A Priori.}

\par\vspace{0.35\baselineskip}\centerline{\rule{5em}{0.4pt}}\vspace{0.35\baselineskip}

The old saying ``\textit{medium tenuere beati!}'' is heard not only where
drinking is concerned, but also where studying is concerned; and so everyday
understanding often recommends this rule when it speaks of the relation
between history and philosophy. ``One must praise history, it is said, for all
thinking would be empty and nebulous if the things themselves were not traced
out with the greatest diligence and accuracy; --- but philosophy is no less
necessary, for the historian needs logical rules, lest the treasury of
knowledge become a formless mass, a \textit{rudis indigestaque moles}.
Prudence therefore bids that one lean too far neither to the one side nor to
the other, so as not to confound the a priori and the a posteriori with each
other.''

But however plausible and reassuring this well-meant and friendly admonition
to sobriety and moderation may seem at first glance, on closer examination of
the matter it nevertheless proves empty and wavering. For if
% --- p. 2 ---
\opage{2} thinking is to be adjusted to experience, and experience conversely
corrected by thinking, what then is the touchstone of truth, where then is the
supreme judge?

Experience in and for itself cannot occupy the judge's seat; for in order to
prove its title to be the judge of truth, it would have to adduce grounds drawn
from thinking; it would thus, precisely in striving for the office of judge,
forfeit it; --- thinking, on the other hand, even if it wished to renounce the
office of judge, could not do so except by itself judging; the final judgment
therefore belongs to thinking; on its truth everything rests. But the
investigation of the truth of our thinking necessarily calls forth the
question how thinking arises at all, and in answering this we then meet the
opposed theories of the sensualists and the idealists.

When the sensualists lay down the principle that everything is owed to
impressions from the existence outside us, that the human soul is a
``\textit{tabula rasa}'', and that the universal concepts therefore arise by an
abstraction, it follows from this that thought is only a shadow, while reality
itself must be placed in empirical objects existing outside us. With them,
then, the a posteriori is the main thing. --- But precisely in thus overlooking
the moment of subjectivity, they fall prey to subjectivity; they do indeed
assert that the human soul is a \textit{tabula rasa}, but since absolute
passivity cannot take place, there must necessarily be a subjective moment
present in all the subjective impressions. One must thus proceed from special
sensations in order to arrive at the universal concepts; but since this cannot
take place, because the phenomena of experience are particular, one is forced
to acquiesce in subjective assumptions, since one has denied the universal
concepts absolute validity. Herein, then, lies the contradiction, that
% --- p. 3 ---
\opage{3} subjectivity on the one hand counts for so little that on the other
hand it counts for everything.

If it is thus evident that one is caught up in subjective opinions whenever one
attempts to comprehend objective things, these opinions are nonetheless
universally valid from the standpoint of the subject. The universal concepts,
therefore, cannot derive from an external experience, but must lie in
subjectivity as such.

But when all phenomena are apprehended by means of subjectivity and led back to
subjectively necessary and universal forms of concept, the universal concepts
are, as a consequence, necessary to every experience, so that the opposition
between subject and object is not an opposition between thinking and
experience, which flow dialectically over into each other, but an opposition
between the phenomena and pure, absolute objectivity, the \emph{thing in
itself}.

If one now investigates what the thing in itself really is, it will turn out
that it has only a thought existence. For one cannot experience the thing in
itself, since every experience rests on subjective concepts; this alone, then,
does one know about the thing in itself, that such a thing is. I think not only
my thoughts, but I also think that something is outside my thoughts; one must
then seek this its being outside, not outside, but within the bounds of
subjective thinking.

It is by this path that the idealists have developed their theory, in that they
teach that the whole world is posited by the human I itself, although the I
was not conscious of this act; the a priori then takes first place, for the
hidden I itself is precisely this \textit{prius}.

The science of our time has, however, fully proved that idealism becomes
nihilism when it denies objective existence and seeks reality only in its own
self-activity.
% --- p. 4 ---
\opage{4} For if all objects are only representations and images formed by the
human soul, what then is this I itself? If it is regarded as a resting reality
and not as an empty image, must not the same be able to be said of the other
objects? And if, on the other hand, one refers the existence of the I too to
representation, one surely declares the I itself a representation formed by
the I. --- The result, then, is an infinite series of representations of
representations.

When, therefore, subjective idealism makes the nature of thinking the object of
a thorough consideration, it is forced to assume an objective existence and to
concede that activity and representation could not be, were there not
something that gave activity a substantial content.

The human I thus becomes, precisely at the moment when it rejoices in the
feeling of perfect independence, the prey of an alien power, and plunges into
the abyss of emptiness just when it arrogates to itself divine majesty.

But precisely this fall of subjectivity prepares salvation and makes a true
reconciliation possible between subject and object. For if sensualism proves
that there is not a single point in the life of the soul that cannot be led
back to an impression from the existence outside us, and if subjective
idealism on the other hand demonstrates the uninterrupted self-activity of the
human soul, then it follows in truth that thinking and objective existence
cannot belong to two opposite spheres; they must then be of the same nature,
and we cannot then, as the saying goes, \emph{know ourselves} (γνᾶθι σεαυτὸν),
except through the cognition of the objective phenomena.

Perhaps someone will object that we thus abolish the important difference
between individual subjectivity and universal objectivity; but this objection
proceeds only from an abstract and external view of the matter: for if we
% --- p. 5 ---
\opage{5} consider only the external world, it cannot occur to us to deny that
every human being has his subjectivity and everywhere finds himself surrounded
by objects that set themselves up as barriers against him; but the matter must
be seen from the inner side, and one will then observe that every subjectivity
owes its formation to objectivity, in that the subjective and objective moments
will be seen in their mutual, infinite transition into each other. --- It is
precisely the nature of objectivity everywhere to develop subjectivity, as it
is the aim of subjectivity to cultivate objectivity. Therefore everything must
be sought \textit{a posteriori} as well as \textit{a priori}, referred to
thinking as well as to experience. And from this standpoint, therefore, the
remarks so often made, that the cultivation of empirical knowledge would be the
restriction of philosophy, or that the flourishing of philosophy would entail
the neglect of empirical knowledge, do not hold either.

From this it follows, then, that the middle way in the realm of the sciences
deserves no particular recommendation, and that, much rather, every endeavor to
keep to the middle way is mediocrity.

\parag{2}{Subjective Faith and Speculative Logic.}

It is one thing to recognize the identity of phenomena and thoughts, of
empirical knowledge and philosophy, and another to subject the moments of this
identity to such a scientific analysis that the identity is seen in its
mediation, seen through and comprehended from all sides. What matters here,
then, is to find the right starting point.

In immediate practical life the relation of the objects met with everywhere to
individual subjectivity has the greatest
% --- p. 6 ---
\opage{6} influence; everything is apprehended by the inner and outer senses,
and one defends everything just as one has apprehended it. The one holds to the
truth of his statement because his immediate apprehension of the truth happens
just now to be so constituted; the other, for the very same reason, defends
the exact opposite. They contradict each other, and both lay their hands on
their hearts in confirmation of their opposed statements. --- What, now, is to
move us to respect their assurances? Everything depends on subjective mood. By
not granting such statements our credence, therefore, we offend not only
against truth in general, but also against the person of the human being who
proclaims the truth. Subjective faith constitutes the strength and, so to
speak, the marrow of every human being's life, and cannot be changed without a
complete change in the human spirit itself; if we therefore deprive anyone of
his subjective faith, we deprive him of life. --- It is therefore not to be
wondered at if human beings, in defense of their various faiths, not only
dispute and contradict one another, but also, as life shows by so many
examples, contend and fight with a vehemence as though the struggle were for
altar and hearth.

But as subjective faith kindles the flames of this war, so it also has the
means to quench them. And had it none such, no contradiction could have arisen
at all, for in order to disagree there must necessarily be something in which
one agrees. An absolute disagreement cannot take place, for if the one did not
understand the other, each of them would have been left to, and could only have
spoken with, himself; in order to be able to dispute, one must then necessarily
have a medium in which one can meet. And what is this medium other than
precisely thinking itself? All must think when they wish to explain their
faith; but what is
% --- p. 7 ---
\opage{7} thinking other than seeking out and uncovering the universal in the
particular, unity in multiplicity, identity in opposition?

Precisely for this reason one must not lose heart amid all these disputes;
there always remains in the disagreement a certain agreement. When the debate
has risen to the highest point of contradiction, when you hear the asserter's
``\emph{it is true!}'' over against the denier's ``\emph{it is false!}'', you
must nevertheless not, say the logicians, give up hope of being able to
reconcile the agitated minds. This can be done by an easy subtraction: just
take away the conflicting predicates ``\emph{true}'' and ``\emph{false;}'' there
then remains only that ``\emph{it is}'', and if you promise to draw from it
many, and indeed important, consequences for both parties, the combatants will
lay down their arms and willingly concede that the thing about which they are
contending, is.

When, then, this debate, this negation, this σκέψις has advanced to that
abstract concept, so trivialized in our time, \emph{being}, whose denial
contains its affirmation, it must necessarily cease. It is here, then, that the
absolute starting point of true knowledge must be sought. But it is quite
another question whether, by this beaten path, namely by pursuing the universal
concepts, we can succeed in fully uncovering the fullness of life from which we
set out.

When the dialectical stimulus lying in the concept \emph{being} had been
discovered by Hegel, so that logicians were now no longer compelled to remain
standing in the tautology of universality, like the Eleatics of old, who
continually repeated their τὸ ἓν τὸ πᾶν, many thought that the fullness of time
had come. They then liked to speak of the wondrous poverty and extraordinary
richness of this concept, since it showed itself at once as the simplest, in
that all preconceived opinions had been cast down
% --- p. 8 ---
\opage{8} into the abyss of abstraction, and yet at the same time enclosed all
the treasures of wisdom within itself; in its bosom, they thought, the present
and the future, God and the world, lay hidden; in a few words, they extolled the
concept \emph{being} as the egg that in truth enclosed the universe within
itself, and only needed to be brooded and hatched by a proper speculative
philosopher in order to let the true shape of things emerge from it. --- Under
this star was born logical mysticism, which, as befits a true mysticism, bade
farewell to life and its actuality.

But we seem to contradict ourselves when we assert that the absolute starting
point must be placed in the concept of being, and at the same time declare that
the analysis of the logical concepts cannot lead to the realities themselves;
for if we cannot reach the goal by this path, we must seek a new starting point
elsewhere; we cannot then rest content with the beginning once made, but must
begin afresh once more, or perhaps several times.

One ought therefore to observe in what relation the universal concepts stand to
the particular objects. The common opinion about this relation is that the
logical concepts are hollow forms into which the solid objects are poured; if
the forms are taken away, there remains only a jumbled mass of objects; if the
solid objects are taken away, the pure concepts will, in the world as in the
logical system, show themselves only as an artificial net. And this separation
between the universal concepts and the particular objects has a true as well as
a false side. It is true insofar as it is maintained that the form of the thing
is not the thing itself, and that the thing is not the form; for no one can deny
that the form is precisely the negation of the thing and thus its limitation;
and how, too, should the universal, the abstractly thought, be able to be
something particular? By the force of their own
% --- p. 9 ---
\opage{9} contradiction, therefore, these moments exclude each other. But the
separation of the moments of existence is fundamentally false if by it is meant
that the things and the forms could have subsisted separated from each other,
had not God united them by a special act of will. On closer consideration of the
matter this will easily become evident.

We are fully right to see in the concept \emph{being} the universal in its
highest abstraction, but logic teaches us at the same time that \emph{being}
cannot be thought, nor can it be, if there is not something whose predicate it
is; in order that \emph{being} itself can be, more than pure being must be
given. The pure universal thus splits into two moments, of which the one is pure
\emph{being} itself, the other this \emph{something} of which being is
predicated; for if these two moments are confounded with each other, being
becomes the same as being; but such a tautological being cannot be. But if they
are kept in their opposition to each other, and we call the one \textit{A}, the
other \textit{B}, it turns out that what is not in \textit{B} must be in
\textit{A}, or --- in order that the universal can be, the particular moment
must necessarily be present as well.

The universal, then, spreads out on all sides into particular objects. But
while it has now become clear that the universal is not present anywhere where
the particular is not also, nothing has yet been said about this particular
moment except that it must be sought everywhere; in other words --- the
particular moment is here apprehended not in a particular, but in a universal
way. What is lacking, then, is the truly particular, and we must then
investigate where this is to be found.

The truly particular is like a fixed and immovable point which, considered in
and for itself, cannot be referred to a universal concept, but must be capable
of being definitely pointed out, so that
% --- p. 10 ---
\opage{10} %
one can say of it: ``it is here! there! thus!'' The particular moment is a
positive and real moment, which stands firm by itself and must be sought in the
realm of empirical knowledge. For the ontological concepts to be grasped,
therefore, a certain empirical knowledge is required, since all such concepts
rest on a negation of empirical objects, just as the concept of an object
cannot be understood either, except with regard to the determinate empirical
object.

But the image of this or that object is presented to the soul only in order
that what is common to all, not what is peculiar and particular to the
individual objects, may be held fast; the special object is thus called forth
by means of representation in order that the particular, that is, the
\emph{empirical} moment may be negated. Since empirical knowledge is thus
appropriated only in a negative way, there can be no question here of knowledge
of the individual phenomena; but when the speculative cognition of the logical
concepts has been fully developed, the result is that the principle of all
existence is the absolute Idea, or, in other words, that the whole series of
real phenomena is rational, so that objective and concrete reason can no more
be grasped except through the cognition of the empirical objects than anything
can exist if it conflicts with the nature of absolute reason. --- But in this
way we then leave logic, whose concepts are both universal and necessary, and
turn to empirical knowledge, whose realm encompasses the particular and the
contingent.

% --- p. 11 ---
\opage{11} %
\parag{3}{Vulgar Empirical Knowledge.}

Vulgar empirical knowledge, in which everything is particular and contingent,
is rightly regarded as the negation of logic; for while in logic we proceeded
by denying the senses all authority, we are now to give up the investigation of
the universal concepts and appeal to the testimony of the senses. --- But here
we seem to come into contradiction with ourselves; for we have asserted that
there is no difference between the a priori and the a posteriori, and must now
nevertheless admit that a difference does obtain. For we know, even before we
have apprehended the individual objects through the senses, that there can be
nothing that conflicts with the ontological laws; everything that ontology
teaches can thus be predicted of the objects, and this is precisely the a
priori; but how things are constituted on their particular side, that we can
learn only by experiencing the things themselves, and this is the a posteriori.

But this contradiction finds its resolution in the fact that the difference
concerns only the phenomena as such, while the identity is to be sought in
their \emph{essential} existence. Now if these were absolutely opposed to each
other, the knot could hardly be untied; but everyone knows that thinking does
not let the phenomena, as they are immediately apprehended by means of the
senses, vanish into essence in order that they should remain sublated, but
precisely in order that they should be restored again out of essence, as out of
their true ground.

Speculative logic teaches that essence itself is revealed in the phenomena, so
that these are the phenomena of essential existence, but cannot be held fast in
and for themselves; and experience shows that even those who with the greatest
zeal devote
% --- p. 12 ---
\opage{12} \setcounter{footnote}{0}%
themselves to empirical knowledge come up, contrary to their own expectation,
against essence in grasping at the phenomenon. What joy, for instance, would he
have who delights in contemplating the rays of the sun, if these all at once,
as with a single stroke, met his eye, so that for sheer gazing he could think
nothing and feel nothing? What joy or fruit would he have who, like one bereft
of his senses, were held so fast in beholding the individual things that he
understood nothing although he heard with his ears, apprehended nothing although
he saw with his eyes? Therefore we also see everywhere in history, where a
certain immediate view predominates and where philosophy has not yet been
developed into an independent science, not the crude and barbaric empiricism
that extinguishes and supplants all thinking, but the freest interplay of the
moments of empirical knowledge and of philosophy\footnote{Cf.\ Herodotus and
others.}. Now, to be sure, there comes in here the important question whether
these or those occurrences really took place as they are narrated, and one then
seeks to get clear about how far the external phenomena have reality, as though
the historical inquirer had no other task than that of proving that they are
reliable and have corresponded to their description; but one must then take
careful note that the phenomena are not present when one thus seeks them most
eagerly, and that it is much rather their absence that makes us feel the want
of them, so that there can be no question here of a sensuous beholding of
present objects. Indeed, were the objects that the empiricists search for so
eagerly already uncovered, precisely delimited and brought forth into the
light, and had they thus already attained their wish, then either their
activity of understanding would have to cease, or they would have to strike out
on another path for their inquiry.
% --- p. 13 ---
\opage{13} %
We must then consider more closely the distance that lies between the external
phenomena and the consciousness of the spectator, in order to see how the
identity can be restored.

The observer, after all, cannot render the phenomenon that has presented itself
to his gaze in the same immediate way in which it presented itself to him; he
cannot describe it by pointing at it; if he therefore wishes to hold it fast,
he is compelled to pass a judgment on it. But what is passing a judgment other
than referring the phenomenon to the totality of the universal concepts and
comparing it with the totality of the experiences already made, in order
thereby to assign it its proper place?

But here again are two moments that must be carefully distinguished from each
other; for what place the phenomenon is to occupy depends on the one hand on its
own nature, on the other hand on the subjective peculiarity of the observer;
and since the phenomenon makes itself known to first consideration only in an
abstract way (§~\textit{2}), it will perhaps come to occupy far too slight a
place in his consciousness, because its true properties have not clearly come
forward to him. Therefore the way in which the phenomena are apprehended is
wavering and indeterminate, and one must use the greatest caution lest an error
creep in.

But this difficulty grows when one considers that it is only the fewest
phenomena that present themselves immediately to the observer, while in most
cases one must see with others' eyes, hear with others' ears, rely on others'
understanding and power of judgment. Since the one observer most often observes
the phenomenon only with the other's eye, one must therefore at the same time
take care to remove scrupulously whatever untruth may have mixed itself into
the apprehension of the thing from one subjectivity or another, and this is
exceedingly
% --- p. 14 ---
\opage{14} difficult. For it is not merely another's mistake that one has to
fear; but since the reporter may not only have apprehended the object
incorrectly, but may also --- by virtue of reflection's capacity deliberately
to present the same object in a different way --- have distorted his judgment
according to arbitrary considerations, one must also fear a deliberate
falsification. One must therefore investigate the author's credibility; and
since the author's credibility depends on the way in which he describes the
thing, while this conversely depends on the author's credibility, one is forced
to seek out another, more general and surer criterion of truth. We are thus
forced to take refuge in the ontological concepts, and in the conditions
without which the thing cannot take place, in order to remove absurdities and
separate history from fable. Thus we do not believe someone who would tell us
of round squares, and are just as unwilling to believe the Decretals of Isidore
when they report that Victor, who was bishop of Rome in the \textit{2}nd
century, wrote a letter to his dear friend Theophilus of Alexandria, although
the latter lived in the \textit{4}th century. --- Now if we do not wish to
remain standing at the trivial tautological rule that the impossible is
impossible, we must determine more precisely, with the help of experience, what
the concept of possibility means. We then compare what we know from an earlier
experience with what experience subsequently presents to us, so that the former
shall not be sublated by, or be found to conflict with, the latter.

Here, then, one reasons by means of analogy and induction, in that one accepts
as true the phenomena that usually take place, but conceives suspicion against
those that are unusual. On the one side, then, one considers the sum of the
phenomena, on the other the particular phenomenon, for the judgment of which
the former are to furnish the standard. But within the
% --- p. 15 ---
\opage{15} likeness there nevertheless always remains an unlikeness; otherwise
nothing new would have been added; it is therefore not the side of the
phenomenon from which it resembles the preceding ones that comes into
consideration, but it is what is unlike in it that makes a judgment necessary.
--- But the judgment grounded on a phenomenon's likeness or unlikeness to a
preceding one is again external and uncertain; for because the object in
question has often taken place before, it does not follow that the same or a
similar one took place at the time when it is narrated to have taken place; and
because a phenomenon has seldom or never taken place before, it likewise does
not follow that the narrative of its having taken place must necessarily be
false. --- One must then take into consideration the greater or lesser degree of
probability, and this consideration now carries the spectator over into the
relative territory of subjective conjectures
(\textit{calcul des probabilités}). For where there can be talk of probability,
there truth must both be present and not be present. It cannot be present; for
if one did not lack it, why would one speak of its external shape, and not far
rather of truth itself; and were it not present at all, how could it occur to
one to seek it, how could one speak of a probability? The matter, then, seems to
stand as follows. Since one recognizes that it is fruitless to compare the
individual phenomena with one another, one compares them with the reality that
stands above them all, i.e.\ with truth itself. Now in thus abstracting truth
from the phenomena, one sets truth over against the phenomena in an external
way, so that it itself emerges in the shape of a new, though more excellent,
phenomenon, ---
for of the immanent relation according to which truth permeates and develops
all the phenomena, so that it makes them forms for itself, there can
% --- p. 16 ---
\opage{16} be no talk from this standpoint. But if one now asks how, in the great
series of phenomena, one can grasp truth's own phenomenon, one will find that,
objectively considered, it lies hidden behind the other phenomena as behind a
veil, while, subjectively considered, it dawns upon the spectator's
consciousness through an anticipation and happy presentiment, as a standard by
which he judges the individual phenomena. But since each spectator now has his
own opinion about truth, these different opinions must also be compared with
one another; and if one cannot then have absolute unanimity, one must at least
have unanimity in the majority, if one is to arrive at any result; and so the
general and public opinion that holds sway in the representation of the age
that has given the spectator his consciousness will be the deciding factor in
the matter. --- And yet we have here still no firm standpoint for our
judgment; for everyone knows how different opinions have been, in the course of
time, about what is possible or impossible. Thus, to name one example, the
realm of possibility had a far greater extent in the Middle Ages than in the
time in which we now live. We see, then, that the further we pursue truth along
this path, the further it withdraws from us.

And lest anyone now object to us that we are merely seeking difficulties in a
sophistical manner where in actuality there are none, we must again and again
recall that we are far from failing to appreciate the sincere search for truth,
the sense for it and the knowledge of it, by which many of the warmest
adherents of this method distinguish themselves; --- for it is one thing to
judge the cognition of truth that one author or another possesses, another to
investigate the method by which he believes he has
% --- p. 17 ---
\opage{17}acquired it. We do not censure the comparison of the individual moments in and
for itself; for we know full well that truth does not come in so palpable a
manner (μετὰ παρατηρήσεως) that one could say: ``Lo, here or there it
is!'' and we do not think that the observation of the individual phenomena may
be neglected, for if the individual parts are taken away, the totality vanishes
as well. But unless we are much mistaken, the error of this method lies in
this, that although the individual phenomena are set up side by side in an
external manner, without permeating one another or being permeated by truth,
one nevertheless believes that by comparing all the individual phenomena one can
grasp truth, while it is yet as clear as day that by this comparison one can
arrive at no other result than this: to see everywhere, not truth itself, but
only its semblance, its alteration. --- It goes with him who seeks truth along
this path as it goes with one who in a foreign land seeks out a friend, enters
the inn, hears that he was there a short time ago but has just this moment gone
aboard, and now, when he has come down to the harbor, learns that he is already
out at sea. For if one seeks truth in subjectivity, criticism refers one to
objectivity; if one then wishes to see the objective, universal idea of truth,
it analyzes the individual phenomena, and if the talk now turns again to this
or that phenomenon, it appeals to the totality.

From this it follows, then, that the rules of criticism are either so general
that they are entirely empty of content, or that, when they determine something
about the more particular relations, they are of such a nature that one can say
of them what most often holds of German grammar in its rules on the gender of
nouns: ``\textit{quot regulae, tot exceptiones!}''
% --- p. 18 ---
\opage{18} \parag{4}{Transition to Speculative Empirical Knowledge.}

The external phenomena, then, must no longer be considered in and for
themselves; what remains of them is only the subjective opinions about them,
and in the sum of these truth must then be placed. --- Thus, even if I do not
know whether I have judged this or that external phenomenon correctly or not, I
know at least that I have judged in this or that way; and even if it is not
certain whether this or that event took place or not, this much is certain,
that it has been told and must therefore also have been thought by someone; and
though I may doubt that miracles took place in antiquity, I cannot doubt that
the ancients believed in miracles. It is thus not the historian's business to
investigate whether the individual things took place in the determinate
external manner in which they are reported; but, just as no deed can do itself,
but every deed presupposes someone who does it, so too no occurrence can be held
fast in and for itself, but must be able to serve to give us a true
representation of the standpoint of the person to whom it is attributed.
Hereby, then, the external phenomena are entirely negated. --- Spiritual life
appears in sensuous phenomena; the phenomena express the essential opinions of
the soul, but the essential opinions are thoughts; it should therefore be the
historian's task to demonstrate the shape of spirit at the different times, the
education of the peoples, reflected in the deeds of the human race as in a
bright mirror. --- If one considers the matter from this side, one arrives at
the standpoint of a, so to speak, historical idealism, which is at least right
in this, that thinking and reason seek themselves in history. It is, however, so
% --- p. 19 ---
\opage{19}\setcounter{footnote}{0}%
far from the case that this idealism denies vulgar empirical knowledge all
authority, that it rather confirms it, only with the difference that the two
moments of history, the real and the ideal, are not taken as though their union
were to take place in an immediate manner, the account of the historical data
being seasoned with accidental subjective remarks, and a spiritual stream
welling up from the dark spring of pious feelings to water the barren fields of
earthly existence. While idealism rejects the holding fast of the external
phenomena as such and reaches for the inner essence, it recognizes at the same
time that it is only through the phenomena that access stands open to spiritual
existence. It would indeed also be impossible to form a picture of the spirit of
individual human beings or of a whole people if one knew only the bare events
and did not also know the names of the persons acting, or if, to name an
example, Caesar's deeds were attributed to Pompey; impossible to know and
understand the activity of spirit if one did not have exact knowledge of what
external barriers it had removed, what raw masses of nature it had worked, what
objective spheres it had formed\footnote{\textit{Avec le monde a commencé une
guerre, qui doit finir avec le monde, et pas avant; celle de l'homme contre la
nature, de l'esprit contre la matière, de la liberté contre la fatalité.
L'histoire n'est pas autre chose, que le récit de cette indéterminable lutte.
Michelet: Introduction a l'histoire universelle.}}. It is, for example, of the
greatest importance for the right understanding of the spirit of the Jewish
people whether we assume that the many miracles attributed to Jehovah actually
took place in external nature, or whether we take them merely to be products of
the subjective representations of pious individuals.
% --- p. 20 ---
\opage{20} True history, then, does indeed require that one seek objective reason in the
external phenomena, but at the same time that the mass of phenomena be
completely consumed by the fire of thought, in order to rise like a phoenix from
its ashes. But the labor that this thorough negation and the ensuing
restoration of the phenomena entail seems to most historians difficult and
hazardous. They do not, indeed, excuse themselves from thinking; they are rather
quite willing to weave the weft of their reflections into the warp of the
reported facts; for this way of thinking seems to be the safest: should it
namely now happen that their subjective reflections were found uncertain and
inconsistent, they still always have the reported facts left; to these they can
then in any case hold and say: ``The fact itself is here, and it is at least
beyond doubt that it occurred at that time or in that place.'' If, on the other
hand, one wishes to leave the empirical phenomena and to turn one's attention
to the objective reason in them as such, they turn their backs on us for fear
of what they call philosophical preconceived opinions and arbitrary a priori
construction.

If it were now really the case that, after having negated all the phenomena, we
turned to this or that system in order to apply certain predetermined formulas
to them arbitrarily and in an external manner, it would in actuality be
fruitless for the philosopher to spend time and diligence on history. But this
is not how things stand with philosophical negation.
--- When in everyday speech one names the concept nothing or negation, one
usually thinks of a, so to speak, empty nullity; but it is easy to see that
there are as many negations as there are modifications of positive things,
since negation is indeed the thing's own limitation or determination. If, for
example, someone were merely to say:
% --- p. 21 ---
\opage{21} \emph{``I have nothing,''} one would not properly understand him; if, on the
other hand, he said: \emph{``before I had wealth, now I have nothing!''}, one
would easily grasp both what this nothing stood in opposition to and of what
nature it itself was, and understand that the man had become poor. --- Daily
life itself also recommends this way of inferring. Thus when we judge the
actions of our intimate friends, we do not stop at their actions as mere
external phenomena, but always look at the same time to the frame of mind, the
causes and grounds, from which they proceeded. Every action has its ground, but
the ground is the negation of the phenomenon; the phenomenon thus indicates the
negative of itself. --- That is also precisely why all the proofs of God's
existence stand and fall with the truth of speculative negation; for each of
them aims precisely at negating a general phenomenon, as the world immediately
presents itself to us, and leads us through its negation to God's essence, from
which the negated phenomenon is then restored anew.

Thus if we first consider the cosmological proof, we meet in it first the
observation of \emph{an empirical phenomenon:} \emph{``Lo, the world is!''}; but
\emph{negation} steps in at once, as one infers from the contingency of the
world: \emph{``If the world is considered in and for itself, it is nothing, has
no substance in itself''}, and only now follows \emph{the conclusion:}
\emph{``The absolute substance is God,''} whereupon \emph{the phenomenon is
restored again:} \emph{``From God's substance the world receives its
substance.''}

From this it is also seen that the essence itself is determined according to the
phenomenon that is negated; for while in the cosmological proof we negated the
world's substance and through this negation came to God's substance, we would,
had we negated another phenomenon, e.g.\ human thought, which
% --- p. 22 ---
\opage{22} is to take place in the ontological proof, have obtained another form of God's
existence. Here too one can clearly see how dangerous it is to stop on the
middle way, something we have already once sought to prove.

\emph{The point} of the ontological proof can be stated thus: \emph{``I think,
therefore God exists.''} The phenomenon immediately present is thus human
thought and self-consciousness. \emph{The absolute negation reads}, then, as
follows: \emph{``Human thought is nothing of itself''}; if this negation is
carried through, the following conclusion emerges: \emph{``Absolute thought,
i.e.\ God's self-consciousness, is''}, and now the phenomenon is restored:
\emph{``If God's thought and the absolute divine self-consciousness is, then the
existence of human thought can also be explained from it.''}

But when Kant now undertook to reject the ontological proof, his separation
between a real and an ideal existence, on which his whole argument rested, led
him precisely onto the bad standpoint of the middle way. ``That man is and can
think,'' he said, ``is a fact; if this is held fast, it is perfectly certain
that man is able, as often as he wishes, to refrain from thinking God just as
well as to think him.'' He thus holds fast to this empirical phenomenon, lets no
true negation come about, and so easily makes the proof waver, since the divine
thought cannot be placed in an external manner beside human thought. Had he, on
the other hand, raised the question whence this human capacity to think had
also arisen, his idealism would have separated out the raw, undigested
phenomena from human
% --- p. 23 ---
\opage{23} consciousness, and he might then perhaps have come to acknowledge and develop
the ontological proof properly.

Just as the bad middle way has had the most harmful influence on theology, so it
has also had it, and will continue to have it, on history, until a true negation
of the phenomena is carried through. For it is with historical ideas quite as
with the theological ideas we have just discussed. One should first of all
strive, through the negation of the reported facts, to penetrate to the
principle itself. If, e.g., someone wishes to comprehend the spirit of
Catholicism, he must first, as far as possible, consider the sum of the
phenomena belonging to it, in order to arrive through their negation at the
Catholic principle; for if he does not do this, he must involuntarily snatch a
principle out of the air, and then it will constantly collide with the
historical phenomena. Just as, then, in the cosmological proof the phenomenon of
the world was first given in its immediacy and contingency, so too the general
phenomenon of Catholicism must first be observed in its contingency before one
can penetrate to its ground. And once one has thus appropriated the principle,
one must again remember well that it must not be held fast in its abstraction,
but that also in its particular development it can be properly appropriated
only through the negation of the particular phenomena, so that thought never
abandons the historical object. --- There are, then, two errors against which
one must especially guard: on the one hand, that one not develop purely abstract
concepts from such a general principle while neglecting the particular
phenomena; on the other, that one not transfer the general principle in an
external manner onto all the details. For there is nothing in the study of
history that more betrays the bad middle way than applying the general
% --- p. 24 ---
\opage{24} categories: freedom, despotism, hierarchy, or others like them, as clothing for
the individual events, so that one does not develop these concepts themselves
but only fills them in with external details. If, e.g., the concept of the
hierarchy is to be developed historically, one cannot know from an abstract
definition that it changes in the course of the centuries, and if one asks
about its individual changes, then it is not enough to enumerate the individual
actions or utterances of Irenaeus, Cyprian, and Gregory. --- There are, then,
three moments to be distinguished in such a development: the first is the
general concept of the hierarchy; the second, the individual actions and events
that are expressions of the hierarchical idea; the third, the identity of both
moments. And what is meant by this identity is not hard to understand; for if
one asks why an application of the universal to the individual should take
place, then, unless we are mistaken, no other answer can be given than that the
general concept, the idea, develops in a peculiar manner in the individual
phenomena, so that the concept of the hierarchy appears in a different manner
in Gregory the Seventh than in Irenaeus. But from this it evidently follows that
Gregory's undertakings can be considered from a twofold side: considered from
the one side, they are peculiar, so that we cannot find them again anywhere else
in history, and, as such, they constitute the fixed and unshakable moment we
spoke of before; but on the other side they reveal in a particular manner the
idea of the hierarchy and are thus everywhere permeated by the general moment.
The individual moments have thus received the stamp of the universal, while
conversely the universal idea has obtained its peculiar modifications. --- But
from this there emerges the third moment, which can neither be determined as a
single object of experience by a ``thus or so!'' nor
% --- p. 25 ---
\opage{25}%
be said to be an empty general concept, applicable in any way whatever, but
must be called the historical limitation of the concept. And such empirical
knowledge, which fully corresponds to and agrees with the underlying reason, is
rightly called speculative.

\parag{5}{The Relation between Vulgar and Speculative Empirical Knowledge.}

The difference between vulgar and speculative empirical knowledge is striking:
the former observes the phenomena from without, the latter from within, and,
while the former heaps together a multitude of contingent phenomena, the latter,
on the other hand, depicts the organic connection of the moments in such a way
that no phenomenon comes into consideration except insofar as it is recognized
to have proceeded from the idea itself. Hence the range of the individual
phenomena can be far greater in vulgar empirical knowledge than in speculative;
for in vulgar empirical knowledge we ask only whether the events took place,
not \emph{why} they took place, and therefore have only the task of observing
and determining exactly as many individual phenomena as possible. The works of
the critics offer many examples of the rich yield of such careful and untiring
investigation. The philosophers, on the other hand, constantly pursue the
constitutive principle of reason and therefore overlook, as contingent and
insignificant, many phenomena in history that are treated by the critics with
the greatest exactness. Yet one must not believe that this indifference is a
necessary consequence of the nature of the matter, as though this must entail a
certain speculative arrogance; for, objectively considered, everything can be
appropriated and comprehended, --- it is rather a consequence of the
imperfection
% --- p. 26 ---
\opage{26} of human thinking, in that this can indeed distinguish the general lineaments
of objective reason from one another, but not the countless nuances in which it
appears, so that philosophy precisely calls for modesty; --- and so far is
speculation from weakening the interest in the empirical phenomena that, on the
contrary, it increases it. A single action can express only one moment of the
universal idea; if, therefore, the other moments are to come in, the other
actions must also be taken into consideration. To find out the objective and
historical concept, one must therefore have an apparatus of many phenomena; the
more phenomena we therefore take up from vulgar empirical knowledge, the better
we will succeed in separating out and developing the essential phenomena.

If, therefore, one wishes to describe Gregory the Seventh's hierarchy in the
vulgar manner, namely by citing his individual utterances and undertakings, one
need only weave in one's definitions here and there and otherwise pick out and
present in a convenient order such passages as are suited to illustrate it;
(and when a goodly number of such passages, heaped together from everywhere,
then make a show in the row of definitions, they will at the same time be a
brilliant proof of the great learning the author possesses). If, on the other
hand, one wishes to comprehend this pope's hierarchy in its whole truth, one
cannot rest content with having a sum of selected passages; one will then find
that one can catch the finer threads of the concept in no other way than by
carefully elucidating all the grounds that have given it its peculiar
character. The concept of truth is thus entirely different in these two kinds of
empirical knowledge. In vulgar empirical knowledge that is true which has
actually taken place, without regard to whether it is something significant or
something insignificant; in speculative empirical knowledge, on the other hand,
truth depends
% --- p. 27 ---
\opage{27} on the relation in which the determinate object stands to the constitutive
principle, and much may have taken place in which speculative empirical
knowledge finds but little truth. --- Thus, while Isidore's Decretals are
absolutely rejected by the critic on account of their untruth, they nevertheless
have a deeper truth if the phenomenon is considered from its other side; and we
have already remarked that these two modes of consideration must be weighed
against each other. --- A similar difference obtains where the talk is of
historical proofs; for though we can often document the matter in question by
definite statements or events, we also often maintain, even when we have no
such things to hold to, that the truth of the matter follows from the general
relation of all the individual objects. It is precisely here that the
difference arises between the so-called juridical proofs and those that rest
only on a moral conviction; and as we assign the former to vulgar empirical
knowledge, so, if we do not wish to remain standing on the shaky ground of
probability (\textit{cfr.~§.~3}), we must assign the latter to speculative.

If vulgar empirical knowledge could itself solve the task it has before it, the
transition to speculative empirical knowledge would be easy and safe; but it is
easy to see that the truth of vulgar empirical knowledge cannot be determined by
the standard of speculative empirical knowledge alone; if one is to be able to
comprehend fully how these two can be neither immediately in agreement nor
absolutely in disagreement, it must be shown that the opposition between them
has a deeper opposition as its ground, namely an opposition between vulgar and
speculative logic.

The external mode of apprehending the objects of existence differs from the one
that penetrates to the interior of the objects in this, that the former
juxtaposes the objects in
% --- p. 28 ---
\opage{28} such a way that they do touch one another at certain points but remain separate
within the boundaries that enclose them, whereas the latter everywhere lets them
coincide and permeate one another. From this it follows, then, that vulgar logic
proceeds from the mutual diversity of the phenomena and ascribes only a formal
significance to their unity, whereas speculative logic leads us from the unity
out into the diversities. In vulgar logic, where the abstract concepts are held
fast by themselves and the real objects by themselves, there is nothing one
must guard against more than contradiction; for when things touch one another
mutually and yet it is the nature of the thing to keep everything foreign
outside itself, how then can the thing \textit{A} be at once the same as, and
yet not the same as, the thing \textit{B}? When, on the other hand, the
consideration is conducted from within, things are not seen in abstract
separation from one another; the moments pass infinitely over into, correspond
infinitely to, one another, and thus the absolute contradiction sublates
itself.

It now becomes easy to account for the many conflicts between everyday
understanding and speculative reason. For everyday understanding is outraged to
see rules such as that \emph{nothing can come to be from nothing}, or that
\emph{three cannot be one and one not three}, rules it has itself recognized as
certain and unshakable, trampled underfoot, if one may say so, by speculative
logic. It goes, however, with these two kinds of logic as with the two kinds of
empirical knowledge: they do not simply cancel or overthrow one another; all
these oppositions arise from one's having a different point of departure for
consideration, so that a remarkable \textit{concordia discors} obtains between
empirical knowledge and philosophy. For when the philosophers speak of the
absolute contradiction and teach the identity lying in it, they do not do so in
the name of everyday understanding; they readily grant that
% --- p. 29 ---
\opage{29} all the rules of everyday understanding must be accepted by one who stands at
the standpoint of everyday understanding; they recognize full well that the
middle way deserves recommendation if one wishes to keep to the common opinions
about philosophy and history; but they want the standpoint itself changed; they
want to change, not the laws of vulgar logic, but vulgar logic itself. The
empiricists, on the other hand, do not want to change their standpoint, because
they cannot be persuaded that their method is really skewed and one-sided. ---
It is, then, self-evident that vulgar and speculative empirical knowledge
cannot be placed side by side in an external manner in order, as by a
subjective trick, to be reconciled with each other, since the wavering between
the two that would thereby arise would call forth a skepticism undermining both,
but that the transition from the one to the other happens, and must happen, in
countless ways by the thing's own objective power. For the outer must always
have the inner in it, --- and it does not rest on a subjective whim whether
speculative logic is to have precedence over vulgar logic or not; it lies
rather in the very nature of vulgar logic that, if only it is carried through
and does not stop halfway, it must lead to speculative logic. It is precisely
this relation that makes speculative empirical knowledge recognize an essential
identity between itself and vulgar empirical knowledge. --- From this it is
also evident with what freedom the speculative concept can develop. Many, to be
sure, are frightened off when one recommends a concept or a system to them,
because they reject every edifice of concepts strictly measured out with inner
necessity as an intolerable and constricting iron harness in which their
thoughts are tormented; but this fear is not grounded in the true nature of the
matter, but is a figment of their subjective imagination.
% --- p. 30 ---
\opage{30} For the system does indeed require that the necessary concepts be held fast,
but the reason for this is none other than that these make their power felt
everywhere in life itself. It requires just as fully that the concept be
elucidated and demonstrated through all its peculiar modifications, and thus
preserves thought in its freedom in the fullest sense. Both the system and the
concept are in truth in possession of the freest mobility; for what else does
the system want than to develop all the moments of the concept insofar as they
mutually take one another up into themselves; and what else does the concept
aim at than precisely the development of the fullness of the pregnant idea?
From the inner fullness of the idea proceed the outer modifications of the
object, and, if these are considered from without, they are indeed countless
and contingent, but if they are considered from within, they are seen to have
proceeded from the idea's own self-limited power and are thus limited in the
midst of their unlimitedness. The necessity of the concept thus stands in no
opposition to the freedom of life.

Nor is the subjective moment of freedom overlooked because we will not leave
the true development of life's moments to subjective arbitrariness. For since
speculative experience at once corresponds infinitely to, and at the same time
deviates infinitely from, everyday experience, what is required, precisely, in
one who is to be able rightly to untie the knot of the intricate moments is a
gifted eye, a capacity to seek the idea immediately.

\parag{6}{The Concept of Science.}

What we have sought to show so far is, in a few words, this: that there are
three stages which our cognition has to traverse. The first bids us leave
everyday experience with its immediacy and contingency and pass over to
% --- p. 31 ---
\opage{31} speculative logic; the second negates logic and leads us out onto the wide
field that the careful testing of everyday experience opens to us; the third
presents objective reason made intuitable in the mirror of speculative
experience.

If one now wished to place these three mutually negating stages side by side in
an external manner, one would need only to seek negation and restoration in
three places and three times. But what is thus in general assigned to three
stages repeats itself again infinitely within each of them; for what holds of
speculative logic, that it looks away from experience altogether, we find again
within the boundaries of this logic, since it develops precisely through an
uninterrupted negation of empirical moments; and what has been said of everyday
experience, that it stands in opposition to logical reason, holds also within
the boundaries of this experience: its moments never pass over into one another,
and in all of them we see the alteration of reason, not reason itself. The
general result at which our cognition arrives through all these stages is, then,
this: to see the oppositions in their infinite passing over into one another.

But from this the concept of science emerges, for it is precisely the task of
science to lead to an absolute identity of the moments absolutely opposed to one
another. If, therefore, one wants a general and formal type of absolute science,
we can find one in our own self-consciousness; for since the human I preserves
its identity with itself in distinguishing itself from itself, we can rightly
say that it knows itself. Now since no self-consciousness can exist without being
mediated through real existence, i.e.\ through the negative of consciousness,
this limitation, if we
% --- p. 32 ---
\opage{32} look to the particular objects, is indeed a sheer negation of consciousness
and thus an annulment of all knowledge; but if, on the other hand, we look to
real negation in general, we can truly say that our knowledge is expanded the
more consciousness is limited.

For if self-consciousness cannot be except through an other that is opposed to
it, it must also at every point be permeated by this foreign moment and, as often
as it grasps itself, grasp this as well; there is therefore in self-consciousness
at once a separation between what is its own, the formal moment, and the
foreign, the real moment, and at the same time an identity of both; it therefore
knows its other at the same time as it knows itself. And from this it is also
easily understood in general that one can at once know God and nature without
any confusion whatsoever of the creature and the Creator. Those, therefore, who
distort the well-known proposition of the philosophers that only like can grasp
and cognize like (\textit{simile modo simile capere et cognoscere}), forget,
alas, that the absolute identity of self-consciousness with the object to be
cognized results solely and exclusively from the absolute opposition between
them.

It therefore follows from the very nature of science that every kind of
cognition that places thinking in an external relation to the objects to be
cognized must be banished from its sanctuary. Whatever in any way belongs to
everyday experience and everyday logic is too immediate and undigested to be
taken up into the realm of science.

But since everyday experience, as we have already remarked before, cannot be
dispensed with and yet stands outside the boundaries of science, one must
carefully distinguish between the elements that prepare science and science
itself; and important are the consequences that flow from this distinction.
% --- p. 33 ---
\opage{33} In the first place, then, we no longer stray onto the wrong path of
peddling, under the name of science, rudiments mixed with trivial concepts, in
the belief that we had appropriated science---a wrong path that would lead us
further and further away from all true science. To expound the science of
history is thus not to tell stories; only what we have conceived and seen
through in history should we reproduce; everything else we must set aside as
apparatus, which we shall have to examine anew in order to become clear with
ourselves whether or not it can occupy a place in science. But the bad habit of
putting forward critical elements under the name of positive science, and of
answering the one who seeks a concept that penetrates all particulars: ``that
is something historical; here there is to be no philosophizing!'', has spread
so far that dogmaticians who teach us virtually nothing else about God than
that there exists a highest being, who in Hebrew is called \emph{Elohim}, in
Greek \textit{Θεος}, in Latin \textit{Deus}, and in German Gott, and so forth,
and that these names have been explained in various ways, win a great name as
cultivators of ``positive science.'' Indeed, when speculative dogmatics
excludes as something foreign everything that the religious self-consciousness
cannot conceive, they even raise vehement complaints against it in the name of
Holy Scripture and the Church, as if speculative concepts did not require
persistent study of the teaching of Holy Scripture and the Church. And when
they see the dogmas proceed from one another in an unbroken sequence of
thought, they accuse the dogmatics of being derived from some philosophical
system or other, while they themselves think they best secure the transition
from one part of dogmatics to the other by citing parallel passages from
Scripture, and thereby betray that they are not free of a certain superstition,
since one no more keeps dogmatic error away from oneself by placing
% --- p. 34 ---
\opage{34} passages of Scripture at the head than one drives away the Devil by
making the sign of the cross.

Secondly, this separation makes evident how much exertion and diligence is
required in order to arrive at true science. When the empiricists and the
critics congratulate one another as the only ones who dare to take on the great
masses of the sciences, while the philosophers occupy themselves solely and
exclusively with their own empty speculations, they truly show that they think
far too highly of themselves and have far too slight a notion of the activity
of thinking. But the same people who pass so harsh a judgment on speculative
thinking set science itself now far too high, now far too low. They have far
too exalted a notion of it when they assert that human beings cannot, by their
thinking, conceive anything fully and with absolute truth, and therefore place
the entire yield of their study in the increase and arrangement of the material
of everyday experience; far too slight a notion of it when they think that,
with little trouble and on their own, they can themselves bring along the
requisite thinking, provided only the positive objects are at hand. Had they
properly considered how much diligence it takes to pursue an absolute concept
in the right way through all its modifications, they would surely admit that
Jehovah's word to Adam: ``In the sweat of thy face shalt thou eat bread!'' can
in this respect too be applied to the philosopher.

Although it is thus not to the critics and the empiricists alone that the
honor of exertion and diligence belongs, in another respect the burdensome part
of the work is nonetheless shifted onto them; for one must distinguish well
between the thought that is thought for its own sake and the preparatory
investigation that has its goal outside itself. The former belongs to
speculation,
% --- p. 35 ---
\opage{35} the latter to the everyday understanding and to criticism. For it
lies in the nature of critical labors never to be satisfying in and for
themselves; and this holds not only of the investigations that aim at inquiring
into the authenticity of a writing, without regard to whether it is a sacred
writing or not, but it holds also of the manner in which they treat the ideas
themselves. When criticism develops dogmatic concepts, e.g., creation,
inspiration \textit{or the like}, it does sometimes get so far as to have the
problem correctly posed, but the problem itself it never solves.

The creation of the world, it is said, must be thought neither in such a way
that we confuse the world with God, as the pantheists do, nor in such a way
that the world is assumed to have posited itself, as the atheists teach; for
such thoughts are impious and pernicious in their consequences; it must be
thought as a free act of God. But this freedom they do not themselves dare to
think; they thus demand something that they do not themselves do
(\textit{λέγουσι γὰρ καὶ οὐ ποιοῦσι}).

Inspiration is treated in the same way: it must be thought neither in such a
way that God has transgressed the laws of the nature of the human spirit, nor
in such a way that the spirit of man is assumed not to have been touched at all
by the Spirit of God, but in such a way that the Spirit of God has enlightened
the spirit of man. They thus demand in plain words that inspiration be thought
in this way, and yet we call on them in vain to think it. We are not, however,
so much surprised that they will not answer the question themselves, --- they
could indeed reply that they have not taken it as their task to solve the
problem, but only to pose it; but since the critical investigations have taken
the place of true science, their supreme contradiction shows itself in this:
that although they call on us to think, they nevertheless themselves declare
that these questions about how
% --- p. 36 ---
\opage{36} God created the world or how he inspired the prophets and apostles
are useless and impossible to answer.

What we have here sought to develop must not be presented as though we had
meant that the investigations of the empiricists and critics were to be
despised or at any rate regarded with indifference; for from the relation that,
according to our preceding development, obtains between vulgar and speculative
experience, it follows precisely that they are altogether necessary; only they
must not be confused with science proper. Should philosophy therefore reproach
criticism with the insignificance of its labors and boast of itself being in
full possession of the truth, criticism could with the greatest right answer it
in the same words with which Themistocles, in Plutarch, has the feast day answer
the day that follows it:
\textit{Αληθῆ λέγεις. ἀλλ' ἐμοῦ μὴ γενομένης σύ οὐκ ἀν ἦσθα.}

\parag{7}{The Scientific Method.}

In its striving to lead to an absolute identity of thought and the object
thought, science always has in view the goal of developing this identity
through all the particular modifications of the object. If, therefore,
self-consciousness cannot be except through an other, and indeed a real other,
the task becomes to bring about a complete reconciliation of this positive,
which comes forth in the whole multitude of positive elements, with
self-consciousness, so that the former can at every point correspond to the
latter. Every science depends on this reconciliation, but the reconciliation
depends in turn on the \emph{scientific method}.

% --- p. 37 ---
\opage{37} The contradiction from which the true method must necessarily
proceed shows itself in this: that while formal knowledge and the real fullness
of existence seem mutually to exclude each other, they nevertheless, as we have
previously shown, agree in this, that they both \emph{are}. And it is in this
point of identity that the whole opposition and every disagreement between them
is to find its sublation. We know, then, that it is objective being that
constitutes subjective being, that formal being proceeds from real being, the
being of our thinking from the being of substance.

But once we have sublated this opposition between the formal and the real
moment, a new contradiction shows itself in turn. For self-consciousness
attains complete agreement with objective reality only by an accommodation, in
that it renounces every other predicate than the simple one, that it \emph{is},
although it nonetheless has many more moments in it than this; and if we look
to the objective, we find that real existence too performs the same
\textit{συγχατάβασις}: although it overflows in infinite fullness, it too lays
aside all its predicates except the one, that it \emph{is}. Hence, as often as
we wish to acquiesce in the concept of being, the contradiction arises that of
an object which has a far richer content than mere being we assert nothing but
this one thing, that it \emph{is}. And we cannot avoid this contradiction,
because self-consciousness, which is the medium through which we arrive at this
assertion, will always assert its right. Now what thus holds of the beginning
holds likewise of every intermediate stage and of the conclusion itself. Every
concrete determination or designation that one is compelled to ascribe to
self-consciousness must also necessarily be asserted of real existence. If,
therefore, it is not sufficient to designate the relation
% --- p. 38 ---
\opage{38} between I and Me, so as thus to name the subject and object of
formal knowledge, by the relation between essence and phenomenon, cause and
effect, we can also be certain that this determination does not even fully
express the nature of the real object. Herein lies precisely the dialectical
stimulus that does not let the thinker find rest until he has arrived at the
developed concept of self-consciousness. But when one has at last reached this
highest point of the development, it will also turn out that not merely a
subjective advance has been made, but that the very same movement has also
taken place in real being. As, therefore, at the beginning of the development
we say that objective and real being constitutes subjective and formal being,
so we now understand and conceive, when we come to the conclusion of the
system, that it is an objective and real self-consciousness that constitutes
the subjective and formal self-consciousness.

And yet the contradiction does not cease even here; for the identity between
formal and real self-consciousness that we have thus established contains at
the same time an essential difference. --- While the absolutely real
self-consciousness has the whole fullness of reality in itself, it turns out,
after all, that the identity of formal self-consciousness with it was found
only in the purely formal. The reality that we previously gave up for the sake
of identity we must now restore anew. But in this way the dialectical stimulus
drives us through the whole of speculative theology.

That God is thus thought as absolute freedom by no means entails that one must
give up dialectical necessity; for where there is freedom, there \emph{must}
also be manifestations of freedom. But on the other hand, God's absolute
freedom is no more sublated by this necessity of the concept; for there is the
following circle in the proof: through the negation of the world we came to the
freedom that must
% --- p. 39 ---
\opage{39} have given the world its origin; as a consequence of this we must,
if we wish to exhaust the concept of this freedom, return once more from
freedom to the world it has created; as long, therefore, as we occupy ourselves
with the development of absolute freedom considered in and for itself, we shall
constantly feel this dialectical necessity of advancing further. The same
dialectical necessity that makes the \emph{ground}, in order to be a ground,
have to have a \emph{consequence}, the same thus also makes freedom, in order
to be absolute, have to have a world. --- And if we look to the special moments
in history, we also find in them the very same relation again. We cannot, e.g.,
conceive all the individual data that depict the relation between Germany and
Italy, between the political and the ecclesiastical power at the time of the
pornocracy, except by penetrating, through the totality (§~\textit{4}) of the
reported facts, to the principle in its simplest and most general form in such
a way that the very same contradiction we spoke of above will return here too,
if we now wish to remain standing at the first, abstract point of departure
(cf.\ §.~\textit{10}). Thus even human arbitrariness, which is so often assumed
to lack all rational necessity and to be entirely accidental, is bound to
objective laws, without thereby ceasing to be arbitrariness; and were it not
so, we could neither accept as true the objections that are made against the
indeterminists, nor rightly understand the nature and essence of obduracy. And
everyone who has noticed that man's resistance to the object of truth is itself
also determined and modified by the nature of this object will grant that the
cruelty of Messalina and Agrippina was of a different nature from that of
Fredegund and Brunhild, and that the cruelty of Theodora and Mariuccia in turn
had a peculiar character of its own. Thus, too, the obduracy that the Catholics
% --- p. 40 ---
\opage{40} showed when they denounced Luther as a heretic was of a different
nature from that which the Jews showed when they persecuted Christ, because the
truth here appears in altered form.

As in the sensuous world we look at the same point with both eyes, and both
therefore constitute, so to speak, only one eye, so on the other hand the soul
is indeed ascribed only one eye, but this in actuality looks at two points at
once and can therefore be taken as two eyes. For as soon as we begin to think,
we always come up against the reality that immediately confronts us, whose
negation thinking seeks to find; therefore the soul of the thinker must always
hold fast the positive fullness with the one eye, while with the other it seeks
to find the simplest and purest moment in the existence of that fullness. The
distance between the point of departure and the fullness that has resulted
from it thus constitutes the fundamental contradiction from which all the
contradictions lying in between the two proceed, so that every science can in
this sense be called hypothetical.

If this is so, one can also rightly assert that the chief task of method is
this: fully to uncover the fundamental contradiction in all the individual
spheres of existence. For the difference between these individual
contradictions rests on the difference between the positive realities from
whose negation one proceeds.

The first positive from which we proceed is, of course, our own
self-consciousness, which rests on objective reality, and since this is the
most universal foundation of all, it follows that its fundamental contradiction
must extend everywhere, and that we cannot take it to be tied to any
hypothesis, since hypothesis is precisely thinking's own necessary condition.
This contradiction we may then call the \emph{ontological}.

% --- p. 41 ---
\opage{41} The next positive foundation emerges when one proceeds by negating
reality itself, whose author is the absolute self-consciousness, and in
particular the reality that we have most immediately before our eyes, namely
our world; from this too there arises a fundamental contradiction, but in such
a way that the hypothetical here comes forward far more strongly; for it is
indeed certain that the absolute self-consciousness can acquiesce in itself,
but what reality is to proceed from it seems to be hypothetical and to
presuppose freedom as its necessary condition, since our world is presumably
not the only one in which the divine life can reveal itself. If, therefore, in
dogmatics we consider theology (the Kingdom of the Father), Christology (the
Kingdom of the Son) and pneumatology (the Kingdom of the Spirit), we may observe
that what we are to negate in the first kingdom is reality in its pure
universality, whereas in the second and third kingdoms we fix our gaze more
specifically on the reality of our world, so that in these the contradiction
and the necessity are set within certain definite limits; for it is a definite
world, namely our world, in which Adam sinned and Christ became incarnate, that
is here negated. We may therefore call this fundamental contradiction the
\emph{dogmatic}.

But as universal reality passes over into being the reality of this world, so
too the reality of this world issues forth in more special realities; hereby we
come to the third form of the fundamental contradiction, in which hypothesis
indeed asserts its power far more still, but which nonetheless, within its
limits, also brings dialectical necessity with it; and this is the
\emph{historical}.

We believe, then, that we have thus proved that every method leads through
contradiction, that every science requires that its moments mutually sublate
one another, and that all dialectic is negative.
% --- p. 42 ---
\opage{42} %
\parag{8}{Conceiving and Intuiting.}

The objective fullness that we seek, with the help of science, to take
thoroughly into ourselves must be considered from the following two sides. On
the one side, one must undertake an analysis of the individual moments in order
to get to see them in their proper mutual light; but on the other side the
totality of all the moments also presents itself to our eye and thus demands at
the same time an immediate intuiting. Every concept is thus determined in a
twofold way; for on the one hand it must, within its own limits, have a
peculiar distinctiveness by which it distinguishes itself from the other
concepts and keeps them outside itself, --- and this constitutes the fixed,
immovable moment of individuality that we have spoken of earlier
(§~\textit{2}), --- but on the other hand, in order to be able to be
established as a concept, it must necessarily also be capable of being
dissolved into other concepts. To grasp the relation between the individual
concepts by comprehending several moments at once may be called
\emph{conceiving}, whereas the peculiarity that belongs to each of them must be
referred to \emph{intuition}; science therefore takes both moments into itself,
both conceiving and intuiting. If one wished to determine a concept solely by
stating its relation to other, kindred concepts, then \textit{eo ipso} all
determination would thereby cease; for if one can determine what \textit{A} is
only by referring to \textit{B}, and again what \textit{B} is only by referring
to \textit{C}, determination vanishes into the infinite, and as shadow follows
upon shadow, every light of cognition is extinguished. And if, on the other
hand, the object is presented as a resting, dead object without inner mobility,
one does indeed get a petrified positivity, but here too no true cognition. The
identity of these two opposite
% --- p. 43 ---
\opage{43} sides, too, is present formaliter in formal knowledge, but is present
realiter and absolutely in God.

As the thinker is, so too are his thoughts. From this proposition one must
proceed if one wishes to determine the nature of divine thinking. The
absolutely real self-consciousness knows itself at the same time as the
absolute positivity. It concentrates itself, so to speak, in its will, in that
it affirms itself and thereby has all the elements of its life in complete
possession. In making itself an object for itself, it beholds not an empty
shadow but the fullest reality, and therefore all the thoughts that mediate
this its self-contemplation also partake of the same reality; for if God did
not himself utter his absolute ``Here am I!'' (\emph{Hinneni}), neither would
the things he posits attain to any true individuality. But in thus positing
himself and positing all, he at the same time intuits all. He intuits in
positing, and thus reveals himself as the absolute substance; he posits in
intuiting, because he is the absolute freedom: ``\emph{Truth is because God
wills it.}'' --- Now if this relation were confined to the life of the Godhead
alone, the whole of existence would split apart into monads which, as Leibniz
says, had no windows through which they could see into one another. But from
God's own individual existence it follows that his thoughts too pass infinitely
over into and interpenetrate one another. Since it is the same God, the same
self-consciousness, that thinks all, his different thoughts must also stand in
mutual connection with one another. For God does not merely intuit the
individual things piecemeal, since his knowledge is simultaneous, but he also
derives the one from the other and considers the mutual relation among them
all; and were it not so, he could neither clearly distinguish nor
% --- p. 44 ---
\opage{44} fully penetrate things with his knowledge; --- without abstraction,
no distinction.

The absolute ideality of divine thinking does indeed require that each of his
individual objects have the totality of all his other thoughts in it; but if
this is not to be a merely tautological thinking, it is on the other hand
equally necessary that the fullness of the ideas be present in the different
objects in different ways. God, then, cannot think the object \textit{A}
without at the same time abstracting from it the object \textit{B}; --- the
point of difference is the point of departure of abstraction. --- Therefore
those who emphasize God's knowledge as that which is simultaneous and intuits
everything at once must not forget that the divine thoughts form a system, and
that God, in intuiting the individual things, runs through the whole long path
of thought from the point of beginning, pure nothing, to the individual
concrete thing, so that his immediate knowledge has at the same time from
eternity been, and will always continue to be, a mediated one.

Should anyone object that we are thereby transferring unfitting
anthropomorphisms to God, and apply to us Athanasius's well-known words:
``\textit{o ludicram doctrinam aedificantem simul et demolientem}'', we must ask
him to remember well that abstraction and mediation are not separated in an
external manner from simultaneous knowledge, but that conceiving and intuiting
rather coincide. In that God conceives himself and his thoughts hang together
with one another by a bond of necessity, he has the absolute truth: ``\emph{God
wills because it is truth.}'' The more precise development of this must,
incidentally, be left to speculative theology; we have touched on these points
only in order to present the archetype of human knowledge.
% --- p. 45 ---
\opage{45} Human self-consciousness, which is an image of the divine, contains
the above-mentioned moments within itself in a formal manner. For the \emph{I}
passes continually over into the \emph{Me}, and the \emph{Me} into the
\emph{I}; the one thought constantly takes up the other, constantly dissolves
itself in the other; there must thus take place a conceiving; --- but intuition
is no less present; for when the I regards the Me, the former must always see
the latter as something positive; for were this not the case, one could not
grasp what the I is either, and there could then be no self-consciousness at
all. But that the positive and intuitable is a necessary moment, if
self-consciousness is to have the fullness of reality, is easy to see. For if
one is to be able to arrive at a conceiving, one must also be able to intuit
the real fullness, and such an intuition is not only necessary in general but
is also required in the individual concepts, in order that each individual
concept may thereby obtain its true intensity. --- That this is so is shown
also by dialectic itself; for when one wishes to develop a concept from the
concepts that precede it, one does indeed first indicate and describe it with
other expressions, but then passes over from these to its own true and
necessary designation. But the addition of such a category can have no other
ground than this: that it brings with it the intuition of a certain positive
and peculiar moment, and that this is necessary in order that the concept may
in actuality obtain individuality. --- One can therefore also demonstrate the
significance of intuition in ontology itself. Does there not, for instance,
seem to lie a far more powerful stimulus to dialectical advance in the doctrine
that pure being passes over into pure nothing than in the doctrine that the
idea of the world-soul must necessarily pass over into the idea of abstract
personality? This can hardly be denied; but the reason for it is not so much
that doctrine's greater conceivability as, far rather, its
% --- p. 46 ---
\opage{46} \setcounter{footnote}{0}%
greater intuitability. For only one who has immediately experienced the
significance of personality can recognize that the universe cannot be held fast
or conceived except by being seen as having issued from a person. Indeed, the
necessity of such a pregiven intuition will always repeat itself anew, since
God's absolute freedom has created not merely reality in general but also this
determinate world and all the determinate phenomena that meet us in it. One
must then have put oneself in possession of the whole doctrine of the Church
before one can develop what this freedom is like, and only after having
developed freedom in general can one pass on to the individual phenomena in
history. From this one sees, then, what it means that it is necessary to
believe, a demand so often made when true science is spoken of. I grasp God's
personality by carrying through the negation of my own personality; now had I
not myself a real personality, I could not have any higher concept of
personality either; but if it is to be real, it must also be religious; --- to
know God is thus to confess God, i.e., to surrender oneself to God. As,
therefore, considered objectively, everything is taken up into an existence
that is not merely thinking but also willing, so too, considered subjectively,
knowing God lies not merely in the cognition of him but also in the will's
readiness to submit to him.\footnote{``Since, then, one comes in the light of
conscience to a true cognition of an \emph{author} of all human
self-consciousness, we can determine it more precisely as the light in which
man is revealed as \emph{God's creature}; in this cognition lies the germ of all
religion and worship.''
\par H.\ Martensen, Selvbevidsthedens Autonomie, translated by Petersen,
P.\ \textit{10}.
\par ``--- --- --- since, moreover, freedom is the principle of every cognition,
in that one rational being can communicate itself to another, and live and move
in it, only through the mediation of freedom: man must freely submit to God in
order that the divine idea may really be able to penetrate into his spirit and
be inwardly united with his consciousness; man must, in other words,
\emph{believe in order to cognize}.'' Ibid., P.\ \textit{15}.}
% --- p. 47 ---
\opage{47} %
It is precisely at this point that the ontological proof of God's existence
coincides with the testimony of the religious self-consciousness; as the former
speaks to our conceiving, so the latter speaks to our intuition. One must then
have religious faith and experience in order rightly to be able to develop the
dogmas of the Christian religion; --- that one must have appropriated the
doctrine of the Church we have indicated above --- for it is through these that
the soul is filled with inner truth; it is therefore also only through them
that we become able to penetrate the dialectical moments. But this faith is
precisely genius in its highest form, i.e., religious genius.

\parag{9}{Our Knowledge Is Relative.}

From the mutual relation between conceiving and intuiting it follows, however,
that our knowledge cannot become complete before the life of God has been
revealed in its whole fullness, or, in other words, before the Kingdom of
blessedness has come. For the clarity of our concept is, after all, always won
through a greater or lesser loss of the real fullness, and the more we turn our
attention to concrete and peculiar objects, the more we are drawn into the
hypothetical sphere of freedom. Indeed, this even has its natural ground in the
relation of mutual dependence in which the disciplines of science stand to one
another, which entails that none of them can be
% --- p. 48 ---
\opage{48}\setcounter{footnote}{0}%
completely closed, because the completeness of the one depends on that of the
other. As long as the two different ways of regarding objects, the vulgar
through the senses and the speculative through thinking, only take each other
up in that they infinitely presuppose each other, that is, as long as we dwell
in this life, so long does a relative darkness continue to reign in our
knowledge (\textit{βλέπομεν γὰρ ἀρτι δι' ἐσόπτρου ἐν αἰνίγματι}). This must not
be ascribed merely to a deficiency in our consciousness; the clarity of our
knowledge is rather darkened by the unclarity of the sensuous world itself. We
cannot cognize anything before it comes forth in the objective
light\footnote{Every science is a knowledge of something that science does not
produce, but that exists independently of it. \emph{Mynster}: ``Om Begrebet
Dogmatik.'' §~\textit{4}. Cf.\ \emph{Sibbern}: ``Om Erkjendelse og
Grandskning.'' §~\textit{1}. \textit{Scimus, quia accepimus}: All cognizing is
a finding, a receiving and a beholding.}. We can therefore expect to arrive at
absolute knowledge neither in nature nor in history, and just as little in
religion itself, since it is precisely the peculiar character of this world
(\textit{τοῦ αἰῶνος τούτου}) that all the elements of life in it are split up
into detached fragments. And how, indeed, should we be able to grasp the truth
before it has itself attained an adequate form in which it can fully reveal
itself, since it is, after all, precisely the truth itself that forms
self-consciousness in us?

It is precisely this limitation of our knowledge that is designated in the
vulgar saying that there are indeed several things that men can conceive, but
that they cannot conceive everything, not the highest.

But if one now tries to refer all things to two classes, the conceivable and
the inconceivable, so that they are to stand separated, each on its own side,
one will in truth arrive at the result that either they are all conceivable, or
that
% --- p. 49 ---
\opage{49} \setcounter{footnote}{0}%
none of them is. For if we first look to the things we call conceivable, they
are surely limited and determined precisely by their relation to those we call
inconceivable. And if we proceed from the inconceivable things, these must
surely of necessity step out into the things that are conceivable, and thus
themselves become conceivable. We can adduce examples of this from dogmatics
itself: precisely by urging that human nature has preserved its powers
unimpaired, the Pelagians think to explain all the phenomena of religious life,
whereas the Augustinians, on the contrary, in order to interpret the nature of
sin and of regeneration, refer to Adam's fall as ruinous for the whole human
race, and to God's grace and justice, with the addition that God had not
predestined Adam's fall; indeed, the Calvinists go still further, in teaching
that God has predestined everything. Thus they all wish to \emph{conceive} the
phenomena of religious life, in seeking to explain them and to reconcile them
with one another; but in striving to reach this goal they all, each in his own
place, come to smuggle in an inconceivable moment, and precisely for that
reason none of them arrives at a true conceiving. It is indeed generally assumed
that the propositions of the Pelagians are easier to conceive than those of the
Augustinians and the Calvinists; but since the relation between divine and
human freedom that the Pelagians teach evidently stands in opposition to the
concept of God's omnipotence and grace, and all human actions are thus
expressions of an inconceivable human freedom, all these actions surely betray
everywhere an inconceivable moment and are therefore themselves
inconceivable\footnote{That \emph{Augustine} was an infralapsarian gives
\emph{Neander} (Allgem. Gesch. d. Rel. und Kirche. \textit{2} Band. \textit{3}te
Abtheil, \textit{p.\ 876}) occasion for the following remark: ``Aber es war dies
eine schöne Inconseqvenz, welche aus dem Siege seines religiössittlichen Gefühls
über seine dialektisch=speculative Richtung herrührte.''}. --- On the other
hand Calvin, since he
% --- p. 50 ---
\opage{50} \setcounter{footnote}{0}%
has given the premises, cannot refrain from completing his inference for the
benefit of human cognition; he therefore refers everything to a \textit{decretum
absolutum} in God that surpasses all human understanding, and so, in his
striving to conceive the relation, arrives at a result that is altogether
inconceivable.

We thus find that it does not depend on the nature of the matter itself where
this birthmark of human cognition, if we may so call it, is to be found, or how
it is to be concealed, but that it rather has another ground, whether we now
wish to call this subjective arbitrariness or religious tact\footnote{Hence the
objections that have been advanced even in our own time against Calvin's
consequences in this respect may also well deserve a hearing. Thus
\emph{Guerike} says, Kirchengesch. P. \textit{1020}, Not. \textit{216}: ``Die
heilsam inconseqvente Vorstellungsweise der Concordienformel im Verhälltnisse zu
der selbst gotteslästerlich conseqventen Calvins ist also kurz diese: Seligkeit
und Hoffnung ist nach beiden nur ein Werk der göttlichen Gnade, Verdamniß eine
Folge der eignen Schuld. Während aber nun die Concordienformel demüthig bei
diesem letzteren Satze stehen bleibt . . . . . argumentirt Calvin weiter etc.''
But when Guerike saw fit to add: ``Gott also hat selbst die Sünde der Menschen
gewollt. Diese letzte, nothwendige, ob auch noch so bemäntelte Schlußfolge
. . . . .'' he seems, by adding this new inference, himself to have sinned
against the pious reverence for the holy that he praised in the Formula of
Concord, and I do not doubt that Calvin, were he still alive, would answer him
with the same words with which he answered the irreligious striving of his time
to conceive the divine mysteries: ``\textit{Arripiunt subito, fatcor, profani
homines in prædestinationis materia, quod vel carpant, vel cavillentur, vel
allatrent, vel subsannent. At si nos absterret eorum procacitas, celanda crunt
præcipua quæquc fidei dogmata. (Instit. christ. rcl. De prædest cap.
XIV.).}''}. It is therefore evident that one cannot draw that separation in an
external manner.

% --- p. 51 ---
\opage{51} %
If the mystery must thus pervade all moments once it is present at a single
point, we can also conclude conversely that our knowledge must be able to
extend itself to all objects as soon as it is present at merely a single point;
for were it bound to succumb to the mystery, it could not exist at all, since
the mystery asserts its power everywhere; in other words: the limitation of our
knowledge is not metaphysical, but empirical. The laws to which our cognition
is bound hold for all eternity, and this accords, too, with the nature of truth;
for of truth the same holds as of God, that it is not far from any one of us
(\textit{Ἡμῖν δὲ ἀπεκάλυψε ὁ θεὸς διὰ τοῦ πνεύματος αὐτοῦ} \textit{1} Cor.
\textit{2, 10}).

Human consciousness, therefore, has only to hold unshakably fast to the divine
and real object and diligently strive to penetrate its moments; it will then
find within itself what it seeks (\textit{quod petis hic est}), and the bounds
of absolute cognition will thus expand more and more for it from within.

\parag{10}{The Relation between Dogmatics and Sacred History.}

In order to develop more fully the relative limitation of our knowledge, we
must consider more closely the relation between speculative dogmatics and
speculative history. --- Our knowledge is absolute, partly because
consciousness is perfectly
% --- p. 52 ---
\opage{52} clear to itself about its content, partly because the real moments
it comprises are altogether universal and universally valid. From this it
follows that dogmatics is the most real of all the sciences. For the concepts of
ontology are altogether abstract, and the knowledge that history imparts is far
too hypothetical. Therefore that positive moment which, as we have remarked
earlier, is always necessary in order to call forth the negative stimulus of
dialectic, does not in the domain of history attain any true intuitive clarity,
but often has the effect that the dialectical sting, too, is blunted, because
it is itself surrounded by sensuous obscurity. Dogmatics, by contrast, does take
account of positive reality, but not of its more special modifications, and
precisely thereby brings it about that the contradiction always comes forward,
so that the concept is preserved. But this peculiarity of dogmatics calls forth
a certain opposition. For if the individual phenomena are effaced, the more
general reality itself disappears as well; one must then take refuge in
everyday experience, in order with its help to gain from every quarter the
missing fullness, --- but since every such experience is a negation and
limitation of science, everyday experience must transform itself into a
speculative experience, i.e.\ history must become science.

The two different kinds of history each require a faith of their own; the
faith necessary in vulgar history is mediated through the spirit of this world,
which considers only trivial actuality; the faith of scientific history, by
contrast, is mediated through the divine spirit, which posits ideal actuality.
About some events the two kinds of faith agree; about others, or rather about
most, they cannot, on the other hand, come to agreement.

We have touched on these relations here in passing, although they will later
become the object of a closer consideration, in order to make it evident that
dogmatics is a necessary
% --- p. 53 ---
\opage{53} presupposition for all historical knowledge. For one cannot conceive
the deeds of the human race without having grasped the idea of freedom in its
truth; but in order to understand the freedom of man it is necessary to have
understood the freedom of God; one can therefore conceive history only once one
has appropriated the idea of religion and of revelation. So far, then, is
sacred history from being violated when we seek to conceive it through
philosophy, that it is much rather profaned whenever the phenomena are
presented in an external manner, torn loose from the idea of God. Since, then,
vulgar history is always profane, whereas speculative history is sacred, sacred
history must be speculative.

\par\vspace{0.6\baselineskip}\centerline{\rule{0.42\textwidth}{0.6pt}}\vspace{0.6\baselineskip}

\deel{Second Part.}{Sacred History.}

\parag{11}{Catholicism.}

Sacred history requires that the invisible divine nature present itself in
visible phenomena; it thus sublates the opposition between the sensuous sphere
of created things and the spiritual sphere of divine existence, in that
spiritual objects come forth sensuously so that sensuous objects may be
explained in the light of spirit; it likewise sublates the antinomy between the
freedom of the Creator and that of the creature, in that God lays upon himself
the bonds of finitude in order to give created freedom occasion to choose and
to act. While it thus entails an immediate identity of the
% --- p. 54 ---
\opage{54} supernatural and the natural realm, it must nonetheless just as
surely let an opposition between the two come into view. For if God's history
is to be sacred, God must be honored as the Holy One; and if God is to be
honored as the Holy One, he must be exalted above created nature in infinite
majesty. But sacred history removes this self-contradiction by referring to an
opposition between different times and places. In some places and at some
times, namely, the life of the world, so to speak, does not accord with the
life of God and shows itself torn loose from it, but in other places and at
other times the two come forth in immediate reconciliation. On this account no
perceptible discord arises between vulgar and speculative experience; for in
the world of sense we behold everywhere, with the eyes of faith, the
supernatural beings, and as often as the laws of the world are interrupted by a
higher order of things, as soon as mysteries make their appearance, everyday
consciousness willingly surrenders without any philosophical tumult; but
scarcely has the natural order returned before it again enters into its old
rights. Wherever there is a revelation, this reciprocal relation between the
two spheres always comes about; in the absolute, the Christian, religion it
therefore appears in the highest measure.

When the apostles saw Christ, of whose appearing the Evangelist solemnly
testifies: ``And the Word was made flesh, and dwelt among us, and we beheld his
glory, the glory as of the only begotten of the Father, full of grace and
truth,'' they in truth beheld neither a merely sensuous figure nor an abstract
spiritual idea, but they beheld in him the divine and the human nature in
perfect identity, so that his history could with full right be called sacred
and divine. And when he had gone away to the Father and had sent out the
Spirit, the life of God so pervaded the apostles and all the faithful,
% --- p. 55 ---
\opage{55} \setcounter{footnote}{0}%
the power to work miracles was present in such full measure, the angels of
heaven revealed themselves so clearly to all, the divine counsels made
themselves known so plainly in the intrigues of the adversaries (\textit{τῶν
θεομάχων}), that we also with the greatest right call apostolic history a
sacred history. --- Even after the death of the apostles, when the machinations
of this world increased and the implacable enemy of Christianity \emph{devised
new intrigues in order, under the very cover of the Christian name, to lead the
thoughtless astray}\footnote{Cf.\ \textit{Cyprian: de unitate ecclesiae.}}, the
source of revelation seemed unable to be exhausted; the more men mingled truth
and falsehood, the greater was the need to have God visibly present in holy
persons, in infallible doctrines, in external criteria of the congregation of
true Christians, in short, in every way. Thus the Church, in the tribulations
and darkness of this world, \emph{taught by the Holy Spirit, who each day
imparted to it the whole truth}, developed its excellent institutions. The
priests came forth as God's representatives, a visible ruler established the
visible unity, so that Christ's bride should not have her fair
\emph{countenance marred by any spot or wrinkle.} In the mystical nimbus of the
sacraments one sees Christ himself, daily the number of miracles grows, and as
the divine nature ever anew descends in the form of flesh, heaven opens itself
to the faithful, and all the barriers of everyday experience are overstepped:
one hears the myriads of angels praise the majesty of the Almighty, and,
surrounded by the host of victorious martyrs, the pious \emph{Mother of God}
smiles upon the praying pilgrims, while the inhabitants of paradise rejoice and
the monsters of hell tremble.
% --- p. 56 ---
\opage{56} \setcounter{footnote}{0}%
Thus the Church holds at once earth and heaven, the present and the future, the
natural and the supernatural in its embrace, so that for Catholics the history
of the Church is, in the fullest sense of the word, a sacred history.

\parag{12}{The Fundamental Contradiction of Catholicism.}

But this glory of Catholicism contains a fundamental contradiction. It is
always the nature of essence to create its phenomenon; since, then, the Holy
Spirit must necessarily create forms and traditions, one rightly says:
``\emph{Where the Spirit is, there must the Church also be}'' (\textit{Ubi
spiritus, ibi ecclesia}).

Now since the institutions of the Holy Spirit are infallible, since the
individual propositions of tradition are words dictated by God and therefore
unchangeable, the object of truth is absolutely present, so that the invisible
and the visible Church perfectly coincide. If the human spirit, therefore,
wishes to grasp the truth, it need only throw itself entirely into the arms of
the Church; its whole knowledge consists in the cognition that the Church has
the truth, and everything that is true it will find expressed in the precepts
and doctrines \emph{of the holy Mother}\footnote{Baur: Gegens. d. Kath. und
Protest. \textit{p.\ 501---504.}\\
Sollte jener freiere Begriff der Tradition seine Gültigkeit haben, so dürfte das
eigene Urtheil des Einzelnen nicht schlechthin ausgeschlossen werden, es müsste
jedem überlassen bleiben, die Unterscheidung, die zwischen Form und Inhalt,
zwischen dem Wesentlichen und Unwesentlichen, dem Allgemeinen und Besondern
gemacht werden soll, auch wieder auf alles dasjenige anzuwenden, was die Obern
der Kirche als allgemeine Lehre der Kirche geltend machen. \textit{p.\ 505.}}.

% --- p. 57 ---
\opage{57} \setcounter{footnote}{0}%
But in this one forgets to carry out the negation and restoration of the
phenomena of which we spoke above (§ 4); a certain ecclesiastical sensualism
then sets in, in that the unchangeable forms of the Holy Spirit are impressed
upon the religious mind as upon a \textit{tabula rasa}; the mind becomes not a
free organ for the truth, but a form for the forms of the Holy Spirit, and
therefore regards the phenomena from outside through a vulgar empirical
knowledge and acquiesces in empty, subjective opinions\footnote{It (i.e.\ the
Catholic Church) does not derive its authority from its Christianity, but its
Christianity from its authority.\\
\textit{Dr.} H. N. Clausen, Cathl. og Protest \textit{p.\ 10. 11.}}; thus the
proposition comes to hold: ``\emph{Where the Church is, there is the Spirit}''
(\textit{ubi ecclesia, ibi spiritus}). If the ecclesiastical institutions are
now to be maintained, the Holy Spirit must have certain immediate organs, i.e.\
the Church must have priests. Since these neither descend from heaven nor are
born as such, they must be chosen from the midst of the laity; hence the
transition from the class of the laity to the priestly estate is of thoroughgoing
significance. Since all laymen are struck with equal blindness, so that, when
questions arise about improvements or corrections in ecclesiastical doctrines,
no account is taken of differences either in mental gifts or in education,
there can in this respect be no question of different degrees of enlightenment.
Before ordination the Spirit is present only immediately; only through lawful
ordination does the Holy Spirit step in and stamp the one consecrated with the
mark of infallibility (\textit{character indelebilis}). It is thus indeed the
Spirit that makes him a priest, but there is no other mark of the Spirit's
presence than ordination, i.e.\ no one becomes a priest because he has the
Spirit, but
% --- p. 58 ---
\opage{58} \setcounter{footnote}{0}%
everyone who is a priest has the Spirit\footnote{\textit{Quodsi rasuram,
unctionem et longam tunicam tantum possunt ostendere pro suo sacerdotio,
permittimus illis gloriari. Lutherus de instituendis ministris.} Cf.\
\textit{Dr.} F. C. Baur. Der Gegensatz des Chatol. und Protest, \textit{P.\
475.}}. Nor is it through the priest's teaching that the truth of the Spirit
dwelling in him makes itself known; for although each individual priest is in
possession of the Spirit, they all nonetheless remain weak mortals who can sin
and go astray. One cannot, then, establish the infallibility of the individual
priest except by the testimony of the Church; the priest therefore owes the
Church the same obedience as the layman owes the priest. But through whose
mouth can the Church speak, if not through that of the priests? Therefore each
individual priest must be subject to the united opinion of all the others:
\emph{the majority is decisive.} For if the one were to convince the other of
the truth by the truth's own power, one would surely have to negate the
infallibility of the form as such; but once such a negation and free inquiry
had asserted itself at a single point, it would penetrate to all parts of the
Church, the lesser along with the greater, and then the very idea of
Catholicism would disappear. Every inquiry rather rests on external traditions,
and in such an inquiry each person's immediate opinion now comes into play, so
that it is always something accidental which opinion will get the plurality on
its side\footnote{Denn ich glaube weder dem Pabst noch den Concilien allein
nicht, weil es am Tag und offenbar ist, daß sie oft geirrt haben, und ihnen
selbst sind widerwärtig geweßt. Luther.}. It contributes but little to the
removal of this difficulty that Catholicism has instituted different grades in
the hierarchy of the priestly estate and of ecclesiastical
jurisdiction\footnote{Cf.\ Clausen Cathol. og Protest. \textit{P.\ 91} f.}; the
possibility of the
% --- p. 59 ---
\opage{59}\setcounter{footnote}{0}%
distortion of the truth is thereby only ramified the further. For what does it
help that the opinions of the presbyters conform to those of the bishops, if
the bishops contradict one another among themselves? It then becomes necessary
to have a visible ruler to whom one can turn in order to hear the living voice
of the Holy Spirit. We have thereby arrived at the well-known disputes between
Episcopalists and Curialists\footnote{The same: Cath. og Prot. Det
curialistiske System \textit{P.\ 14---26.} Det episkopale Kirkesystem \textit{P.\
43---52.}}. But that episcopalism is the system that strives to keep to the bad
middle way in the Catholic Church is well known.

The glory of Catholicism, like its fundamental contradiction, culminates in the
person of the Pope; on the one hand the Pope is a weak and sinful man, on the
other hand he is God's representative, in whom the fullness of the living
Spirit takes up its dwelling; but who is now to decide when he speaks as Pope,
driven by the Holy Spirit, and when he acts according to his own subjective
arbitrariness? The councils, which were supposed to have the whole fullness of
spiritual life within them, are surely subordinate to the Pope. Protestant
science, therefore, has long since proved\footnote{Clausen \textit{l.\ l.\ P.\
58.}\\
Curialism forms in the Church a theocratic monarchy so unlimited that in
actuality it must pass over into despotism, for it combats not only external
relations and connections, but also the spiritual personality of every
individual human being, every expression of independent thought and will; and
in return the Church relieves its citizens who are obedient in faith of all
responsibility of conscience, and grants them undisturbed security and peace;
but it is the peace of death on the grave of slain enemies.}, that the objective
idea of Catholicism --- precisely where it has reached the highest summit of
objectivity --- has passed over into, and disappeared in, an arbitrary
subjectivity, whose whole infallibility consists in its declaring itself
infallible. Therefore
% --- p. 60 ---
\opage{60} \setcounter{footnote}{0}%
all Catholic institutions, although they are of such a nature that they can be
forms for the Holy Spirit, can at the same time, when they are appropriated in
an external manner, be held fast in and for themselves.

The Church thus often presents empty forms, and the falsity in Catholicism lies
precisely in this, that whoever is in possession of the forms is also assumed
to possess the Spirit. Under such a relation it is no wonder that the external
phenomena have often been forsaken by spirit and truth to such a degree that
not even the semblance of decency has been preserved, and that Catholicism has
often had to blush at the testimony of history, although it is precisely in
this testimony that it has put its trust\footnote{\textit{Baronius}
``\textit{significavit suo Jac. Sirmondo, nihil se æque reformidasse, quam
pervenire ad hoc tempus (i.e.\ magni Schismatis Occidentis) de quo, quid
statuendum, non esset libere pronuntiaturus}'' \textit{l.\ l.\ p.\ 49.}}.

\parag{13}{The Transition from Catholicism to Protestantism.}

In the attempt to present the absolute object of truth to the beholder,
Catholicism constantly strives to bring human self-consciousness to turn its
gaze away from itself toward external objects. But from this it follows that
the objects in the kingdom of God, by being fixed in concrete images, come to
be ranked alongside earthly objects, so that the holy is profaned, while on the
other hand the descent of heavenly beings into earthly forms entails that
temporal objects are accorded divine holiness, because they enclose the divine
within themselves, so that what is in itself unholy becomes an object of
% --- p. 61 ---
\opage{61}%
\setcounter{footnote}{0}%
worship\footnote{Cf.\ Hegel, Geschichte der Philosophie. \textit{3}ter
Band, P.~\textit{253} f.}. When this barbarous confusion of opposed elements
had become plainly manifest in the theocracy of government, in the
scholasticism of science, in the significance ascribed to good works,
self-consciousness, which had long been banished from its own interior, began
by various paths to return to itself\footnote{Hegel \textit{l.~l.}
P.~\textit{267}. Die andere Seite ist, daß das Ewige, was an und für sich wahr
ist, auch erkannt, aufgefasst werde durch das reine Herz selbst; der eigene
Geist mach sich für sich das Ewige zu eigen. Das ist der lutherische Glaube ohne
anderes Beiwesen (die Werke, wie man es nannte).}.
If one now investigates how it was possible that so great a split could arise
in the Church, one will find that the only means of thoroughly reforming the
Church was to negate the phenomena in order to penetrate to the essence, to
draw a sharp separation between the idea and the form, between the invisible
and the visible Church, and in consequence to enter into open struggle against
the hierarchy. For where consciousness can do nothing but become a slave under
a foreign authority, no true unity of subject and object can ever emerge; one
therefore sees in such a positivism a constant wavering between obedience and
rebellion. The obedience is blind so long as the subject can be bound to seek
the substance outside itself; rebellion arises when this substance is not seen
to constitute the subject's own nature, but, as a foreign power, only strikes
fear; therefore the \emph{faith} of the Catholics is a relation of
\emph{bondage} and always contains a certain mistrust.

If the true is to be beheld in the light of truth, objective reason must
necessarily pervade subjective reason; otherwise we cannot be as certain and
firm in possession of the thing itself as we are in possession of our
% --- p. 62 ---
\opage{62} \setcounter{footnote}{0}%
own self-consciousness. Only where this relation is immediately present is there
a true faith. --- With such a faith, so far is subjective life from being
hampered by the objective bonds, that the subject rather surrenders itself with
full confidence to the object, in the cognition that it is precisely from the
objective that it draws its true nourishment. This is the Protestant faith,
which breaks the chains of positivism and founds an actual freedom. What
deceptive phenomenon could mislead consciousness, what external power could
shake the strength of the soul, once one has found this palladium of freedom?
Therefore the man of God, whom we praise with grateful hearts as the prince of
Protestantism, stood unshakably firm; girded with the belt of truth, he
dispersed the fog of traditions and fables; protected by the shield of faith,
he quenched the fiery darts of ecclesiastical tyranny\footnote{Es thut dir wehe
im Herzen mein fröhlicher großer Muth. Ich bin aber, und will, ob Gott will,
auch bleiben, gegen dir und Ecken, Pabst und eurem Hauffen, auch dem Teufel, mit
Gottes Hülfe, in einem beständigen, hochmüthigen, unerschrocknem Geist, und euch
trotzen und verachten \ldots\ Luthers Schutz=Schriften v. An.~\textit{1524}
wied. Hier. Emsern u. Paris. Theologos etc.}: with the drawn sword of the
Spirit, the word of God, he attacked the unholy idols of
superstition\footnote{Hörst du es, Pabst, nicht der Allerheiligste, sondern der
Allersündigste, daß Gott deinen Stuhl vom Himmel auf's schierste zerstöre und in
Abgrund der Hölle senke. Cf.\ Pfizer L.~L. P.~\textit{161}.}. From him went
forth the first crisis, --- the false beauty vanished, and the history of human
errors was drawn forth into the light.
% --- p. 63 ---
\opage{63} \setcounter{footnote}{0}%
\parag{14}{Faith and Holy Scripture.}

And yet this first crisis did not dissolve the bond between the natural and the
supernatural realm; faith did indeed exchange the external elements for the
inner moments of spirit, so that the motley splendor of images disappeared; but
nonetheless the good angels still continued to look down upon men with their
fair countenances through the clouds of heaven\footnote{Luther: Daß man also
Cherubim verstehe für Engel, die da erscheinen, nicht mit einem runtzlichten,
noch traurigen Angesichte, sondern mit einem fröhlichen Geberde, und vollem
feistem Angesichte \ldots\ Ausleg. des \textit{3}ten Capit. des \textit{1}sten
Buch Mose.}, the devil, \emph{this caricature of God}\footnote{``\textit{Simia
dei}'' an allusion to Luther's words about the devil: ``der Affe
Gottes.''\hfill Translator's note.}, continued under countless transformations
to tempt the children of men\footnote{Luther: Er (der Teufel) ficht mich selbst
offtmahls so gewaltig an, und überfället mich so hefftig mit schweren und
traurigen Gedanken, daß ich meines Herrn Christi gar vergesse \ldots\ ist er
dennoch heimlich bei uns, schleicht Tag und Nacht um uns her \ldots\ Ausleg. d.
Epistel a.~d. Galater Cap.~\textit{3}.}.
Nor did the freedom with which the truth of the phenomena was negated go so far
that no infallible form of the truth should have remained; --- the more the
corruption of the existing Church gave offense, the more gladly one sought out
its original pattern: the bolder one was in rejecting the divinity of the
traditions, the more firmly one attached oneself to the Gospel in its original
form. The Catholic Church had in a special way exalted the apostolic age and
acknowledged its writings as canonical; since, then, the religious
consciousness had beheld the truth in these books in unadulterated form, one
eagerly seized upon \emph{Holy Scripture} as the true pillar of faith.
% --- p. 64 ---
\opage{64} \setcounter{footnote}{0}%
Yet this is not to be understood as if, in urging Holy Scripture, one merely
wished to set up a petrified norm in a positivistic manner, so that
Catholicism, which had long prided itself on the agreement of its doctrine with
that of the apostles, could have continued this self-praise undisturbed for a
long time yet, had not some theologians happened by chance upon Holy Scripture,
read it quite carefully, and undertaken a comparison. That was impossible; for
had they not already harbored great doubt about many of the propositions of
Catholicism, they would surely have expounded Scripture according to the rule
laid down by the Church. The truth is present everywhere and at every time;
therefore the Spirit, which never forsakes its Church, revealed itself at once,
as soon as the false veil that had covered it had been cast off. --- One must
therefore bear well in mind that when the champions of the Protestants
constantly appealed to Scripture as the highest testimony of truth, they could
do so solely and only because the truth itself was the witness of Scripture, so
that they honored the spirit in it and not its letters. --- When the Reformers
had noticed the disagreement of the traditions with
Scripture\footnote{Was Christum nicht lehret, das ist noch nicht apostolisch,
wenn es gleich S.~Petrus oder S.~Paulus lehrete. Wiederum was Christum prediget,
das wäre apostolisch, wenns gleich Judas, Hannas, Pilatus und Herodes thät.
Luther: Vorrede auf die Epistel St.~Jacobi. Cf.\ with this: Ein Stück der
Vorrede auf das N.~Testament: Welches die rechten und edelsten Bücher des Neuen
Testaments sind. An.~\textit{1521}.} and found that they had distorted rather
than developed the Christian idea, it lay in the nature of the matter that they
had at once to ascribe to Scripture \emph{clarity} and \emph{completeness}
(\textit{perspicuitas et sufficientia})\footnote{Da erdenken sie eine neue Lüge,
finden Degen und Spies und dergleichen Narrenwerk, und sprechen: Die Schrift sey
so finster, daß wir sie nicht mögen verstehen ohne der heiligen Väter Auslegung
\ldots\ Wenn sie denn nun einen Spruch der Väter wieder mich aufbringen, so
läuten sie alle Glocken, schlagen alle Drummeln, und schreien feindlich, sie
haben gewonnen. Luther: Von dem bleyern Degen Bocks Emsers.

And the more everything was referred to the substantiality of faith and not to
the contingency of works, the easier it was to show the completeness of
Scripture. Therefore he said, when the Catholics referred to John
\textit{21, 25}: Es ist nicht alles geschrieben, was wir thun sollen. Danck habt
ihr guten Gesellen, ihr wisset der Schrifft Auslegung wohl zu geben.
S.~S.}; but one must note well that they, brought up as they
% --- p. 65 ---
\opage{65} \setcounter{footnote}{0}%
were, in the Catholic Church and devoted to it with their hearts, were so
thoroughly permeated by the Church's teaching in its entirety that, as their
interpretation too sufficiently proves, they did not so much draw the
fundamental dogmas of faith out of Scripture as recognize them as deposited
in it; and although, in opposition to the allegories the Catholics employed
so as to be able to interpret every passage to their own advantage, they
insisted on the literal meaning and on a comparison of parallel passages of
Scripture, they nonetheless wanted Scripture explained in the light of the
totality, so that the Old Testament could not be unfolded except with the
help of the New Testament\footnote{Die heilige Schrifft dringet
vielmehr auf den Sohn, als auf den Vater: denn die ganze Schrifft ist um des
Sohnes willen geschrieben: darum sind auch im alten Testament mehr Sprüche
oder Zeugnisse vom Sohne als vom Vater. --- Luth. Schrifft. \textit{III.~4.}
Denn da steckt es, da bleibt es, wer diesen Mann, der da heist Jesus Christus
Gottes Sohn, den wir Christen predigen, nicht recht und rein hat, noch haben
will, der lasse die Bibel zu frieden, das rathe ich, er stöst sich gewisslich
und wird, je mehr er studiret, je blinder und toller.

Wenn es nun sollte Wünschens und Wählens gelten, entweder, daß ich
S.~Augustini und der lieben Väter, das ist der Apostel Verstand in der
Schrifft sollte haben, mit dem Mangel, daß S.~Augustinus zuweilen nicht die
rechten Buchstaben oder Worte im Ebräischen hat, wie die Jüden spotten
\ldots\ ich liesse die Jüden mit ihrem Verstande und Buchstaben zum Teufel
fahren, und führe mit S.~Augustini Verstande ohne ihre Buchstaben zum Himmel.
Auslegung der letzten Worte Davids \textit{2}~Sam. \textit{23, 1---7}. Schr.
\textit{IV. 303}.}. Nor was this to be wondered at, for despite the
% --- p. 66 ---
\opage{66}\setcounter{footnote}{0}%
differences of the times the Holy Spirit had been the same, and it could in
truth have had Christ in mind, although it addressed its words to the
Jews\footnote{Thus Luther himself cannot refrain, quite contrary to his own
principle, from sometimes adding an allegorical meaning, e.g.\ Genes
\textit{c.~9} on Noah: Darum willst du es ohne Gefahr deuten, so deute es auf
den Herrn Christum. Denn wie \emph{Noah den Weinberg pflanzet und des Weins
trincket,} davon \emph{truncken wird} und einschläfft, \emph{und bloß in der
Hütten lieget,} und wird von diesem verlachtet, aber von andern
\emph{zugedecket:} also ist es auch Christo ergangen.}.

Here, then, is where the negation and restoration of the phenomena of which we
spoke above has its beginning. For when Luther undertakes to read Scripture
and to judge the greater or lesser truth of the traditions, the Spirit of
truth is present in his consciousness, so that he does not judge merely by a
subjective estimate. But how did he obtain this Spirit? Not through any
peculiar inspiration, like the \emph{Enthusiasts}\footnote{Thus he writes to
Amsdorff about the new prophets: Durch die Zwickauer Propheten lasst Euch nicht
so schnell bewegen! Ihr habt ja Zeugnisse der Schrifft, die Euch sicher
machen, daß Ihr nicht sündiget, so Ihr sie hinausschiebet und zuerst die
Geister prüfet, ob sie aus Gott sind \ldots\ Mir wahrlich ist das dem ersten
Anblick nach gar verdächtig, daß sie sich der Gespräche mit der göttlichen
Majestät rühmen..}; but through the total impression he received in his
persistent study of the traditions and of Scripture. A true reconciliation of
the objective and the subjective moment had thus taken place; that he asserted
his subjective freedom is seen from the fact that no opinion given from
outside produced a blind submission in him; but objective reality was
preserved just as well, for he recognized that only through the appropriation
% --- p. 67 ---
\opage{67}\setcounter{footnote}{0}%
of the objective phenomena can one come to the truth\footnote{So Gott
sein Heiliges Evangelium hat aus lassen gehen, handelt er mit uns auf
zweierley Weise. Ein äusserlich; das andere mahl innerlich. Aüsserlich handelt
er mit uns durch mündliche Wort des Evangelii und durch die leiblichen
Zeichen, als da ist, Tauffe und Sacrament. Innerlich handelt er mit uns, durch
den H.~Geist und Glauben \ldots\ also, daß er's beschlossen hat, keinem
Menschen die innerlichen Stücke zu geben, ohne durch die äusserlichen Stücke.
Wieder die himmlischen Propheten. Schrifft. \textit{XIX., 186}.}. At this
standpoint, then, a true transition is also possible from the pure
universality of spirit to its particular concrete determinations; for when
the universal idea of truth has by its light negated the traditions and the
passages of Scripture in their particularity, it will always strive anew to
restore the negated forms; and when Luther insisted that Holy Scripture
should be the norm for the unfolding of the opinions of the Fathers, and
faith in Christ a necessary condition for the unfolding of the individual
passages of Scripture, it was not the idea of Christ torn loose from its
particular determinations that he insisted on; this careful examiner of the
very least utterances of Scripture\footnote{Ach es ist Dollmetschen ja nicht
eines Jeglichen Kunst, wie die tollen Heiligen meinen; es gehöret dazu eine
recht fromm, treu, fleissig, furchtsam, christlich, gelehrtes, erfahren, geübt
Herz. Cf.\ Pfizer S.~S. P.~\textit{843---852}.} never made Christ the object
of his thought without at the same time looking to all the individual
passages of Scripture that portray Christ for us\footnote{Also weiset er uns
allenthalben in die Schrifft, warum das? Darum daß du kannst bewahren den
Christlichen Glauben. Es ist mir wohl selbst begegnet, wenn ich das Wort habe
fahren lassen, daß ich Gott, Christum und alle mit einander verlohren habe.
\emph{Es ist auch kein leichterer Weg alle Artickel des Glauben zu verlieren,
denn ausser der Schrifft daran zu gedencken.} \textit{III.} Predigt über das
Evangelium am Ostermont. Luc. \textit{24, 13, 25}, zu Coburg geh.
\textit{1530}. Schrifft. \textit{XV. 631}.}.
% --- p. 68 ---
\opage{68} \setcounter{footnote}{0}%
As the moments thus flow infinitely over into one another, we see the
objective spirit really present in Luther, we see it make his spirit into an
organ for itself, inasmuch as he, as an \textit{ἄνθρωπος πνευματικὸς}, tests
everything \emph{through faith}\footnote{Ich will mich hören lassen, und wie
St.~Petrus lehret, meiner Lehre Ursach und Grund beweisen für alle Welt, und
sie ungerichtet haben von jedermann, auch von allen Engeln. Denn sintemal ich
ihr gewiß bin, will ich durch sie, euer und auch der Engel (wie St.~Pauius
spricht) Richter seyn, daß wer meine Lehre nicht annimmt, daß der nicht möge
selig werden.}.
It therefore lay in the nature of the case that preachers of the Word who
came forward, driven by the unshakable conviction of its truth, had to
declare their opponents either deceived or deceivers, even though they
thereby seemed to stand in open contradiction with themselves, since they
themselves were weak mortals and were preaching precisely against the
infallibility of any mortal. This contradiction was resolved, however, by
their sharply distinguishing between the teaching insofar as it was their own
and the objectively grounded truth that had its immediate and original
expression in Holy Scripture\footnote{Nicht also du Narr, höre und laß dir
sagen. Zum Ersten bitte ich, man wollt meines Namens schweigen und sich nicht
Lutherisch sondern Christen heissen. Was ist Luther? Ist doch die Lehre nicht
mein. Luthers tr. Ermahn. a.~all. Christen Schrifft. \textit{XVIII.}
P.~\textit{289}.}. For the sake of their hearers' welfare they had to make
this separation, which was as far from weakening the force of their own
conviction as it was from having arisen out of a feigned modesty; for even if
their words were the most perfect expressions of the truth, the hearers could
nevertheless not make what they had heard their true property except by
themselves making it the object of an investigation; nor was the speaker's
conviction
% --- p. 69 ---
\opage{69} \setcounter{footnote}{0}%
of the truth of his teaching thereby shaken, since he hoped precisely,
through a thorough and open investigation of the matter, to be able to
convince others too of the truth of that at which he himself had arrived by
the way of conviction, and to be able to trace back whatever difference of
opinion might possibly remain to a merely subjective peculiarity in the
contending parties; unity, then, was sought not in everyone's holding the
very same opinions, but in the fullness of truth satisfying everyone's
religious need. Truth, then, exists only for the testing
consciousness\footnote{Für das erste (viz.\ werden die Geister
geprüfet) durch ein innerlich Urtheil, da ein jeder Christ durch den Heiligen
Geist und Gottes Gnade, für sich und sein Gewissen, also erleuchtet ist, daß
er aufs allergewisseste schliessen und urtheilen kann von allen Lehren.

Und diese Gewißheit gehört zum Glauben, und ist von nöthen einem jeden
Christen, ob er gleich nicht ein Prediger oder in öffentlichem Amte ist. Luth.
w. Crasm. Schr.: \textit{XIX}.}, is not content with dead propositions but
desires living subjects; it therefore appears in a peculiar shape in each
individual's consciousness; there are, then, as many truths as there are
subjectivities. But so as not itself to vanish through subjective
particularism in the dispersal of its elements, it gathers its confessors and
summons them to mutual communication and discussion; thus the religious
colloquy is called forth. Hence the age of the Reformers may rather be called
the age of colloquies than the age of councils. Councils suit hierarchs, who,
although they lay claim to a priestly infallibility, treat truth in the
manner of laymen, by comparing opinions and counting votes in an external
way; colloquies, on the other hand, suit laypeople, who lay aside the
borrowed nimbus of infallibility and, in the name of objective truth, set
faith over against faith, subjectivity over against subjectivity, in order
through this ordeal by fire to come
% --- p. 70 ---
\opage{70} \setcounter{footnote}{0}%
to certainty as to who are the true priests\footnote{It was therefore
natural that the Pope did not hasten to convene a council for the
Protestants' sake; for he wanted to command, whereas they wanted to dispute
and ``to bring their case before a general, free and Christian council.''}.
But in such religious colloquies two moments are to be observed: on the one
side, namely, each of the disputants has his own spiritual fullness, to which
his consciousness is \emph{inwardly} bound in a universal way; on the other
side he turns \emph{outward}, in that he wants to describe and defend this
subjectivity of his; each therefore endeavors to determine his subjectivity
more precisely through particular opinions and to secure it through
particular grounds, so that they may thus correct one another and arrive at
the perfect truth. --- But if all difficulties are to be removed along this
path, then, in the first place, the mystical and obscure side that is usually
found in every human being's interior must come forth clearly before
consciousness; for only when one is perfectly clear to oneself can one give
others a perfect portrayal of one's subjectivity, --- but this will hardly be
able to happen before the kingdom of blessedness has come ---; and next, the
utterance of subjective truth must not be a merely subjective confession, but
must be traced back to an inner necessity, as the true method demands, ---
but method was foreign to the spirit of this age.

Those sermons and religious colloquies aimed more at becoming conscious of
their mutual spiritual kinship than at arriving at true conceptual
clarity\footnote{Therefore at the council in Marburg, after they had
discussed everything without coming to agreement, from the metaphysical
properties of the body down to the most insignificant particles of
Scripture, Luther pushed away the hand that Zwingli, with tears in his eyes,
held out to him, with the words: ``Ihr habt einen andern Geist, denn
wir!''}. It was the feeling of this spiritual kinship, awakened by the
Master's mighty voice, that
% --- p. 71 ---
\opage{71}\setcounter{footnote}{0}%
founded the Church to which we confess; and so quickly did this feeling
penetrate that, although the subjective Luther, if we may so put it, left the
council at Worms as an outlaw, scarcely \textit{10} years had passed before
the objective Luther returned to the council at Augsburg, strengthened by the
rich blessing of the Reformation. And just as Luther constantly appealed to
the objective testimony of God's Word, although he held unshakably fast to
his faith, so too the Augsburg Confession arrogates to itself no authority
beside or above Scripture, but on the contrary derives its whole validity
from it\footnote{\textit{Offerimus in hac religionis causa
nostrorum concionatorum et nostram confessionem, cujusmodi doctrinam ex
scripturis sanctis et puro verbo Dei hactenus illi in nostris terris
\ldots\ et in ec clesiis tractaverint. Præf. ad Cæs.}}.

\parag{15}{The Symbols. Lutheran Scholasticism.}

The external organism that the objective spirit forms for itself from within,
it breaks down again from within in the course of time. So long as it is
still occupied with this, the basic features of the organism remain the same,
and all is calm, although the form grows old and the vigor of youth withers;
but once time has carried its fruit to term, then it is seized by labor
pains, then the elements of life are shaken, and amid violent pangs of birth
the new shape of things breaks forth. The spirit of Protestantism, which had
not been willing to bear the Catholic yoke, had given subjective freedom such
slack reins\footnote{Ulrich von Hutten's words to Luther: Wache auf die edle
Freiheit! were a cry that had scarcely echoed across Germany before it called
the most diverse combatants to arms.} that even its champions
% --- p. 72 ---
\opage{72} \setcounter{footnote}{0}%
often shrank back; one cannot wonder, then, that they now sought, as far as
possible, to raise a dam against the rushing torrent. But the best means of
giving the teaching the necessary firmness was that public confession of
faith. For people soon became conscious of the distance that, by the nature
of the case, must always obtain between the objective truth of Scripture and
the subjective religious consciousness, and as a consequence the general
endeavor aimed not so much at proving the unconditional validity of
Scripture --- for on that everyone already agreed ---, as at interpreting it
rightly; but the only infallible criterion of a true interpretation of
Scripture is that consciousness be able to find itself again in it, and just
such a criterion was furnished by the symbols of Protestantism.

Thereby a bond was also laid upon the opposition that had arisen in the
individual congregations between the various subjectivities. For since
subjective faith, which Protestantism cannot combat with external weapons,
had been led by the inner power of spiritual kinship to give its assent to
the public confession and to acknowledge the teaching set forth in it in its
entirety, it followed that individuals no longer contended with one another
about their subjective faith, but that the individual subjectivity was judged
by the universal objectivity\footnote{\textit{Omnino Philippus Melanchton
privatus hanc insignem mutationem (de confessione variata dicitur) fecit,
inconsultis, quorum communis causa agebatur. Quare statim ac ista innotuit,
improbavit illam elector Saxoniæ, Johannes Fredericus, misitque ad
Melanchtonem cancellarium Pontanum, qui eum objurgaret: \ldots\ Walch:
Introductio in libr. eccl. Luth. symb. p. 188.}

How often he had to defend his teaching against objections is shown by the
following words of his own: \textit{Non gigno novas opiniones, nec aliud
majus scelus esse in Ecclesia Dei sentio, quam ludere fingendic novis
opinionibus, et discedere a Prophetica et Apostolica scriptura et consensu
vero Ecclesiæ Dei. \emph{Sequor autem et amplector doctrinam Ecclesiæ
Vitebergensis et conjunctarum, qui sine ulla dubitatione consensus est
Ecclesiæ catholicæ.} Melanch. Opera p.~I. p. 148 Locor. Theol. præfatio.}}.
But thereby there also arose, little
% --- p. 73 ---
\opage{73} \setcounter{footnote}{0}%
by little, a certain external limitation; the bold negation of forms and
phenomena ceased, and since a violation of the validity of the confession of
faith was now also regarded as a violation of truth itself, so that it was
held to be a duty to exclude from the Church, or to imprison, those who
twisted truth into a lie, the State too was at once at hand with its
weapons\footnote{Compare with this --- not to speak of the Calvinists'
church discipline --- the events in the crypto-Calvinist controversies.
Walch. Historische und theologische Einleitung in die
Religions=Streitigkeiten. \textit{1} Th. \textit{2} Cap.}. --- And yet the
need to institute investigations and to raise objections had a far deeper
ground than empty hair-splitting or quarrelsomeness. It lay, to be sure, in
the nature of Protestant faith that at the beginning of the Reformation it
addressed itself almost exclusively to the immediate religious
consciousness, so that it not only separated the dogmas from purely
historical objects\footnote{Denn wo ich je deren eines mangeln sollte, derer
Wercke oder der Predigt Christi; \emph{so wollte ich lieber derer Wercke, denn
seiner Predigt mangeln.} Denn die Wercke hülffen mir nichts; aber \emph{seine
Worte, die geben das Leben,} wie er selbst sagt Joh.~\textit{5, 51}. Vorrede
auf das N.~T. \textit{XII. 55}.}, but even began, within the dogmas
themselves, to separate the moral and practical moments from the objective
and mysterious ones\footnote{\textit{Mysteria divinitatis rectius
adoraverimus, qvam vestigaverimus \ldots\ Proinde non est, cur multum operæ
ponamus in locis illis supremis, de Deo, de unitate, de trinitate Dei, de
mysterio creationis, de modo incarnationis \ldots\ Reliqvos vero locos,
peccati vim, legem, gratiam, qvi ignorarit, non video, qvomodo Christianum
vocem; nam ex his proprie Christus cognoscitur, siqvidem hoc est Christum
cognoscere beneficia ejus cognoscere. Melancht. Loci com. præfatio ed.
principis (ad ann. 1521 ref.)}

Hinfürder lehret er (Carlstad) uns, was Frau Hulda, die natürliche Vernunft,
zu diesen Sachen sagt: gerade als wüsten wir nicht, daß die Vernunft des
Teufels Hure ist und nichts kann denn lästern und schänden alles, was Gott
redt und thut. Wieder die himmlischen Prophet. Schrifft. \textit{XIX. p.
199} cf.\ Von Frau Hulda der klugen Vernunft Doctor Carlstads. Schrifft.
\textit{XIX. p. 208}.}; but if this faith (ὑπόστασις) is
% --- p. 74 ---
\opage{74} \setcounter{footnote}{0}%
to be able to find perfect rest in its substance, it must also have the
fullest confidence in the firmness of the foundation on which it rests; for
if it has its ground only in itself, then it cancels itself. Therefore faith
cannot meet contradiction without passing over to testing the objectivity of
its object as such, that is, to thinking\footnote{Therefore Luther too cannot
refrain, although he combats reason, from defending his teaching with
objective grounds and analogies drawn from metaphysical relations; therefore
he expressly takes refuge in thinking, so as not to confuse
``\textit{usus passionis}'' with ``\textit{factum passionis}'' or
``\textit{quod}'' with ``\textit{qualiter}'': ``Gilt denken und ist genug, so
will ich auch denken, besser denn sie \ldots'' cf.\ Bekenntniß vom Abendmahl
Christ. \textit{XIX. 440}. Melanchton too shows this transition from
subjectivity to objectivity, in that, e.g., in the first edition of his
\textit{loci} (\textit{1521}) he passed over the mystery of the Trinity and
that of the Incarnation, among others, but in the later ones not only added
them but even sought to explain them. Cf.\ \textit{Edit. ann. 1548 absoluta
(In Opp. Viteb. 1580 I. pg 148)}.

\textit{Filius dicitur Imago et} λόγος. \textit{Est igitur imago cogitatione
Patris genita, qvod ut aliqvo modo considerari possit, a nostra mente exempla
capiamus \ldots\ Mens humana cogitando mox pingit imaginem rei cogitatæ, sed
nos non transfudimus nostram essentiam in illas imagines \ldots\ At Pater
æternus sese intuens gignit cogitationem sui, quæ est imago ipsius non
evanescens, sed subsistens communicata ipsi essentia \ldots\ Dicitur} λόγος,
\textit{qvia cogitatione generatur. Dicitur Imago, qvia cogitatio est imago
rei cogitatæ \ldots\ De Filio pg. 152.}}; by what path this transition takes
place is another question.

Absolute faith surrenders the whole human being to God in unconditional
submission; it not only recognizes, in a practical
% --- p. 75 ---
\opage{75} \setcounter{footnote}{0}%
respect, the great corruption of nature and the impotence of good works, but
also, in a theoretical respect, calls darkened reason back to obedience
under faith\footnote{\textit{Hominis intellectus et ratio in rebus
spiritualibus prorsus sunt coeca, nihilque propriis viribus intelligere
possunt. Form. Concord. Cfr. Hase: Hutterus red. 3} Aufl. Pg. \textit{65
etc.}}. Nor was the severity with which that age sought to suppress reason
mere subjective arbitrariness; for by reason was understood the everyday
understanding, which abstracts all its concepts in an external way (cf.\
§~\textit{5}), and must thus necessarily enter into sharp opposition to
faith, which considers everything from within. Hence the theologians of that
time laid so much weight on the proposition that a true and saving cognition
of God could be drawn only from a special revelation\footnote{\textit{Datur
qvidem revelatio naturalis, sed prorsus mutilata: Ut post splendidæ domus
conflagrationem favillæ et adustæ particulæ, aut post collapsum gravissimo
casu ædificium rudera nonnulla remanent, ita quoque vix tenues qvædam
reliqviæ ac} ἐρείπια \textit{divinæ imaginis et parvæ scintillæ primævæ
lucis, qvæ in mente hominis ante lapsum fulgebat, in natura jacente et misere
prostrata supersunt. Qvenst. Theolog. Pars I p. 7.} Θέσις \textit{XIII not.
III.}}. --- But how can darkened consciousness be receptive to a true
revelation? The greater the separation and opposition between the holy God
and corrupted nature, the more wondrous must be the activity by which the
dividing wall between the two is broken down. If, then, a revelation is to
become possible, nature must become other, he who is above nature must show
himself visibly in it; there must appear an immediate identity of the
mutually opposed moments, i.e.\ a miracle must take place. But if God is to
make himself known through miracles, he can be cognized in his phenomena only
in the following way. One hears God speak, one
% --- p. 76 ---
\opage{76} \setcounter{footnote}{0}%
sees heavenly figures reveal themselves; everything that is heard with the
ears and seen with the eyes is fact for the pious believer. In everything
that shines forth in the phenomena, the mysteries of divine truth are
revealed. --- But a contradiction lies hidden in the concept of a revealed
mystery; for insofar as a revelation takes place, to that extent the mystery
descends into the sphere of phenomena; through this descent a natural moment
shows itself in the mystery, and in this respect it cannot then be
distinguished from other phenomena; but insofar as it is a revelation of the
mystery, it also contains a supernatural moment and is infinitely exalted
above all other phenomena. One must therefore know how to distinguish the
phenomenon of the mystery from all other phenomena. But if one cannot
acknowledge the divinity of the phenomenon without knowing the mystery, and
cannot know the mystery without acknowledging the divinity of the phenomenon,
then surely a revelation is required in order to cognize the revelation, then
surely the cognition of the mysteries and miracles is itself a miracle.

Every special revelation of God can thus be referred to the two principal
miracles, inspiration and the illumination mediated from inspiration through
the Church. --- As the former has given us a perfect presentation of divine
things in infallible Scripture\footnote{\textit{\ldots\ revelationes certis
testimoniis in verbo suo proposuit, in qvo ceu foetus in alvo materna
nutrimentum attrahens ex umbilico et cotyledonibus nos qvoque sedemus
inclusi, et attrahimus agnitionem Dei et vitam ex verbo Dei, ut cum
invocemus, sicut se patefecit. Melancht. l.~c. cap.~I pg. 151 \ldots\ in
Scriptura Canonica nullum est mendacium, nulla falsitas, nullus vel minimus
error sive in rebus, sive in verbis. Qvenst. l.~l. p.~I. p. 77.}}, so the
latter enlightens consciousness
% --- p. 77 ---
\opage{77} \setcounter{footnote}{0}%
to see clearly and rightly\footnote{\textit{\ldots\ adeo ignorans et perversa
est ratio illa, ut etiamsi ingeniosissimi et doctissimi homines in hoc mundo
Evang. de Filio Dei et promissiones divinas de æterna salute audiant, tamen
ea propriis viribus percipere et vera esse statuere neqveant. F.~C. cfr.
Hntt. red.} Pg. \textit{67.}

\textit{Fides divina, qva qvi infallibiliter credit, hoc verbum, qvod
Scriptura proponit, esse vere Dei verbum, immediate oritur a vi et potestate
S. Scripturæ intrinseca, qva Scriptura ipsa ad generandam fidem sibi debitam
efficax est. Qvenst. p.~I. pg. 98.}}. It is then easy to understand what is
meant by the demand that we take reason captive. For all the phenomena of
revelation must be considered from a twofold side. Considered from the
objective side, God reveals himself in them: he truly thinks as he speaks; he
is in himself as he wills to be represented by human beings; he is triune,
for he has revealed himself as triune; he is
\emph{truthful}\footnote{\textit{Causa certitudinis est revelatio Dei, qvi est
verax. Melanch. l.~l. pg. 148.}}. But although God is truth, and his nature
cannot contradict itself, human reason nevertheless grasps only those
properties in God that resemble its own and belong to our world; since,
then, God has a nature of his own, reason itself recognizes that with respect
to this nature it must deny its own validity\footnote{\textit{Qvando
objiciunt adversarii (Socin, Armin.) religionem multa habere supra, nihil
vero contra rationem, respondeo: 1) Articuli fidei in se non sunt contra
rationem, sed solum supra rationem; per \emph{accidens} vero fit, ut sint
etiam contra rationem, qvando ratio judicium sibi de illis sumit ex suis
principiis, nec seqvitur lucem verbi, sed eosdem negat et impugnat. Qvenst.
p. 45 XIX Distinct.}

\textit{Cognitio Dei naturalis non extendit sese ad essentiam divinam,
qvatenus illa relativa seu in Tribus Personis consideratur \ldots} Ibid.\
\textit{p. 265} Θέσις.

\textit{Congruentiæ naturales et analogia creatarum rerum cum hoc mysterio
non divinam fidem, sed opinionem tantum humanam generant \ldots} Ibid.\
Ἔκθεσις \textit{p. 266 IV. Obst}}. But on the other
% --- p. 78 ---
\opage{78} \setcounter{footnote}{0}%
side we are surely to know the phenomena of revelation and apply them to our
own life; if, now, human reason had to declare them to contradict one
another, so that we would always have to be in doubt about their truth, then
in actuality no revelation would fall to our share at all. It will therefore
be expedient, in the realm of the divine phenomena, to employ reason in a
formal respect\footnote{\textit{Aliud est principia et axiomata philosophica
in Theologia adhibere illustrationis, declarationis et secundariæ probationis
gratiâ, ubi res e Scriptura definita est; et aliud, adhibere eadem decisionis
et demonstrationis causa, vel principia philosophica et ratiocinationes ex
illis factas pro principio in Theologia agnoscere, aut ex illis de rebus
fidei judicium ferre.} Ibid.\ \textit{p. 42 II Obs.} cf.\ \textit{III. IV. V
VI. Distinct.}}. --- What, now, is the supreme law that reason, reinstated in
its lawful dominion, must lay down? Since the revealed truths cannot be
conceived in the sense that one could demonstrate the dialectical transition
between them, and the ordering activity has only the task of showing that
they can all exist alongside one another, because the certainty of their
existence derives solely from their having been revealed\footnote{\textit{Etsi
autem omnes Angelicæ et humanæ mentes obstupescunt admiratione hujus arcani,
qvod Deus Filium genuit \ldots\ tamen sic sentire necesse est, qvia, ut
toties jam dictum est, de Deo sentiendum est, sicut se patefecit \ldots\
qvanqvam hæc arcana non penitus perspicimus, tamen in hac vita Deus inchoari
vult hanc agnitionem, et invocationem nostram a falsa discerni.} Ibid.\
\textit{p. 15.}}, one must above all hold fast to the well-known proposition
that what is actual must also be possible (\textit{ab esse ad posse valere
conseqventiam}). This
% --- p. 79 ---
\opage{79} \setcounter{footnote}{0}%
law entails that logic itself now renounces the use of its own rules, now is
summoned to use them. Thus in the dogma of the Trinity logic falls silent;
for one God has revealed himself in three persons, therefore One can be in
Three; on the other hand, one ought to listen to reason when, according to
its own rules, it sets forth the doctrine of the omnipresence of Christ's
body, although this is not expressly set forth in Scripture, for Christ has
surely been revealed as the one who at once sits at God's right hand and is
really present in the Lord's Supper\footnote{\textit{Gerhard: Qværitur, sitne
corpus Christi vere et substantialiter in coena præsens? Affirmant Eccles.
nostræ, qvia Christus dicit: \emph{Hoc est} etc. Negant adversarii ex hoc
principio, qvia verum et naturale corpus non potest simul et semel esse in
pluribus locis. Addunt, rationem renatam non posse aliter de corpore
statuere, testimonium ergo ejus audiendum esse. At inqvam, ratio renata,
qvatenus talis de art. fidei statuit ex Dei verbo, et limites ejus non
egreditur. Cfr. Hutt. red. p. 71.}}. Hence there also arises the division of
the \textit{articuli fundamentales} of faith into \textit{primarii} and
\textit{secundarii}, and of its \textit{articuli conservativi} into
\textit{antecedentes} and \textit{consequentes}\footnote{Divisions which show
that our fathers did not neglect to take account of the principles of
reason: ``\textit{Articuli \emph{conseqventes} sunt, qvi explanant fidei
salvificæ conseqventia, qvibus positis fides confirmatur, \emph{non positis}
evanescit.}'' \textit{Art. \emph{secundarii}, qvi} ``\textit{licet ignorari
possint illæso fidei salutisque fundamento, pertinaciter tamen negari salvo
illo non possunt, nedum impugnari.}'' \textit{Hutt. rediv. §~17 p. 24.
25.}}. --- Since by this mode of treatment one could arrive at no insight
into the essence of the divine nature, but only consolidated an external
acquaintance with its mysteries, and these were not drawn forth from the
phenomena in which they had appeared but rather remained enclosed in them, a
separation between the idea and its external form was naturally impossible,
and one only tied
% --- p. 80 ---
\opage{80}\setcounter{footnote}{0}%
the external bond between dogmatics and sacred history all the tighter.
Human traditions had been rejected and the Catholics' τερατολογία given up;
one was therefore all the more zealous in guarding the foundation of
supernatural facts that had been left untouched by doubt, and could not
concede, even in the most insignificant points of sacred history, the
possibility of any error in inspired Scripture\footnote{\textit{\ldots\ omnia
et singula sunt verissima, qvæcunqve in illâ traduntur: sive dogmatica illa
sint sive moralia, sive Historica, Chronologica, Topographica, Onomastica,
nullaque ignorantia, incogitantia aut oblivio, nullus memoriæ lapsus
Spiritus Sanctus amanuensibus in consignandis sacris litteris tribui potest
aut debet. Qvenst. p.~I pg. 77} Θέσις.

And it was not only dogmaticians toward the end of the \textit{17}th century
who put forward such opinions; historians too, in the latter part of the
\textit{19}th century, have in like manner distinguished between divine and
human sources:

Ihre (i.e.\ Holy Scripture's) Verfasser haben nicht allein nie geirret,
sondern auch nicht einmal irren können, und ihre Erzählungen sind daher nicht
allein so gewiß, als jede andere historische Wahrheit, sondern auch
untrüglich gewiß. Ch.~W.~F. Walch, krit. Nachricht v.~d. Qvellen der
Kirchengeschichte. Leipz. \textit{1770} p, \textit{23}. --- They therefore
not only, following Luther's example, made manifold fruitful applications of
sacred history to the life of Christians in a moral respect, but even
grounded important dogmas on events narrated in Scripture:

Thus, e.g., E. Brochmand, by comparison and ingenious analysis of a few
passages of Scripture (Gen. \textit{1, 26. 27}; Eph. \textit{4, 24}; Col.
\textit{3, 10}), refutes with little trouble all of Bellarmine's objections
against the doctrine of the perfection and immortality of the first human
being. \textit{Univ. Theol. Syst. p. 124. Quæst. III. IV.}

In his proof that there was in the first human being's cognition a special
agreement with God ``ὁ πρωτοτύπος'', Qvensted says: \textit{\ldots\ Ex
Adami factis qvalia sint 1)} Ονοματοθεσία, \textit{sive nominum conveniens
impositio Gen. 2. 19, qvæ erat non solum \emph{Grammatica}, qvoad
nomenclaturam animalium, sed vel maxime \emph{Logica}, qvoad verissimam
definitionem.} Indeed, the personality of the Holy Spirit was proved in the
same historical manner; so already Melanchton: \textit{\ldots\ impio
sophismati Ecclesia vera opponit testimonia in scripturis tradita, qvorum
primum et maxime illustre est revelatio divinitatis, facta in baptismo
Christi, ubi clare discernuntur tres personæ. Pater inqvit: Hic est Filius
meus dilectus. Est igitur alia persona Patris, alia Filii. Descendit autem
Spiritus Sanctus in specie columbæ. Jam si Spiritus esset agitatio in animis
creata, non appareret peculiari specie corporali.}

\textit{Sic et in Pentecoste apparet Spiritus Sanctus peculiari specie
corporali. Hæ patefactiones non sunt frustra factæ, immo sunt excellentia
beneficia Dei, in qvibus Deus se patefecit Ecclesiæ, et testificatus est
Spiritum Sanctum esse personam.} S.~S. Pg. \textit{137}.

And following his example Qvensted too proves the Trinity ``\textit{ex
singulari illa} θεοφανείᾳ \textit{facta in baptismo Christi \ldots\ ut vere
dictum sit ab antiqvis:} `\textit{abi Ariane ad Jordanem! et videbis
Trinitatem}'\,'' Pg. \textit{324}.}. Since the theologians thus regarded
Scripture's
% --- p. 81 ---
\opage{81} letters as the sole and sufficient foundation for positive doctrine,
while they denied the Enthusiasts (the Schwenkfeldians and Weigelians), with
their spiritual inspirations and angelic revelations, all credence, their whole
endeavor went toward fixing the objective forms of faith. They therefore did
indeed abstain from empty speculations, for they would not, like the
Scholastics, ``set up a centaur-like theology, a \textit{mixtum compositum} of
divine revelations and philosophical inferences of reason''; but they
nevertheless left to reason the important calling of ordering and arranging
the material. For in the great mass of objects one might fear that the one
would seem to contradict and cancel the other, and one therefore had to take
care to demonstrate, as best one could, the possibility of the mutual
subsistence of all the moments.
% --- p. 82 ---
\opage{82} \setcounter{footnote}{0}%
Reason then had to undertake to carry out this laborious work by its
affirmations, negations, distinctions, divisions, and exceptions. And as it
now, in the subtle analysis of the individual moments, stated τὸ ζητούμενον,
set up τὴν θέσιν, added τὴν ἀντίθεσιν, took on τὴν βεβαίωσιν, the work grew
into such a mass of propositions that we are uncertain what we should most
admire, --- the perseverance with which reason, admittedly quite
mechanically, fitted together all the individual parts, or --- the strength
that enabled faith to bear this great building, erected out of the most
opposed elements\footnote{Das dogmatische System des \textit{16}-und
\textit{17} Jahrhunderts kam mir vor wie einer unsrer alten deutschen Münster
mit seinen himmelstrebenden Spitzbogen und wunderlichen sinnvollen
Zierrathen \ldots\ Einen Dom wie unsre Vorfahren kann auch unsre Zeit nicht
wieder bauen \ldots\ es wird einem aber doch ganz besonders wie in einem
Gotteshause darin zu Muthe. Hase: Der theol. Facult. zu Tübingen
(\textit{Hut. rediv,})}. This magnificent temple in the grand style, which the
human spirit, returning from exile, built for itself, was built up and
fortified amid the manifold storms of a turbulent age; always in arms, always
ready for battle, these men raised their building; --- whether it was gold and
silver and precious stones of which they built it, or whether it was wood and
hay and straw, only the dawning day, the testing fire, will show.
\parag{16}{Continuation. Transition to Criticism.}

The absolute agreement between the essence and its phenomena, on which this
building rests, sublates itself.
% --- p. 83 ---
\opage{83} \setcounter{footnote}{0}%
For since the divine nature has revealed itself not in one but in several
phenomena, and is expressed not in one but in many mutually contradictory
ways, there must truly be in every expression of the idea of God something
more or less inadequate\footnote{\textit{Quia intellectus noster finitus
infinitam et simplicissimam Dei essentiam uno conceptu adæquato adæquate
concipere nequit, ideo distinctis et inadæquatis conceptibus essentiam
divinam inadæquate repræsentantibus eandem apprehendit. Quenst.} P.
\textit{284,}}; one must consequently be able to take the individual word in
different senses in its different places, and to distinguish between what is
said in a proper and what in an improper sense\footnote{Here belong all
anthropomorphisms; thus Qvensted distinguishes ``\textit{ea, quæ
imperfectionem aliquam subinferunt et Deo proprie tribui nequeunt, sed
improprie et per} ἀνθρωποπάθειαν \textit{ipsi adscribuntur, ut affectus
etorgana corporea, ut ira, odium, manus, pedes, oculi}'' \textit{etc,} \ldots{}
P. \textit{296}.}. One then reverses the rule and infers from a thing's
impossibility to its non-actuality (\textit{a non posse esse ad non esse valere
consequentiam}), --- but in this way reason recovers its lost freedom, for it
is, after all, reason that must compare the individual expressions. Where two
expressions are so opposed to one another that the one obviously excludes the
other, the result will of course differ according as one takes the one or the
other of them to be certain. But which is to remain standing and which is to
go can be decided as little by the sum of the revealed utterances --- for it
is precisely the individual utterances that are at issue ---, as by reason
itself, for what subjective reasoning at this standpoint is concerned with,
solely and exclusively, is precisely to have a fixed and unshakable point of
departure; but in this way the entire system of objective doctrine vanishes,
and only subjective faith in its full indeterminacy remains.
% --- p. 84 ---
\opage{84} \setcounter{footnote}{0}%
One can then grant that one God has revealed himself in three persons, and
yet at the same time, by the rule cited above, conclude that where three
persons in God are spoken of, this can, on account of God's unity, only be
taken in a figurative sense\footnote{Thus the Socinians' objection:
\textit{Quoniam essentia divina tantum unica est numero, ideo nullatenus
plures in eâ possunt esse personæ. Catech. Racov. Cap. 1, p. 21.}

Indeed, Volkelius \textit{Libr. V. de V. R. cap. IX. fol. 401} goes so far as
to ascribe the doctrine of the Trinity to ``the old liar Satan''; and Joachim
Stegmannus teaches that this mystery conflicts both with \emph{Scripture} and
with \emph{reason}. Cf.\ Quenst. P. \textit{345}.}. --- Since one sought to
fix doctrine in this way, it was a matter of course that subjective faith also
had to hold fast to the individual propositions of doctrine, just as it sought
to defend the truth revealed in Scripture in general; for if the truth of one
article of faith could be called into doubt, it would of course also be easy
in the same way to weaken faith in each of the others. Therefore in the
religious colloquies one could not give way on any point whatsoever; disputes
were always conducted \textit{de omni et de nullo}. Hence the unhappy outcome
of the irenic colloquies\footnote{Thus, e.g., complaints are made in vehement
terms about the Colloquy at Thorn:\\
\textit{Quæ sunt gesta tribus Thoruna mensibus urbe,}\\
\textit{Hæc potuere tribus cuncta diebus agi.}\\
and\\
\textit{Quid synodus? nodus. Patrum chorus integer? æger.}\\
\textit{Conventus? ventus. Gloria? stramen. Amen.}\\
Cf.\ \textit{Dr.} K.~H. Hagenbach: Vorles. üb. W. Geschicht. d. Reform.
\textit{4} Theil, P. \textit{148---149}.

On the syncretistic controversies, which furnish rich contributions to the
characterization of that age, cf.\ J.~G. Walch. R. Str. \textit{1} Theil,
§~\textit{5}, P. \textit{216} ff.}; hence the fruitlessness of the
Syncretists' efforts to keep to the middle way; hence the vehemence, the
partisanship, the fanaticism
% --- p. 85 ---
\opage{85} and the furious rancor of the theologians of that time. It is only
by developing to the extreme that such systems are completed and sublated;
therefore one must always observe them in the course of their own
development.

The fundamental principle of faith is thinking; therefore faith demands an
immediate identity of the subjective and the objective moment; but in this
there also lies hidden an immediate opposition. The former emerges first, the
latter later. When the subject, in immediate agreement, gives itself over to
the object, it does not occur to the believer to institute an external
comparison between himself and the object of his faith, --- for the thought of
having to institute a careful and petty examination of the properties of both,
or that these should need to be adjusted to one another in order to accord
with one another, is unnatural to the first, immediate faith ---; but, gripped
by the greatness of the object in which he believes, his whole endeavor goes
toward determining and fixing the forms of objective faith. From this arose
precisely the Protestant faith's desire to set up systems. --- But when faith
has thus immersed itself in pure objectivity, the relation changes. For the
absolute validity of the objective system shows itself in this, that the
objects hold not because they are conceived, but because they are factually
given. The system then appears as something foreign that is to be obeyed
because it demands obedience, i.e., objectivity swallows up subjectivity. But
thereby the system also tears itself down; for though its subjective origin
has been forgotten, it is in actuality faith that created it; therefore the
subjective consciousness feels itself hampered by foreign, external bonds when
the system would assert itself over spirits in an external manner; it begins
to reflect, and now the oppositions come forth; but then the objective system,
too, will dissolve and perish,
% --- p. 86 ---
\opage{86} %
because it is only on the immediate sympathy of subjectivity with it, if one
may say so, that its whole power rests.

Since, therefore, the individual propositions of faith had been immediately
set up in this theology, faith itself passed over into being a sum of external
propositions, a petrified formalism. One collected with the greatest avidity
\textit{dicta probantia} and \textit{loci communes}, not in order to justify
the statements of Scripture from within, in the name of truth, but in order, in
the name of Scripture, to come in an external manner into tangible possession
of the truth: one worshiped the individual words and signs in Scripture.
Therefore spirit withdrew from this objectivity of spirit; therefore the true
light of the Bible was darkened under this \emph{bibliolatry}. --- A negation
of the positive forms followed upon this formalism; the religious
consciousness began to seek the truth in its own interior, and the powers of
our Scholastics were exhausted in struggles against the Pietists.

Thus the second moment in faith now emerged. The subject is indeed to give
itself over to the object with its whole heart; but --- once bitten, twice
shy; therefore the believer must first turn his eye inward into the depths of
his own consciousness and examine his own heart. --- Protestantism thus
contains not only its own negation but in equal measure its own restoration;
as the real moment had at the beginning been the preponderant one, precisely
because it immediately gripped conviction, so now little by little the formal
moment gained greater significance, so that philosophy, at the very same time
as Protestantism, after the end of the Thirty Years' War, obtained full
recognition and legal existence in the state, could turn away from theology
and teach \emph{that one should doubt everything}. In the name of concrete
faith Luther asserted freedom of conscience in the \textit{16}th century; in
the name of freedom, in the \textit{18}th and \textit{19}th centuries, his
% --- p. 87 ---
\opage{87} %
\setcounter{footnote}{0}%
true followers defended freedom and thinking, and so make possible a rebirth of
the objective, concrete doctrine of faith\footnote{\ldots\ even in the case
that the Reformers had meant not only to depict the spirit and tendency of
Protestantism, and along with it their own and their contemporaries' religious
conviction, but also to set up in this conviction an unalterable norm of faith
for the following generations, one would have to lament the lot of humanity,
which is weakness even in its greatest strength; but no obligation for the
Protestant Church could be derived from this. Clausen: Cath. og Protest.
P.~\textit{331}.\par
Der Protestantismus hat, was vom Anfange in ihm lag, offenbart, die
Entwicklung der religiøsen Selbständigkeit, die durch freie Liebe und Einsicht
am kirchlichen Vaterlande festhält so daß die Treue an Luthers Geiste zum
Abfalle von seinen Buchstaben \ldots\ geworden ist. Hase: Kirchengesch.
(\textit{1834}) §~\textit{545}.}.

\parag{17}{Subjective Rationalism. The Profanation of Sacred History.}

After the path on which one had previously sought to set up doctrine had been
abandoned, the objective forms of faith, too, could not, at least for a time,
escape being disdained, since these after all made up the proper basic
material of the obsolete system; eager for occasion to undertake
investigations and to exercise its thinking, the human spirit leaves untouched
only that part of divine truth which immediately grips the religious heart,
and then turns to other regions that can offer a new and richer material for
it to work upon. --- The objects of faith recede into the background, while
the objects of this world lie nearest to the observer; the sensuous
perceptions and impressions, through
% --- p. 88 ---
\opage{88} %
which access to the wide field of experience stands open, are the first to
meet the gaze that seeks an immediate identity of subject and object. Yet the
observer will not long be satisfied with roaming about in external actuality,
foreign to consciousness; he will soon recognize that it is only empirical
things that are here the object of his thought. Thinking is a negation of the
thing; therefore he turns his gaze away from sensuous nature and seeks a
rational world order, in order to find himself again in this objectivity. Now
although empiricism and abstract rationalism are wholly opposed to one
another, since the one occupies itself with individual, concrete objects, the
other, by contrast, solely with universal and abstract ones, they nonetheless
readily join hands, for the elements of both attract consciousness, and
abstract reason needs empirical knowledge in order to give its hollow concepts
content out of the latter's fullness. As for the doctrine of faith, both stand
in a wavering relation to it; on the one hand they are akin to faith: the
content of faith is, after all, nothing other than the sum of spiritual
realities, the same holds also of the objects we meet in our conscience, and
finally both demand that the thinking spirit be able to find itself again in
objectivity, but this the nature of faith also demands; --- on the other hand
they stand in direct opposition to faith. For faith derives all its elements
from the holy God, exalted above our nature; these, by contrast, proceed from
the subject and created nature and never overstep the profane bounds of this
world of ours.

On account of this dubious kinship they enter into a double relation to the
Church. For as soon as they betray the profane nature of their principles,
they are turned away in the name of offended religiosity, and if they seek,
with banner raised and drums beating, to storm the fortress of Christ, they
fall before the sword of the mightier
% --- p. 89 ---
\opage{89} %
\setcounter{footnote}{0}%
Spirit; if, on the other hand, they succeed in sneaking into the fortress in
disguise, they begin a secret undermining. Yet the cunning of these enemies is
not to be understood as though they had by some means or other bribed
conscienceless (κεκαυτερισμένους τὴν ἰδίαν συνείδησιν) theologians to cast
suspicion on the truth of religion; on the contrary, theologians of this sort
generally act from the fullest conviction and with a bold countenance. The
struggles that arise from this are therefore also not waged directly over who
is Christ's friend or enemy, but over what belongs to the essence of religion
and what belongs only to its phenomena. For those who remain with the older
method of teaching urge the supernatural moment in the sacred phenomena,
declaring that these constitute the objective forms of faith and therefore
cannot be classed with the ordinary events that are subject to our judgment,
any more than they can be altered or volatilized without the greatest harm to
religion itself\footnote{The controversies that took place in the Reformed
Church on the occasion of the critical treatment of Holy Scripture furnish, so
to speak, a prototype for all similar debates. Thus Buxdorf declared Ludv.
Capellus a \textit{\emph{novator, propheta, videns}, qui somniis
indulgeret}. Capellus, by contrast, replied \textit{(advers. Censorem Buxdorfium
\ldots\ brevis oeconomia): Satis, mihi fuisset id ostendisse adversus Censorem
meum \ldots\ nisi illi ipsi et aliis multis altius animo insedisset et infixa
esset hæc opinio: Actum esse de Dei verbo ejusque apud nos auctoritate, atque
ex consequenti de hominum salute, (quæ non aliunde quam ex Dei verbo pendet),
si apiculi isti non statuantur esse a Mose et Prophetis, libris suis, Deo eos
docente, additi et adseripti.} And rightly does Georgius Pasor put the
following words in the mouth of Sebastianus Pfoechenius, who is famous from the
purist controversies:\par
\hspace*{2em}\textit{Diligo Philologum, qui eviscerat abdita verba,}\par
\hspace*{2em}\textit{Jure novum Chrtsti sed mage \emph{fedus} amo,}\par
\hspace*{2em}\textit{Ergo \emph{novo} quisquis de \emph{federe} tenuia dicit,}\par
\hspace*{2em}\textit{En pupillam oculi vulnerat ille mei.}\par
\centerline{\rule{0.16\linewidth}{0.4pt}}\par
\hspace*{2em}\textit{Falx critica in quosvis per me grassetur. At unis}\par
\hspace*{2em}\textit{A Bibliis fas est abstinuisse manus.}}.
% --- p. 90 ---
\opage{90} \setcounter{footnote}{0}%
The fundamental error in the investigation thus lies in this, that one does
not infer from the truthfulness of the self-revealing God to the truth of the
phenomena and forms through which the revelation takes place, but believes that
through an investigation of the phenomena one can determine which moments are
essential to the revelation. They therefore, trusting in the infallibility of
human investigation, urge the natural side of the phenomena, and declare it a
\textit{petitio principii} if anyone, in the name of faith, holds fast to the
individual forms and phenomena, since only by an investigation of them can one
come to determine the objects of faith. And it is on this conviction, that
truth cannot contradict truth, that empirical truth cannot sublate the truth of
religion, that this school rests its authority\footnote{\textit{Sunt tamen et
pii et docti viri, qui cum vix habeant, quod tam manifestæ veritati opponant,
pergunt nihilominus periculosas inde metuere consequentias. Quasi vero a
veritate ullum jure metuendum sit veritate periculum! Verum certe vero semper
consonat fertque opem, nunquam repugnat aut officit \ldots\ Ergo candido et
christiano pectore, et forti constantique animo dignius esse existimavi, fateri
in hac causa, quod res est, quam cum nonnullis id oppugnare, quod sine rubore
ab eo, cui ingenua est frons, negari non potest. Ludv. Capellus. Arcan. punct.
Præfatio ad lect.}\par
That thinking, like cognition itself, is universally valid, since it has faith
as such within itself, Hegel has fully proved in his Phænomenologie des Geistes
(Werke P.~\textit{410}): Von dieser Seite, daß beide wesentlich dasselbe sind
und die Beziehung der reinen Einsicht durch und in demselben Elemente geschiet,
ist ihre Mittheilung eine unmittelbare, und ihr Geben und Empfangen ein
ungestörtes Ineinanderfließen \ldots\ Die Mittheilung der reinen Einsicht ist
deswegen einer ruhigen Ausdehnung oder dem Verbreiten wie eines Duftes in der
wiederstandslosen Atmosphäre zu vergleichen \ldots\ ein unsichtbarer und
unbemerkter Geist, durchschleicht sie die edlen Theile durch und durch, und hat
sich bald aller Eingeweide und Glieder der bewußtlosen Götzen gründlich
bemächtigt, und ``an einem schönen Morgen giebt sie mit dem Ellbogen dem
Kammeraden einen Schubb und Bautz! Baradautz! der Götze liegt am Boden.''}.

% --- p. 91 ---
\opage{91} %
Under such a relation there could naturally no longer be any question of fixing
the positive forms of faith; the phenomena of revelation, which stemmed from
grey antiquity, and the forms which had already been fixed in the course of
time, became merely external objects for the investigator, and the results of
his investigations could therefore consist only in finding out which events
could have taken place, which had actually taken place, and which, finally,
were to be regarded as products of a superstitious fantasy. Faith was then no
longer the constitutive principle in theology, and although that which is
taken as the standard can by its nature admit or recognize nothing other than
what lies within its own bounds, nature was nevertheless made the standard for,
and the judge over, supernatural things. One cannot therefore wonder that the
positive elements of religion became fewer and fewer and at last vanished
entirely. While empiricism thus proved that the Hebrew vowels had been
invented neither by Adam nor by Moses, nay, not even by Ezra, that the Greek
of the New Testament was not classically pure, that the reports of Holy
Scripture were both self-contradictory and at variance with the testimony of
credible profane writers, and that many of the individual books were by no
means authentic%
% --- p. 92 ---
\opage{92}%
\setcounter{footnote}{0}%
\footnote{Cf.\ Selmers Abhandlung von freier Untersuchung des Kanon.},
subjective rationalism, on the other hand, (which shows precisely the mutual
relation between empirical knowledge and thinking), demonstrated the
unreasonableness of assuming a supernatural inspiration\footnote{Cf.\
Eckermann: Beweis, daß der biblische Begriff der Offenbarung kein andrer, als
Offenbarungsbegriff der gesunden Vernunft sey.}, raised doubts about the
possibility of miracles\footnote{Die Verächter der Bibel nehmen gerade die
gefährlichsten Waffen, womit sie das Ansehen, die Ehrwürdigkeit und
Glaubwürdigkeit der Bibel bestreiten, aus der Behauptung der Vertheydiger
derselben her, daß ihre Verfasser sich einer unmittelbaren und übernatürlichen
Offenbarung und Eingebung rühmen. Die Verfasser der Bibel werden deswegen als
eine Reihe von Schwärmern angeklagt und herabgewürdigt. The same: Handbuch der
christ. Glaubenslehre \textit{1} Theil P.~\textit{410}.\par
Aber sagt man, wenn ein Mann von erprobter Frömmigkeit, Rechtschaffenheit und
Wahrheitsliebe \ldots\ versicherte, daß Gott dies unmittelbar und übernatürlich
bewirke \ldots\ sollten wir denn das einem solchen Manne nicht glauben? Die
Vernunft antwortet: Nein! Es ist unmöglich, daß ein Mensch unmittelbaren
Wirkungen Gottes erkenne! Auch der redlichste Mensch irrt, wenn er etwa
dergleichen zu erkennen glaubt \ldots\ Ibid. P.~\textit{456}.}, and declared
that regeneration was to be thought of solely and exclusively as a subjective
change in man's free will\footnote{Thus Paulus (Comment. over d. N. Test.)
paraphrases John \textit{3, 6} τὸ γεγεννεμένον ἐκ του πνεύματος: der Mensch,
welcher durch seine Geisteskraft, Einsicht und Wollen des Guten sich selbst
zeugt \ldots\ which he elucidates further by referring to Seneca, \textit{de
brev. vitæ}: ``\textit{Solemus dicere, non fuisse in nostra potestate, quos
sortiremus parentes \ldots\ nobis vero ad arbitrium (!) uasci licet.}''}.

The higher the sun of reason climbed in the heaven of enlightenment, the farther
the light of experience spread its rays, the more the radiance of the heavenly
figures paled, the more lights
% --- p. 93 ---
\opage{93} %
\setcounter{footnote}{0}%
were extinguished in the supernatural heaven. The acuity of reason come of age
discovered that the voice of God and the appearances of angels, of which
Scripture so often reports, were indeed something that could appeal to minors
(οἱ νήπιοι), but in actuality were for the most part only consequences of
natural events, thunder and lightning or the like, and products of an excited
fantasy; such revelations, then, consist not in the divine mysteries revealing
themselves in transfigured nature, but in quite natural events being wrapped in
the misty veil of human fantasy; one then need only draw the veil away in order
to grasp the naked truth\footnote{Excellent examples of such a revelation are
found in G. L. Bauer (Mythologie des alten und neuen Test. Leipzig
\textit{1802}), and so ingenious are these that we cannot refrain from sharing
some of them with the kind reader. Thus on Moses, P.~\textit{271}. \textit{1}
Band: Moses stellte sich vor, ein brennender Dornbusch sey Jehova. \emph{Kann
das Geschichte seyn}? Kann der brennende Busch würklich Gott, oder nur Gott in
demselben gewesen seyn? Ist Gott ein Feuer? --- P.~\textit{272}: Ein
Blitzstrahl fuhr in einen Dornbusch, als Moses weidete am Horeb mit seiner
Heerde: Es war aber ein kalter Strahl, der nicht zündete \ldots\ Nach Hezel
wars electrisches Feuer.}. God, then, did not with an audible voice from the
thorn bush choose Moses as his messenger in spite of his refusal, no! Moses
imagined he heard Jehovah in the thorn bush set alight by the lightning,
because his head was full of great plans; and neither is it God who dictated
the Law through Moses, no, it was Moses who did so in God's name, since it is
beneath God's dignity to speak in articulated tones\footnote{Ibid.
P.~\textit{225} \textit{2} Band: Und daß Gott selbst vom Himmel geredet habe,
ist auch schwer zu glauben, wenn reden heissen soll articulirte Töne, Worte der
damaligen palästinischen Landessprache aussprechen; and P.~\textit{296}.
\textit{1} Band: Hätte Gott selbst gesprochen, sollten denn die Israeliten so
gar dumm, unempfindlich und hartnäckig gewesen seyn, daß sie gleich darauf sein
Verbot des Bilderdienstes vergassen? Und für die zehn Gebote, welche zwar gut
und meistens Natur= und Vernunftgebote sind, zu deren Erfindung aber kein
göttlicher Verstand gehört, sollte eine Gesetzgebung vom Himmel nöthig sein?}.
% --- p. 94 ---
\opage{94} \setcounter{footnote}{0}%
Not only the angelic appearances of the Old Testament but also those of the
New are to be regarded as human illusions. Sound reason teaches us, after all,
that since no angels appear to us, neither can Gabriel have announced to Mary
that she was to bear Christ; and since it is certain that nothing can come from
nothing, and since what is natural, moreover, can be no shame to anyone, we do
best to admit that Jesus was begotten in the natural way by Joseph and
Mary\footnote{Soll die Erscheinung nur von einer Vision oder innern Intuition
oder den eigenen Gedanken der Maria erklärt werden, so fällt zwar der Einwurf
von der persönlichen Erscheinung eines Engels weg, aber dafür ist
unaufgeklärt, wovon denn die Maria ist schwänger geworden. Durch die bloße
Einbildung hat sie doch wohl nicht schwanger werden können \ldots\ Weit
natürlicher ist's, Jesum für einen Sohn des Josephs und der Maria, unsere Sage
aber für eine später entstandenes Raisonnement über die Entstehung des Messias
zu halten Ibid. P.~\textit{192}. \textit{1} Band.\par
\textit{Neque vero fieri potest, ut religionis christianæ auctori major dignitas
aut excellentia concilietur, si vitæ ejus partes et casus non secundum
naturalem rerum humanarum viam a Deo constitntam, sed supranaturali sive
miraculosa quadam ratione evenisse dicantur. Nam si hoc esset, consequeretur,
naturales rerum in orbe gerendarum leges, a divina sapientia et omnipotentia
constitutas, quasi minorem quandam dignitatem habere, hominumque præstantiam ex
iis potius, quæ ipsis accidunt, quam ex ipsorum consiliis et factis debere
existimari. Wegscheider: Institutiones Theol. Halæ 1833. p. 442 not. a.}}. But
how can one who is born in the natural way employ supernatural powers in his
service! --- John the Baptist likewise imagined that he saw the Holy Spirit,
while in actuality he saw nothing but an ordinary dove that happened to fly
% --- p. 95 ---
\opage{95} %
\setcounter{footnote}{0}%
past\footnote{Der heilige Geist kann nicht wie eine Taube mit sichtbarer
Gestalt herabgefahren seyn, denn er ist Gottes unsichtbare überall wirksame
Kraft, wie hätte dieses sollen sichtbar werden? Bauer: \textit{l.~l.}
P.~\textit{225}. \textit{2} Band.}. And the apostles no doubt believed that
Christ, transfigured on the mountain, had spoken with Moses and Elijah; but if
one considers the matter a little more closely, it is easy to see that the
disciples had merely slept and dreamed that night and then, when they had been
awakened, rubbed their eyes and seen Christ lit up by the morning
sun\footnote{The same, P.~\textit{213}. \textit{2} Band.}. One will therefore
also see, on considering the true connection of the matter, that sacred history
is not very different from profane history. For the heathen, too, believed
they heard the voice of God in sudden thunder; their lawgivers, too, gave out
their laws as dictated by God; indeed, these too took omens from the flight of
birds (as John the Baptist did from that dove), just as they also readily
believed in virgin-born heroes with divine fathers. Nonetheless one must not
disdain the stories of Holy Scripture; for while many people awaken from
apparent death without one's being able to see the purpose of Providence in
it, one must praise divine Providence, because it has visibly so directed
things that the most perfect and most pious man, who is set before us as a
model for imitation, returned to life in order to raise up again the sunken
courage of his disciples; and even if we can hardly infer from his resurrection
to the immortality of our souls, we can
\stepcounter{footnote}\footnotetext{Ibid. P.~\textit{230}. \textit{2} Band.}%
\stepcounter{footnote}\footnotetext{Ibid. P.~\textit{192}. \textit{1} Band.}%
% --- p. 96 ---
\opage{96} %
\setcounter{footnote}{0}%
nevertheless now have this confidence in God, that we too, should we be carried
off by an all-too-early apparent death, will, if God so wills, return to this
life in order thereby to promote the cognition of a rational world order; and
although Christ did not ascend to heaven in a visible manner, but merely,
hidden by a cloud, went down the mountain by an unknown path, he has
nevertheless in this way received a far more glorious apotheosis than Romulus
and Drusilla\footnote{Therefore Paulus says, not without religious emphasis, on
the occasion of Luke \textit{21, 51}: Welch ein Gegensatz gegen dergleichen
Apotheosen lag in Jesu ἀναλήψις, auf welche \ldots\ das alte Römergesetz würdig
anzuwenden ist: \textit{Divos et eos, qui coelestes semper habiti,
\emph{colunto, et illos quos endo coelo merita locaverunt}, Herculem, Liberum,
Æsculapium, Castorem, Pollucem, Quirinum, et Horatii}:\par
\hspace*{2em}\textit{``Virtus, recludens immeritis mori}\par
\hspace*{2em}\textit{Coelum, negata tentat iter via:}\par
\hspace*{2em}\textit{Coetusque vulgares et udam}\par
\hspace*{2em}\textit{\emph{Spernit} (!) humum fugiente penna.''}\par
(Comment. üb. d. N. T. \textit{3} Theil. Leip. \textit{1805}).}, and rightly do
we admire his virtue and strive to follow his example; only let one beware of
worshiping him as though he were a god.

\parag{18}{Faith and Knowledge in Critical Theology.}

After criticism had thus sufficiently shown what it had in view, one would
expect that it would either have been expelled from the Church or would itself
have overthrown the Church's building; and yet, by the nature of the case,
neither was possible. Criticism does indeed contain a negation of faith, but
the question is at what point the negation lies. Inasmuch as empiricism and
subjective
% --- p. 97 ---
\opage{97}\setcounter{footnote}{0}%
rationalism urge the worldly and, if one may say so, purely human moments,
they at the same time negate the reality of the objects of faith, i.e., declare
them to lie outside the limits of their own cognition. But this negation is not
yet an assertion that that reality does not exist at all; for just as within
the limits of their own cognition they hold fast in a positive way to the
objects that belong to their sphere, so they can also very well acknowledge the
existence of a reality lying outside the limits of their cognition. --- But in
this way a new reconciliation between faith and
knowledge is prepared\footnote{How Kant, Jacobi, Schleiermacher, Friis and the
eclectic theologians have in different ways modified these theorems of critical
reason, it is not for us to develop more closely here, since our inquiry
concerns only the nature of criticism in general.}. One cannot, to be sure,
think the sum of the phenomena without assuming a unity of essence in them,
cannot think the finite without thinking the infinite, cannot think the visible
without thinking the invisible; but what can be conceived must belong to the
finite sphere, for where there is conceiving, there one must also be able to
separate the moments, and where the moments can be separated, there they must
also be finite. Cognition therefore subjects to itself the whole sphere of
finite objects and moments, but willingly renounces every conceiving of the
realm of infinite and invisible existence, although from the limitation and
restriction of its own sphere it concludes with certainty that this sphere must
have its ultimate ground of existence in that realm. If, therefore, access to
that realm is not to be closed to man, he must possess another faculty of
reception than cognition; but in this way there arises a new separation in
human consciousness itself, between cognition on the one side and will and
feeling on the other side.
% --- p. 98 ---
\opage{98} By the \emph{will} man is able, with the whole force of his subjectivity, to
surrender himself to the eternal and infinite realm, and by this surrender both
the finite and the infinite moment are preserved; for as the finite subject is
made subject to the dominion of the infinite subject, the latter is satisfied,
while on the other side finite man satisfies himself by surrendering himself of
his own free will. --- By pure \emph{feeling}, purified of all particular
affections, the same reconciliation falls to the share of the religious mind.
For in feeling there is an immediate identity of the subjective and the
objective moment; one then need only abstract from all particularities, and one
will then have the universal subject and the infinite object on which the
former is absolutely dependent. --- But from this there follows another,
extremely important separation, namely a separation between the religious idea
itself and the form in which it appears; and so great is the significance of
this separation that from it one can derive the great change that has taken
place in more recent theology.

For since forms and phenomena by their nature belong to the realm of the finite
and must therefore be conceivable, one can in truth conclude conversely that
God cannot be conceived by us, and that it cannot be the purpose of religion to
make God an object of our conceiving; we must then distinguish between God as
he is in himself, infinite and inconceivable, and the forms and phenomena in
which he has revealed himself for our salvation. It is thus at this point that
the newer theology departs from the older. In the older theology God finds his
true determination through his phenomena; he is himself present in them; they
must therefore be received as they are given to us in revelation; --- in
critical theology, by contrast, the infinite God cannot be contained in or
limited by any form; therefore the phenomena
% --- p. 99 ---
\opage{99}\setcounter{footnote}{0}%
and forms through which divine things are revealed to us must be regarded only
from man's standpoint. The visible and the invisible, the realm of cognition
and that of faith, therefore cannot, however much they may diverge, be kept
perfectly separate from each other; like France and Spain, Spain and Portugal,
they have the same borders in common, and since the concept of space cannot
fittingly be applied to this relation, we rightly say that in the finite realm
there are symbols and signs of, marks and traces of, the infinite realm. ---
Only by laying down these conditions did criticism, after so many fierce
battles, bring about a peace between \emph{faith} and \emph{reason}, according
to which everything that belongs to experience and logic was to fall to
critical reason, yet on the condition that it abstain from everything that
could in any way injure faith; on the other side faith was to content itself
with its own realm, not to begrudge free reason its earthly goods, and to
demand nothing of it except that it be ready for its defense and watch over
its interests\footnote{The inviolability of this pact is depicted by De Wette
in the following strong expression: Der Verstand
muß wissen, was er zu bestreiten hat, und was nicht; er muß nicht, vom blinden
Wiederspruchsgeist getrieben, seine Waffen gegen dasjenige richten, was er,
als hoch über ihm stehend, achten und schonen soll; und der Glaube, seiner
hohen Abkunft sich bewußt, muß es verschmähen, sich in die niedere Sphære des
Verstandes herabzusenken und diesen freigebornen Sohn der menschlichen
Vernunft in Sklavenfesseln legen zu wollen. Herrschen soll der Glaube im
Menschen, als der der himmlische Genius, gesandt, dessen dunkles Daseyn zu
erhellen und zu verschönern: aber die wahre Herrschaft lässt alle unter ihr
stehende Kräfte frei walten in dem ihnen angewiesenen Gebiet, und findet in
der freien Harmonie des Ganzen ihr schönes Ziel. Ueber Religion und Theologie.
Vorred. P. \textit{VII.}}. According to this agreement the character of the
theological method can
% --- p. 100 ---
\opage{100} be set forth in the following way: Everyone who has Christian faith
acknowledges that he has been brought up in this faith in the Church, since the
Christian idea of religion shines forth through the life all around him,
through family life and the sacred assemblies; but once he has reached a
certain maturity of spirit, he at the same time perceives that not even these
forms are perfect expressions of the objects of faith, since the Church's
ministers and teachers are human beings who can err like anyone else; he
therefore subtracts the universal idea from the forms the Church has handed
down, and seeks out for himself others that are more adequate; for though he
cannot himself set up true forms, he can nonetheless, when more perfect forms
present themselves to him, rejoice in them and acquiesce in them, since he
himself has faith. In his search he is supported on the one side by reason and
external empirical knowledge, on the other side by faith, internal empirical
knowledge. --- The Christian religion is a positive religion, which arose at a
certain time and under certain external conditions; one must therefore examine
its historical documents. We thus come to the testing of the authenticity,
axiopisty, and integrity of Holy Scripture, --- and although by these
investigations we do not exactly arrive at any mathematical certainty that
Scripture is the pure and unadulterated source of religion, we do arrive by
them at a high degree of probability of it. This probability becomes absolute
certainty when faith is added. For one can rightly conclude that as surely as
God has communicated a revelation to us, so surely does he also will that it
should come unadulterated to all by the surest way; since human tradition, the
oral tradition, as the theologians call it, can go astray, we hold ourselves
assured that Scripture is the true source of religion. But Scripture itself has
a double side from which it ought to be regarded, the historical and the
religious. From its historical side
% --- p. 101 ---
\opage{101} \setcounter{footnote}{0}%
Scripture must be an object of criticism's investigation, and one must grant
here that one may hold different opinions and voice doubts about much
concerning the origin of the individual writings and the disagreement between
the individual accounts, without thereby harming religion; in the religious
respect, by contrast, everything that Scripture teaches about God and our
relation to him is so deep and clear that it will always constitute the surest
foundation of faith\footnote{Since the Protestant Church grounds the authority
of the biblical books chiefly on the power with which they themselves bear
witness to their higher character and origin, it is, at least immediately, as
little touched by the doubts that scientific criticism raises against the
genuineness and integrity of the New Testament as one who is absorbed in glad
contemplation and enjoyment of a work of art lets himself be disturbed by the
reasons with which the art critic disputes its standing as a work of this or
that famous master. \textit{Dr.} Scharling. (Theol. Tidsskr., ed.\ by
\textit{Dr.} Scharling and \textit{Dr.} Engelstoft, \textit{4}th vol.,
\textit{1}st no., P.~\textit{16}).}. The religious person therefore, on his
first reading through Scripture, observes the general type of the doctrine, in
order under its guidance to pass on to an examination of the individual
dogmas, so that there comes to be a reciprocal relation between the universal
and the individual moments. This must be the right method for the treatment of
theology, since it attends equally to all the moments of truth; for it holds
fast to the positive and objective moment, --- one idea that pervades the
Church, one and the same need in the religious mind, one Holy Scripture,
finally, amid the vicissitudes of the centuries ---, and yet preserves, in the
midst of this unity of faith, the freedom of the individual; for within the
general limits everyone can form the doctrine for himself in his own way, and
choose for himself what best accords with his
inclination\footnote{\textit{Nemini totum est necessarium, alium alia pars
ad salutem ducit. Bengel.}}.
% --- p. 102 ---
\opage{102} \setcounter{footnote}{0}%
Just as Holy Scripture has two different sides from which it ought to be
treated, so it now goes with sacred history too. If one regards it from its
rational side, if we may say so, one will know the events set forth in the
accounts as such; then it is subject to the same relations as profane history;
if, on the other hand, it is regarded from its religious side, one sees at once
that it stands infinitely high above the latter. For although we do not
conceive history, we cannot, when it presents its sublime examples to us, avoid
receiving into our consciousness, through the religious view, an image of the
sublime idea itself, whereby the matter itself is confirmed and at the same
time takes root in our conviction. Therefore we should above all strive to
appropriate the image of Christ as it is painted in the N. Test., and we shall
then reap from it the most glorious fruits for a religious
life\footnote{For this way of treating sacred history De Wette has earned
extraordinary merit, something we gladly acknowledge, even if we do not approve
his metaphysical system. Ethics in particular owes him much gratitude, since it
was he who led the principle of morality back to Christ himself, whereas before
it contained nothing but general laws of conscience and reason, confirmed with
the seal of Christ and the apostles. Cf. Lærebog i den christelige Sædelære og
sammes Historie, translated by Mag.\ Scharling P.~\textit{3} §~\textit{4}
\ldots\ Christian (ethics) presents the (moral) tasks already actually solved
in revelation in definite exemplary images \ldots\ or the existing Christian
morals not merely in empirical-historical (positive) form, but also in a
distinctive material \ldots\ what is distinctive about it lies in the same
sphere in which the purely Christian doctrines of the Christian doctrine of
faith lie, in the domain of presentiment, where the idea and actuality
interpenetrate each other. P.~\textit{36} §~\textit{62} Christ is \ldots\ the
principle of wisdom, but primarily for feeling, which apprehends in him the
perfection of life. In order to become clearly conscious of this principle, we
must resolve it into a concept of the understanding.

The highest that Christ is \ldots\ consequently the principle of feeling that
transcends every concept is the idea of the Son of God \ldots\ the Son of God
is at the same time the Son of Man \ldots\ and the divine image is the pure
archetype of humanity.}. --- It lies in the nature of the case that
% --- p. 103 ---
\opage{103} this theology cannot disdain philosophy, much less exclude it from
its precincts, since its aim is precisely to bring cognition to ever greater
development; the spirit of man constantly goes forward, and what power could
halt it in its progress! But philosophy too is a human work, and we all know
that philosophers constantly waver between the most opposite theorems.
Therefore this method has also taken care to find a remedy against this
inconstancy of the philosophers. When the theologian takes up a philosophical
system, he must not intend to pursue the system itself through the countless
labyrinths of its thoughts, in order with the very power of dialectic to refute
its content and expose its metaphysical errors; for that is an insuperable
labor, and all the more troublesome since most systems are intricate and
obscure and in addition engage in so many abstract considerations about God
and his inner life, wholly superfluous for the development of Christian piety,
that the true theologian, in pursuing them, could easily run the risk of being
himself swept along into the abyss of philosophy and becoming a philosopher
instead of a theologian.

Since it is the Church, and above all Scripture, that communicates the
Christian faith to us, we consequently have in them certain fixed and
unshakable points of the doctrine of faith with which we can compare the
results at which the philosophical systems arrive by the way of philosophy. We
can then adopt these insofar as they accord with those articles of faith,
reject what stands in contradiction to them, and for the rest let the
% --- p. 104 ---
\opage{104} philosopher keep what belongs to him as such; for thus we give each
his due. --- How sure this procedure is can be seen most clearly by citing
examples of the test that the theologians, according to this method, make the
philosophers undergo. \emph{Spinoza} thus taught that God was substance, and
that the phenomena of the world were nothing but forms and modifications of
this substance; but by this pantheism the whole doctrine of the personality and
immortality of human beings is abolished, and as a consequence the whole of
ethics; his system thus conflicts both with faith and with Holy Scripture; it
ought therefore to be excluded from theology.

\emph{Fichte} taught that the human I was the author of all things, and thus
fell into atheism and crass Pelagianism; but a philosophy that conflicts so
openly with all sound reason and can have so pernicious an influence on
religiosity naturally cannot be tolerated in Christian theology.

Theology consequently stands with absolute authority over philosophy; and in
this way one can test ten philosophers in one hour.
\parag{19}{Examination of Critical Theology.}

We readily grant that this theology has taken all the moments of truth up into
itself; we allow that it is right that the Christian religion is a positive
revelation whose ideas cannot be found or developed by abstract, a priori
speculations; we admit that Holy Scripture is the rule for
% --- p. 105 ---
\opage{105} Christian doctrine, and so forth; --- but because it touches on all
the most important moments of truth, it is not yet settled that truth itself is
actually present in them, any more than one can say that a man lives because
all the individual parts of his body are present.

The task therefore becomes to examine how the moments are united and how they
are to be united, to test whether a true transition takes place between the
individual moments. And it will then soon become apparent that the alliance
that was concluded between faith and reason is not of a lasting nature, and
that they did not enter into this peace because it actually satisfied both
parties, but rather had seen themselves compelled to enter into an armistice in
order to avoid the danger that hung over them both. --- If we thus consider
first the relation between subjectivity and objectivity, we see that the
subjective freedom this theology extols is at once far too unbridled and far
too restricted and bound. One starts from the faith that lives in the
consciousness of the person brought up in the Church, thus from the general
influence that Christian life exerts upon everyone; one is thus agreed not to
go further into the particular modifications under which this faith appears in
individuals, for that each individual has, and has the right to have, his own
distinctive subjectivity is a settled matter. But since the same circumstances
and the same upbringing act differently upon different minds, and since no
doctrine can assert itself except through subjective faith, it is not objective
truth that corrects and transforms the mind, but rather the subjective spirit
that chooses and determines what seems to it to be the good and the true.
\emph{I will believe what most appeals
% --- p. 106 ---
\opage{106} \setcounter{footnote}{0}%
to me}\footnote{Compare with this \textit{Dr.} H. Martensen's remark on the
solipsism of critical theology: ``Since, then, the idea of God first acquires
validity through practical reason or through absolute feeling, not conversely,
it loses the character of the constitutive, and if one ascribes such a
significance to it, this happens only nominally; for what constitutes and
\emph{determines} everything is surely rightly to be regarded as the
\emph{actual} God, the omnipresent Godhead in consciousness, even if it does
not receive the name of God.'' Selvbev. Autonomie, translated by Petersen.}.
And yet this freedom is not very far from being servitude; for precisely
because subjective nature is urged to such a degree that it keeps one element
for itself as its incontestable property, namely the right to preserve its
immediate distinctiveness inviolate, objectivity on the other side determines
and fixes its limits all the more firmly, so that the positive moments, which
are given in a wholly external way, demand an external obedience; but in this
way the objective, positive elements that remain do not become flesh and blood
in us. This self-contradiction will emerge still more clearly in what follows.
--- As a great merit of this theology, one emphasizes the separation between
the essence of faith and its external form, according to which Christian faith
is indeed to be sought in the Church, but cannot acquiesce in its forms. The
ecclesiastical forms, it is said, are not perfectly adequate. But what is an
adequate form? If one means by it such a form as immediately expresses the
divine essence itself, then according to this theology there can be no
adequate forms at all; if, on the other hand, one means such forms as preserve
and indicate faith in its truth, one will hardly be able to avoid the dilemma:
either to assume that the Church is in possession of the true and adequate
forms for doctrine, and then the Church's doctrine must have the same
% --- p. 107 ---
\opage{107} validity as the doctrine of Scripture; --- or to take the Church's
forms and modes of expression to be erroneous in a greater or lesser degree;
but then the faith that the Church communicates must also be more or less
distorted, and one will then also be unable to arrive at any true
interpretation of Scripture, since such an interpretation requires that the
Church have handed down to us the true faith. An attempt to hold fast to the
faith the Church has called forth in us, while rejecting or doubting the
individual propositions of the Church's doctrine, will prove the truth of this.
Can one who removes all definite representations of faith, and then reads
Scripture doubting everything that faith has taught him about God, Christ,
etc., really have faith? Must he not then also doubt the existence of God, the
redemption and salvation of human beings through Christ, --- in other words
begin his study of Scripture with an absolute negation of faith, something that
conflicts with all reason? If one considers this, one can hardly deny that the
faith received from the Church is bound up with the general articles of faith,
so that it is impossible for anyone to take up Scripture without preconceived
opinions about the truths he will find expressed in it. The main question thus
becomes how the fundamental articles of faith can put off the concrete but not
perfectly true forms of the Church and put on the authentic forms of Scripture.

The Church teaches that God exists; but what attributes he has, and what
feelings he bears toward us, Scripture is to teach. One then demands of whoever
would read Scripture in the right way that he abstract from all God's
attributes insofar as the Church has described them, and find out from
Scripture itself the concrete concept of God; but it is impossible to think God
without having a certain representation of him, and the Church's doctrine of
God's essence will thus always, voluntarily or involuntarily, make its
influence felt on the reader's spirit.
% --- p. 108 ---
\opage{108} For if he is satisfied with the description of God that the Church
has given him, he will find the Church's God again, in complete likeness, in
the God of Scripture; if, on the other hand, before he reads Scripture, he
departs in one point or another from the Church's doctrine of God, he
nevertheless cannot negate the Church's doctrine without a positive concept of
God appearing with greater or lesser clarity in his consciousness. (Thus if
someone cannot in the negative respect acknowledge the Trinity of God as it is
taught by the Church, he must nonetheless in the positive respect worship
another God, whether it be that of the Jews or of the Muhammadans, or rather,
whether it be God as he is depicted by the Jews or by the Muhammadans; and if
he has the Jewish or Muhammadan faith, it goes without saying that he also
cannot find the triune God in Scripture, since precisely Christian faith,
according to this theology's own principle, is the first condition for the
right reading of Scripture). But the Christian reader, it is said, must not
urge, or with presumptuous self-confidence hold fast to, his preconceived
concepts of divine things, even if he cannot wholly abstract from them; his
spirit must have a willing mobility, always to adjust its subjective opinions
to the statements of Scripture: he must have the true method of
interpretation. In what, then, does the method consist by which we are to seek
to find out the meaning of Scripture?

First of all one must clear away the obstacles that the difference in language
and in time puts in the reader's way; and how great the merits of the learned
critics and philologists are in this respect is well enough known. But once
one has examined the general conditions of speech and action to which the
sacred authors were subject, added the necessary psychological remarks, and
thus come to know with fair certainty the meaning of the individual words and
expressions
% --- p. 109 ---
\opage{109} \setcounter{footnote}{0}%
that the authors must be assumed to have wished to put into them, --- one is
still very far indeed from having penetrated to the eternal significance of
their subjective opinions and determinations. Here, now, a double difficulty
arises; for, first, the words and expressions of Scripture are nothing but
imperfect designations of the highest, divine objects; and, second, the
particular circumstances of that age made it a necessity for the authors, for
both inner and outer reasons, to avail themselves of a mode of expression
wholly peculiar to those times. The whole yield of such acute exegetical
investigations is then confined to finding out how the sacred authors
expressed their objects; how these objects themselves are to be thought, they
cannot teach us, for since the difference between the idea and the forms in
which it is expressed is so great, it was conceivable that the apostles had
set forth the doctrine in this or that way, while we cannot think it in the
same way\footnote{An example of this we find in the doctrine of the devil. As
regards the principle of separation itself, cf.\ \textit{Dr.} Hagenbach:
Encyclopædie u. Mthdl. der theologischen Wissenschaften P. \textit{174} and
\textit{175}. Sache der Auslegung ist es, in Stellen, wie Rom. \textit{III,
24. 25. IV, 25. V, 9. 10. VIII, 34. 1} Cor. \textit{XV, 3.} Eph. \textit{I, 7.
1} Pet. \textit{I, 18. 19. II, 29. 1} Joh. \textit{I, 7.} Hebr. \textit{VII,
26. IX, 28} u.~s.~w. die Versöhnungs= und Opfertheorie als eine den
Schriftstellern des N.~T. geläufige Idee factisch anzuerkennen. Aber
verschieden kann die weitere Erklärung der Stellen ausfallen, je nachdem man
sich entweder damit begnügt solche Ideen in ihrem historische Zusammenhange
mit dem alttestamentlichen Opferwesen u.~s.~w. aufzupassen, oder je nachdem
man sie auf ein tieferes religiøses Bedürfniß nach Erlösung zurück zu führen
und in ihrem innersten, geheimnißvollen Grunde erfassen zu müssen glaubt.}.
One must thus grant, as far as dogmatics is concerned, that these
investigations have in the negative respect removed many a false
% --- p. 110 ---
\opage{110} application of individual passages of Scripture, but in the
positive respect can accomplish nothing at all; if one would find the dogmatic
ideas themselves, one must therefore strike out on an entirely different path
for one's investigation. We thus come here once again to the question from
which we started, namely how we can find out from Holy Scripture the positive
forms for faith.

Interpretation is to move forward step by step from the easier and more
comprehensible expressions in Scripture to the more difficult and obscure. But
since the immediately evident concepts are almost always hollow, empty, and
therefore meaningless, the result of such a comparative procedure is only all
too often that the deeper ideas vanish entirely. Thus if it is asked what is
meant by Christ's sayings (John \textit{15, 1, 4}):
Ἐγώ εἰμι ἡ ἄμπελος, ὑμεῖς τὰ κλήματα, and (v. \textit{9. 10}): Ἐὰν τὰς
ἐντολάς μου τηρήσητε, μενεῖτε ἐν τῇ ἀγάπῃ μου, one is in the greatest
uncertainty as to where the development should start: whether one should
assume such a relation between Christ and his disciples that love could have no
true moral force, and Christians no true love, unless his life were, in the
proper sense of the word, the source of their life, unless they drew their
whole strength from his marrow, --- or whether one should say: what a real
connection between Christ and those who believe in him means, we do not
understand, but what love and obedience to Christ's commandment are, everyone
understands; one cannot then take this and similar expressions for anything
but images drawn from an Oriental circle of representations. Just as uncertain
is the interpretation when Christ (John \textit{17, 11}) says that he is united
with the Father (ἵνα ὦσιν ἓν καθὼς ἡμεῖς); for one cannot know whether it is
the meaning of Scripture that there is an immediate identity between the
Father and the Son, and that Christians, therefore, when they have faith, have
entered into the
% --- p. 111 ---
\opage{111} \setcounter{footnote}{0}%
same real relation of unity with God and Christ, --- a relation
incomprehensible to all who do not, by being born of God, become θείας κοινωνοὶ
φύσεως ---, or whether the meaning is that Christ and the Father are one in
the same sense as Christians and Christ, namely by a moral bond of unity. What
such a moral union means, anyone who has sound reason can conceive; one ought
therefore to stop at this concept, for everything that aims at a substantial
unity is said only in a figurative and pregnant sense and has no significance
for religiosity as such.

The investigation of the concept of the Holy Spirit fares the same way. It is
described (John \textit{14, 16}) as a person, and Christ says of it (John
\textit{14, 26}): ὑμᾶς διδάξει πάντα, and (John \textit{16, 13}): ὁδηγήσει ὑμᾶς
εἰς πᾶσαν τὴν ἀλήθειαν (cf.\ \textit{1} Cor. \textit{2, 10. 11}); in Luke, by
contrast, it is called a δύναμις ἐξ ὕψους; one can therefore interpret this in
a double way. Those who believe that the Holy Spirit is a person do not deny
that it is also a power, but only maintain that it is more than a power; those,
on the other hand, who start from the immediate conceivability of the thing
declare that they cannot understand how the Spirit can be a particular person
in God, and think it sufficient to take it for a divine power without
determining its nature more closely\footnote{Numerous examples of this circular
course in argumentation can be taken from Bretschneider's Dogmatik; cf., e.g.,
on the Trinity P. \textit{495}. Handbuch der Dogmatik, \textit{1}ster Band,
\textit{2}te Aufl., Leipz. \textit{1822}; on the person of Christ §~\textit{140},
P. \textit{167}, \textit{2} B.}.

To this is added, further, what dogmatics shows: that every doctrine Scripture
sets forth contains moments that are mutually opposed, indeed even directly
cancel each other. Thus Scripture teaches, for example, that God is not
% --- p. 112 ---
\opage{112} merely loving and merciful, but also that he is angry with those who
have sinned. One cannot then know from Scripture whether one should urge
absolute mercy and take God's wrath only for a subjective representation of
human beings, or whether wrath is an objective concept that in actuality limits
the infinite mercy, or whether, finally, one should think God's nature in such
a way that he is indeed \textit{ad intra} the Blessed One, but \textit{ad
extra}, on account of the sins of the world, enters into such a relation of
atonement, of wrath and love, that his infinite love, through the incarnation,
removes the objective wrath from those who believe in Christ. One must
therefore here apply the analogy of Scripture and look to Scripture's doctrine
of the atonement through Christ. It is, however, so far from being the case
that the passages dealing with Christ untie the difficulty of the knot that
they rather make it still more intricate. For Scripture does teach that
Christ's death is a proof of God's love, but it also teaches that the love of
God and of Christ shows itself chiefly in this, that the punishment for our
sins is transferred to the crucified Christ; if one now defends the
objectivity of wrath, one arrives at the dogma of \textit{satisfactio
vicaria}; if, on the other hand, one cannot assume that God is angry, one must
refer the whole doctrine of the atonement to a change in human beings'
representations of God. The concept of wrath is thus dependent on the concept
of Christ's \textit{satisfactio}, but the latter is in turn dependent on the
former; one sees from this, then, that the one analogy always requires another
analogy on which it can rest, or in other words: analogy itself shows that one
cannot acquiesce in any analogy.

It will perhaps be objected that this economy of Scripture is the best proof of
the truth of the remark one so often hears, that God will not give human
beings any objective
% --- p. 113 ---
\opage{113}\setcounter{footnote}{0}%
insight into divine things, so that all the difficulties we mentioned before
arise precisely from our having ascribed far too great a significance to the
objective doctrine, --- and that Scripture itself confirms this opinion by
presenting God's nature, his love, and the atonement in images and general
concepts adapted to men's powers of comprehension, so that, strengthened by
the doctrine of the forgiveness of sins, they might be ready for every good
work. --- But have we not thus come back to the very same point from which we
set out, to the general articles of the Church?\footnote{The uncertainty in
such a mode of interpretation can be shown in several ways. It says, for
example, in John \textit{4, 24}, πνεῦμα ὁ θεός, but in John \textit{1, 5}, ὁ
θεός φῶς ἐστι; now if someone could not understand what it means that God is a
spirit, but believed he could more easily grasp that God is a light, might one
not then be tempted to believe that the God of the Christians and the God of
the Magi were the very same God?}. The whole result of the laborious work and
manifold investigations of criticism is thus the cognition that Scripture
itself proves that in dogmatic respects it can teach us nothing other than
what we have already learned from the Church. The critics, then, with their
great exertions, have accomplished nothing but to fall into contradiction with
themselves, for they promised us that the reading of Scripture would furnish
us with a more concrete form of faith, if only we on our side brought faith
along in its pure generality. We are thus called upon to investigate more
closely whether it is also actually in the true interest of faith, as they
recommend to us, to renounce every conceiving of the objective
objects.~--- It is evident that a faith in something altogether unknown is
impossible. I thus contradict myself if I say that I believe in the Creator of
the world while at the same time declaring that I can know nothing whatever
about the manner in which the world was created. For if I investigate what
% --- p. 114 ---
\opage{114} \setcounter{footnote}{0}%
creation is, I cannot avoid declaring creation to be an act of God's absolute
freedom; I then think, as far as possible, absolute freedom; but this
requires that I must have a representation of freedom, and if now this
representation, this concept, is false, then that in which I believe is also
false, i.e.\ I have a false faith. And if someone professes that he believes
Christ to be at once in possession of a human and a divine nature, he deceives
himself if he at the same time declares that he cannot conceive a union of
them; for insofar as he has a concept of the divine and human nature, to that
extent he can also think them united with one another; if, on the other hand,
he has no opinion at all about the properties of either of these natures, so
that the human nature is for him an unknown \textit{x}, the divine an unknown
\textit{y}, --- then surely he believes in a union of \textit{x} and
\textit{y}, believes in something of which he himself does not know what it
is, i.e.\ he believes nothing\footnote{The difference between critical and
speculative dogmatics does not consist in this, that only the latter, not the
former, recognizes the necessity of inner experience: they both assume that no
doctrine can be developed without regard to the positive elements of the
spiritual life; they both strive to give the representations and concepts
inner firmness. But the difference lies in the different way in which each of
them makes use of the same positive elements. Critical theology treats the
representations and ideas only negatively; therefore, when, after citing the
passages of Scripture and passing judgment on the various ecclesiastical
opinions, it has stated how creation is not to be thought, it contents
itself with that, leaves the more precise representation of it to each
person's subjective view, and passes on from there to treat another dogma in
the same way. And when the critics have in this way touched upon all the
individual scattered fragments of the elements without letting them take up
and permeate one another, they are highly astonished that speculative
theologians can think of speaking of an objective conceiving of divine things,
as if the elements of such a theory could only be snatched out of thin air.
``How can anyone want to think how creation took place,'' it is said; ``does
it not lie outside all experience! Centuries lie between the beginning of the
world and us; --- can any man experience what lies centuries back in time!''
--- Indeed, were such leaps of thought necessary, then not only would
metaphysical investigations be exceedingly unsatisfactory, but philosophy
would become more difficult day by day; for time hastens past us without
pause, and we are removed further and further from the point of beginning! But
the matter does not stand so; true philosophy, to be sure, never ventures to
overstep the bounds of human experience; but it makes a different use of the
positive elements of life. What makes it possible for us to have a
representation of creation and freedom is the very same thing that makes it
possible for us to arrive at the speculative concept of creation. Critical
dogmatics touches only in passing on the concept of absolute freedom, but
speculative dogmatics, precisely because it is science, demands that we pursue
it to its utmost limits, and it is along this way that we arrive at the
metaphysical concept of creation. If, therefore, one will purify in the fire of
dialectic the positive moments which critical dogmatics leaves undetermined,
one arrives at speculative dogmatics.}.

% --- p. 115 ---
\opage{115} One often hears that Holy Scripture does not constitute a philosophical
system, that it is popular; but this remark, as trivial as it is true, is not
taken in the same sense by everyone. For if those who speak of the simplicity
of Scripture mean thereby to prove that one may leave the objective forms of
doctrine unilluminated, provided only that one, according to the greater or
lesser enlightenment of one's age, engages in reflections on pious feelings
and moral precepts, then surely they act contrary to the whole economy of
Scripture itself. For Scripture itself calls precisely for taking an entirely
opposite path. It is true, to be sure, that Scripture gives the simplest
depiction of divine things, but one must equally remember that its development
of piety and the ethical life bears the very same stamp of simplicity. So far
is it from
% --- p. 116 ---
\opage{116} proceeding from human subjectivity that, on the contrary, it derives all piety
and morality from this, that the objective God has revealed himself in a
visible manner and by his Spirit has implanted in men a new principle of
\emph{wisdom} and life. Scripture, then, knows nothing at all of that
separation between the objective and subjective moments of cognition. Or can
anyone perhaps suppose that Paul, for example, in setting forth the doctrine
of predestination, thought to himself as follows: ``I do not, to be sure, know
what God is in himself, or what counsels he has resolved upon, but since the
idea of predestination is an excellent means of impressing modesty upon the
Jews and imparting comfort to the Christians, I will nonetheless set up this
doctrine!'' --- μὴ γένοιτο! If, therefore, you will be a child in
understanding the objective objects of the apostolic age, you must also
regard yourself as a child with respect to subjective, critical reflections!
But if your age calls upon you to undertake investigations and solve
problems, you must not, directly against the example of Scripture itself,
believe that you can do so without regard to the objective realm of divine
things!

The same indeterminacy, --- a consequence of criticism's external juxtaposing
of the individual moments ---, must naturally show itself in the way in which
they treat sacred history. Criticism can divide the whole of sacred history
into two parts, of which the first comprises all the accounts of the holy
persons that are of such a nature that one believes one can still today point
to examples of the same kind (e.g.\ that Christ went from Judea to Galilee,
that he addressed the people, that he showed great love toward all men).
Events of such a nature are found in profane history just as well as in sacred
history. To the second class, on the other hand, one must reckon all those
events
% --- p. 117 ---
\opage{117} that are peculiar to sacred history, namely those in which God is immediately
present. With regard to these one has two demands to make: first, it must be
shown that these events actually took place --- otherwise they would not
belong to history at all ---, and secondly, the supernatural ideas must not be
attached from outside to events of an entirely ordinary nature, but must be
derivable from the events themselves; the supernatural and the natural moment
must coincide, the former must appear as really and distinctly as the latter.
But this entails the greatest difficulty for the interpreter. For the
supernatural form in which the events, e.g.\ Christ's birth or resurrection,
appear makes it impossible for criticism to find out the purely natural
moments, and even if it were able to bring these to light, it could still not
do so without violating the covenant it has made with faith; it must
therefore leave the natural moment in the narrative standing altogether
undetermined. And as for the supernatural moment, criticism does not, to be
sure, always deny the possibility of miracles in general; for according to the
principle from which we set out, the supernatural realm, and therefore also
miracles in general, belong to the domain of faith; but since the finite
cannot contain the infinite, and the invisible cannot be present in the
visible, the truth of the miracle will nonetheless, in the individual cases,
always be subject to great doubt. Now since one cannot acquiesce in the way in
which the extraordinary event is narrated in Scripture, the supernatural
moment too is reduced to something altogether undetermined. But in this way
this part of sacred history vanishes from the realm of history itself. The
positive moments that remain belong to the outward garb, which is owed solely
to the religious imagination. The visible form of the Son of Man
% --- p. 118 ---
\opage{118} could not in actuality be the true expression of the idea of the Son of God;
that Son of God was then only a thought conceived in the religious minds of the
apostles.

If the matter stands thus, one can hardly deny that the predominant
characteristic of the critics is an external combination of positivism and,
if we may put it so, liberalism; they thus stand on the standpoint of the bad
middle way, a standpoint that cancels itself. This shows itself also in the
fact that, when a system which, as such, can be tested only by way of
thinking sets itself in opposition to criticism, one not seldom hears the
critics appeal to the authority of Holy Scripture, hears them invoke the bare
name of Scripture against the speculative dogmaticians, drawing a sharp
distinction between the word of God and that of man, just as the Catholics
of old, in the struggle against Luther, constantly appealed to the Fathers.
\emph{It is true because it stands in Scripture.} --- But scarcely is there
talk of freedom in critical investigation before all the scruples of
positivism at once vanish, before they concede that it makes little
difference whether Scripture was written by the apostles or by other pious
men, before they frankly and openly distinguish between the peculiar form of
Scripture and the word of God. --- \emph{Scripture is holy because it speaks
the truth.}

We by no means doubt that a reciprocal relation obtains between the testimony
of the truth concerning Scripture and the testimony of Scripture concerning
the truth; but we very much doubt that altogether opposite principles can be
united with one another. Either we must, like the Fathers, unconditionally
let our judgment be subject to Scripture and, as a consequence, also assume
a verbal inspiration; or we must everywhere receive the truth in the name of
truth, and then we have no reason to assume that God, from the time of Christ
right up to Trajan, spoke
% --- p. 119 ---
\opage{119} \setcounter{footnote}{0}%
immediately through certain pious men, but from that time on let men's speech
about divine things be independent of any divine influence. --- The same
uncertainty shows itself when the talk is of the relation between the general
idea of Scripture and the individual passages of Scripture. They themselves
urge, where it suits them, the individual statements of Scripture, and always
say, when one puts forward something that conflicts with their views: ``Where
is this written? There is not a word about it in Scripture!'' and thus in this
case dispute κατὰ ῥητόν; --- if, on the other hand, one cites definite
passages of Scripture against them, they always know how to distinguish
between the thing itself and its outward garb, defending their freedom of
thought with the assertion that one may only reason from the general idea and
the whole economy of Scripture; here, then, they dispute κατὰ
νοητόν\footnote{Examples of both these ways of disputing are found in
Bretschneider: \textit{De obedientia activa} §~\textit{157}. (Pag.
\textit{285}. \textit{2} B.).}. But such a procedure is a
``\textit{mystificatio veri}'', which can easily bring upon the authors of
this school the suspicion that they said one thing but meant another; for the
extraordinary reverence for Scripture often looks like a mask, put on so as to
be able, under its cover, to give subjective freedom all the more scope.

\parag{20}{Continuation. Transition to Panlogism.}

Tell me which metaphysics you follow, and I shall tell you, with the
exception of your purely private opinions, which dogmatic interpretation of
Scripture you will put forward! If you say that you follow no metaphysics at
all,
% --- p. 120 ---
\opage{120} \setcounter{footnote}{0}%
we do not believe you; if, on the other hand, you say that you yourself do not
know which metaphysics you follow, we will believe you\footnote{Es hilft der
Theologie nichts, sich dagegen zu sträuben, zu sagen: sie wolle nichts von
Philosophie wissen, es seyen Philosopheme, also auf der Seite liegen zu
lassen. Sie hat es immer mit Gedanken zu thun, die sie mitbringt; und diese
ihre subjectiven Vorstellungen, Gedanken, ihre Haus= und Privat=Metaphysik
sind dann die Reflexionen, Meinungen u.~s.~f. der Zeit \ldots\ diese
allgemeinen Vorstellungen sind nichts, als was sich auf der Heerstraße
findet, was auf der Oberfläche der Zeit umherschwimmt. Hegel: Geschichte der
Phil. (W. \textit{15} B. Pag. \textit{268}).}. It is not, then, by renewed
historical or philological investigations that criticism can hope to make any
progress, but solely and only by a change of its metaphysics. If, therefore,
criticism can defend its metaphysics through all its individual moments
against the objections of dialectic, then critical theology too will stand
firm and maintain its name as the highest form of science to which men can
attain in this life; if it cannot do this, we are entitled to seek out
another form for true science. In any case this much stands firm: that it
must not seek its defense in Holy Scripture, but must remember what Clement
says --- a saying to which appeal is often made in our time\footnote{Compare
Fr. Baader, Vorles. üb. Dogm. \textit{1}. H. Pag. \textit{17}. J. Martensen.
S. A. §~\textit{1}.} ---: ``\emph{One must philosophize, even when one wishes
to combat philosophy!}''\footnote{The truth of this has been acknowledged
also by the most brilliant and most celebrated critical theologians
themselves: e.g.\ De Wette: \ldots\ ich schmeichele mich mit der Hoffnung,
daß ich wenigstens die allgemeine Ansicht, welche für die Theologie von
entscheidender Wichtigkeit ist, klar genug gemacht habe. Und diese ist keine
andere, als \emph{die Unterscheidung der Verständigen, idealen und
ästhetischen Ueberzeugung}, welche ich für den Schlüssel der ganzen Theologie
halte. Relig. und Theol. Pag. \textit{IV}. compare also Pag.
\textit{162---172}.

Thus a critic, too, seeks to combat Weise with metaphysical weapons in his
review of the latter's book ``Die ewangelische Geschichte'' (Kritische
Prediger=Bibliothek v. D. J. F. Röhr, \textit{20}. Band, \textit{29}. Heft,
Pag. \textit{200---219}). He seems, to be sure, in censuring Weise's doctrine
of a substantial unity of Christ, men, and God, to appeal to the authority of
Scripture: ``Diese Einheit aber der Menschen mit Gott ist nicht eine
substantiale, ein Hineinbilden unserer creatürlichen Selbstheit in die göttl
che Substantz, es ist eine moralische nach der ausdrücklichen Lehre Jesu und
aller seiner Apostel'' \ldots\ but if we consider his objections more closely,
it is readily seen that they belong to metaphysics; for he first rejects that
substantial relation because he cannot conceive it: ``Was sollen wir uns aber
auch unter jenem Hineinbilden in die göttliche Substantz Vernünftiges
denken?'' Next, he is convinced that the substantial relation conflicts with
the true idea of God and therefore cannot serve to further morality: ``und
indem so der Rationalismus die pantheistischen Gedanken von einer
Menschwerdung Gottes, einem sich Hineinbilden in die göttliche Substanz
verwirft, weil sie einerseits unerweisbar und der reinen Gottes=Idee
wiedersprechend, desshalb aber auch andererseits für Beförderung der Moralität
gleichgiltig erscheinen \ldots'' And lest we should object that his assertions,
of which the first may seem to the reader more reasonable than the last, are
without real probative force, the author declares that he wishes every
investigation to be carried through philosophically; indeed he even assures
us that all such questions would long ago have been satisfactorily answered
had not speculative philosophy by chance displaced criticism for a time:
``Hätte man hier die Bahn verfolgt, welche der unsterbliche Kant und die
Anhänger seiner Schule betraten, die völlige Aussöhnung der philosophischen
Speculation mit dem reinen und einfachen Christenthume würde zum Heile der
ewangelischen Kirche bald erfolgt, aller unnützen philosophischen Grübelei
über Dinge, die wir nun einmal nicht zu erkennen vermögen, und die darum auch
für unser wahres sittlich=religiøses Leben keinen Werth haben, würde
vielleicht für immer vorgebeugt, und somit ein Höhepunkt deutscher
Philosophie erreicht worden seyn, den diese Wissenschaft in keinen unserer
cultivirten Nachbaarstaaten so bald erreichen dürfte(!!!). Da trat die
Speculation des absoluten Philosophie dazwischen \ldots\ Trotz dem aber werden
diese Extreme philosophischer Speculation ihre heilsamen Folgen haben; sie
werden darauf zurückführen, daß wir den Grund der Erkenntniß göttlicher Dinge
allein in den Anforderungen unserer sittlichen Natur suchen \ldots'' After
adducing all these weighty reasons he takes his last objection from
Scripture, disputing κατὰ ῥητόν, in order to give his remarks greater
emphasis and to reduce Weise to silence by the silence of Scripture: ``Wo
lesen wir in den heiligen Urkunden ein Wort von einem Hineinbilden des
creatürlichen Selbst in die göttlichen Substanz?''}.
% --- p. 121 ---
\opage{121} %
We are thus called upon to make the critics' metaphysics the object of a
metaphysical investigation. --- It follows
% --- p. 122 ---
\opage{122} from the very nature of criticism that it separates the positive moments from
one another (for κρίσις is precisely the activity by which one separates,
sets apart, the objects). The realm of cognition touches the sphere of faith,
but not in such a way that any identity of the two ever comes about; for what
makes them touch one another is the very same thing that at the same time makes
them repel one another. Each of them seems to acquiesce in itself within its
own bounds and to take up no foreign elements into itself: they are thus
indifferent to one another. Now were this indifference absolute, theology
would be altogether superfluous, indeed contrary to reason; for if theological
science has no other object than the realm of finite things, and if on the
other hand faith, exalted above all cognition, must itself look after the
interest of its objects, then faith too is secure within the bounds of its
own realm, then it can and must also apprehend the divine ideas on its own
account. But this mutual exclusion cancels itself on closer consideration.
For if faith, left to itself, is to present its objects in their order and
right relation, then the religious feelings too must stand in
% --- p. 123 ---
\opage{123} opposition to, and restrict and limit, one another, and in this way
cognition, which is the orderer and governor of all finite, limited objects,
will break through all the bounds of faith. On the other hand, the realm of
cognition is not sufficient unto itself within its own bounds either; for if
its finite elements are to determine one another mutually, the bounds between
them are canceled, since the finite as such is like the finite and cannot be
limited by anything other than its own opposite, i.e.\ by the infinite. It is
thus altogether certain that the infinite realm of feeling touches the finite
realm of cognition at every point, and that as a consequence they cannot go on
being indifferent to one another.

And yet, on the other hand, one can just as little hope that they will
furnish any limitation for one another, for they repel one another as often
as they touch one another; one might therefore rightly call the theology that
proceeds from this method bewitched, since it at once shuns and loves
thought. It is cognition that determines forms and shapes, so that anyone who
would hold fast to the objects of faith is compelled to appeal to it; but
since the theological self-contradiction holds not merely in general but also
in every individual case, and since the proposition that the finite cannot
contain the infinite returns in manifold ways, cognition always reaches
\emph{for} the objects of faith but in actuality never \emph{grasps} them. In
vain does it test the events, for natural pragmatism is the guiding rule of all
its investigations: if, following this path, it could prove that the events
took place, then surely everything supernatural in them would vanish; --- if,
on the other hand, it finds that they stand in contradiction to the ordinary
order of things, then surely it must declare them
% --- p. 124 ---
\opage{124} untrue. --- Cognition thus demands that events which are supernatural be seen
to have taken place in a natural manner, while it at the same time declares
that these realms are mutually opposed to one another. If, therefore, we wish
to intuit the objects of faith, we must renounce thinking; only on that
condition do they descend from heaven in limited, determinate forms. Since
the acknowledgment of divine things thus presupposes a giving up of
cognition, theology as a consequence stands outside the sphere of thinking;
indeed, the more often one thinks about this relation, the more entangled it
looks, and the entanglement goes on to infinity. Christ is divided in this
theology between the two realms; to the realm of feeling and faith belongs
the Son of God, to the realm of cognition, on the other hand, the Son of Man.
One would thus expect that these natures of Christ, owing to their relation
of mutual indifference, would isolate themselves from one another, so that it
would be no concern of faith and feeling in what manner the man Christ was
apprehended, and faith would not be needed in order to conceive his human
nature. --- But if we consider the matter more closely, we readily see that
the human nature has no other purpose than to determine the divine, and that
it cannot at any point whatsoever be held apart from it, since the divine
glory shines through it everywhere, so that in the man Christ the whole
Christ is revealed to aesthetic feeling: we thus seem immediately to behold,
think, conceive Christ. But as often as we now try to think this thought of
ours, the indifferent and opposed in the relation between Christ as man and
Christ as God returns again; the man Christ is then bound to the lower sphere,
while the divine Christ ascends to heaven? Cognition again urges the man
isolated from the divine nature, while feeling reaches for the God isolated
from the
% --- p. 125 ---
\opage{125} human nature: Christ becomes an ordinary man, an unnameable God: he vanishes
like a deceptive phantom of mist as often as we reach for him; and so we
stand struck by the Apostle's words: Ὠ ἀνόητοι, τίς ὑμᾶς ἐβάσκανε; οἱς κατ'
ὀφθαλμοὺς Ἰησοῦς Χριστὸς προεγράφη.

This difficulty has its ground chiefly in this, that profane cognition, which
cannot conceive the divine nature but grasps only the things that belong to
this world, remains behind untouched and cold, while feeling and will are
consecrated to Christ, in order through him to be sanctified and animated
with the fullness of the Godhead. As long, therefore, as the opposition
between Christ and the world persists, so long will that discord persist as
well; if faith is to be determined, it is compelled to take refuge in pagan
cognition. Since, then, feeling is infinite and therefore without determinate
shape, it is cognition that fixes the bounds and determines what lies within
its own domain and what lies within the domain of faith. It can then also
easily, when it sees fit, banish Christian faith from its realm and seize the
whole dominion. Once reason has become conscious of this its authority, it
will soon constitute itself absolute sole ruler and declare without
hesitation that all discrete moments, heavenly and earthly, supernatural and
natural, are to be sought within the bounds of its own thinking; but in this
way there arises \emph{objective rationalism}. --- Thinking is reason, reason
is God; but thinking is a negation of phenomena; thus the last trace of
positivism vanishes. When phenomena are negated, one has arrived at
speculative logic, that is, the concept of pure being, God as he was before
the creation of the world. But just as pure being cannot be except through an
actual phenomenon, so too God is not except through the world, and thereby
one has freed oneself from the
% --- p. 126 ---
\opage{126} last remnant of the older supernaturalism. The world, next, is a complex of
finite phenomena; considered immediately for itself, it is in opposition to
God. But since this opposition concerns things only in their immediate
existence, a perfect agreement between God and the world will set in when one
gives up every, if one may say so, crude immediacy. The individual moments of
existence are then, considered in and for themselves, nothing; they are to be
considered and judged only in their dialectical transition into other
moments; there is thus only \emph{a logical progress}. The revelation of
absolute spirit is, to be sure, posited as its final goal; but the absolute
criterion of this revelation is solely and only the immanence of ideality and
reality, essence and phenomena. One may thus rightly call this system
\emph{panlogism}, since its development is a movement from logic through
logical moments to logic.

The true apprehension of things in this their new, cosmic mediation requires
that one distinguish between the potential being of the principle and its
revelation. In the dark womb of the principle the immediate identity of
essence and phenomenon is present in hidden form, but only through its going
forth into its manifold, various moments does a true, concrete immanence come
about. Essence must therefore descend from its sphere of pure being into the
outer realm of phenomena; there is thus something more to be sought in
phenomena than bare phenomena, and herein lies the spur to further
advance. For insofar as essence needs the phenomenon, to that extent it
posits and preserves it; but when it has developed its existence through it,
as far as was possible within its bounds, it breaks down the bounds and
annihilates the outworn phenomenon in order to let a new and more excellent
one
% --- p. 127 ---
\opage{127} come forth in its place. And if one asks why the more excellent appears only
now, after the former, and is not present at once in the beginning, the
answer must be that the excellence of the second phenomenon is mediated
precisely through the existence and destruction of the first; for precisely
by positing and again destroying the first has essence gained new strength.
Amid the alteration of phenomena essence always holds fast to itself, and,
like substance, preserves in every new phenomenon it creates all that was
essential in every preceding phenomenon, while it destroys all that is
corruptible and inadequate.

The manner in which this development takes place differs in the different
spheres. In the lower realm of nature the power of essence to maintain itself
is weaker; hence the tautological repetition and uniform circular course
observed in the phenomena of this realm. Progress in its true development must
be sought in that realm of existence where, in the midst of the disagreement
between essence and phenomenon, there appears at the same time the firmest
bond between the two: where essence returns to itself through a negation and
restoration of phenomena: where the alteration of the outer moments is at the
same time an assimilation of the inner: where phenomena can be preserved in a
recollection, i.e.\ in \emph{the history of the human race}.

God, then, having gone forth from the realm of his pure ideality, has posited
visible phenomena in the realm of nature in order, through the negation of
these, to come to subjective spirit, i.e.\ to human self-consciousness. Our
self-consciousness belongs to two spheres. Bound to the realm of finite and
sensuous phenomena, we are sinners, forfeit to death; but raised up to the
sphere of the absolute essence through the negation of phenomena, we have God
present in us and are thus reconciled with him. \emph{The absolute}
% --- p. 128 ---
\opage{128} \emph{foundation of history is therefore religion, its absolute goal the
revelation of God.}

Since the whole cosmic mediation must thus be considered in the light of
religion, what we have here developed about the relation between essence and
phenomenon in general holds also of the bond by which religion unites men with
God. God is at first immediately present to human self-consciousness, so that
he can be pointed out visibly in a determinate, individual phenomenon, i.e.\
God is revealed. But in every revelation a twofold moment is to be observed, a
phenomenal and a spiritual. For the organ of revelation, the phenomenon, can
neither be acknowledged nor worshipped as such unless the same principle of
religion that is present in the phenomenon is also present in the one who
beholds the phenomenon. The idea of religion is therefore formed in every
beholder from both sides; on the one side he beholds the particular
phenomenon and attaches himself firmly to it; --- on the other side he seeks
to find out and explain to himself the essence of the phenomenon, and thus
negates the particular, contingent phenomenon in order to make the general,
the essential, that shines forth in it his true property. And here the
proposition holds: \emph{as God is in the phenomenon, so is he also in
consciousness.} Just as, therefore, no religion and no revelation of God can
be given except through a sensuous phenomenon, so on the other hand no
phenomenon is an immediate or absolute expression of the idea of God. Here,
therefore, a mediation must take place, in order that the empirical phenomenon
may be negated and a new, more perfect phenomenon, more satisfying to
religion, may by the idea's own power take its place. This negation and
restoration human self-consciousness undertakes according to a certain
instinct, so that
% --- p. 129 ---
\opage{129} \setcounter{footnote}{0}%
the same thing that causes the phenomenon to be embellished in the interest of
the Idea at the same time procures it credence as objective and real.

From this derive all the accounts in sacred history of God's revelations, and
the miracles; and one will now easily understand the agreement and the
disagreement between panlogism and criticism. For both insist in the strictest
manner on a pragmatic development of nature, which cannot possibly be
interrupted by the intervention of a higher power\footnote{Diese Ueberzeugung ist so sehr Bewußtseyn der
neuen Welt geworden, daß im wirklichen Leben die Meinung oder Behauptung, eine
übernatürliche Ursache, eine göttliche Wirksamkeit habe irgendwo unmittelbar
eingegriffen, geradezu als Unwissenheit oder Betrug betrachtet wird. Strauß.
L.~J. § \textit{14} Pag. \textit{91}.}; both call miracle into doubt. For the
true concept of miracle requires that an actual opposition must take place
between God's will and the laws of nature, so that the moments of nature
proceed from one another according to a certain law of necessity; for were the
objects of nature, in their general development, immediate expressions of
God's will and counsels, they would surely all be miracles, and then no miracle
at all would take place; but in the opposition there must at the same time be
an identity present, for had God not himself established the laws of nature,
and did he not know them, neither could he alter them. Faith in miracles thus
presupposes an identity of God's will and the laws of nature, but assumes at
the same time that in historical development this identity passes over into an
opposition, so that collisions can occur, --- assumes that in divine providence
in general there shows itself a care for the human race in particular,
according to which God, although our phenomenal world does not correspond to
God's own essence, where he finds it needful, uses for special purposes and
% --- p. 130 ---
\opage{130} glorifies the natural phenomena: so that not only do human beings
believe that they see in one or another phenomenon a special revelation of God,
but God himself also, according to a special counsel, presents himself to human
beings in this phenomenon.

Criticism doubts the possibility of miracles because God, although he is
perfect in himself, is so exalted that it would be beneath his dignity to
descend into sensuous phenomena, --- so wise and unchangeable in his counsels
that he would never alter the laws to which he has from the beginning bound the
development of natural objects. The sharp opposition between God's abstract
essence and the realm of phenomena thus makes it impossible for the critic who
inclines to deism to believe in miracles.

The panlogists, by contrast, posit an immanence of God and world; yet the
immanence holds only of the general world order, --- if natural objects are
taken in their singularity, the opposition asserts itself. That God himself
should in actuality be present in a special way in this or that phenomenon is
impossible for two reasons. First, God would then have to have a supramundane
self-consciousness, reflected into itself, in order to be able to think of a
special revelation --- but this obviously conflicts with the cosmic mediation
we spoke of earlier ---, and second, the infinite God would have to be present
in a finite phenomenon; but this is equally impossible. The transfiguration of
phenomena, the thought of special revelations, cannot then be referred to any
objective activity of God, but can only derive from subjective, human
representations. Now although the sacred accounts must thus be assumed to
derive solely and entirely from subjective religiosity, they could nonetheless
not arise without some sufficient ground, since they hold good and find credence
among many at once. The ground and occasion is thus
% --- p. 131 ---
\opage{131} \setcounter{footnote}{0}%
the objective moment in the miracle accounts, whereas the forms in which they
are clothed are owed to the spectators and listeners: one may therefore rightly
call such accounts \emph{myths}.

In myths, then, there are two moments to hold fast, an objective and a
subjective; in manifold ways these pass over into one another, and it is in the
development and assessment of these moments' mutual transition into one another
that criticism diverges from panlogism. For the critics commend everyday
experience and regard objects from outside; they therefore want a tangible
actuality; for since the concept of mediation is obscure to them, they cannot
understand why one does not straightaway, from the very beginning, make this
separation between the objective actuality and the subjective representations
of it. They do acknowledge that it was necessary, on account of the spiritual
immaturity of that age, or at least in keeping with its peculiar spiritual
standpoint, to spread such traditions; but deism often drives them so far that
they would rather declare those accounts products of cunning and deceit than
assume God to be their author\footnote{Cf.\ Strauß. Leb.
Jesu (\textit{1838}) Einleit. §§ \textit{5}. \textit{6}.}.

In panlogism they are judged in an entirely different way. The separation that
criticism insists on between the idea of God, as it hovers before the pious
mind, and God himself, conflicts completely with the principle of panlogism.
If, therefore, criticism concedes that the traditions of miracles and
revelations have issued from an idea of God, panlogism proves that they have
come about through an activity of God immanent in human consciousness:

But since these revelations of God, of whatever
% --- p. 132 ---
\opage{132}\setcounter{footnote}{0}%
nature they may be, thus constitute a necessary stage in the history of
revelation, they also belong to the cosmic mediation and are not to be
despised. It is so far from being the case that they must be regarded as
invented fables, that they rather hand down to us phenomena restored and
embellished in the interest of religious truth. Myth thus contains, according
to its concept, not only a moment of revelation but also a moment of
illusion\footnote{Wenn man die
Religion im Verhältniß zur Philosophie bestimmt als das Bewußtseyu desselben
absoluten Inhalts, aber nicht in Form des Begriffs, sondern der Vorstellung: so
ist leicht zu sehen, daß nur unter und über dem eigentlichen Standpunkte der
Religion das Mythische fehlen kann, innerhalb der eigentlich religiøsen Sphäre
aber dasselbe wesentlich und nothwendig vorhanden ist. S.~S. Pag. \textit{98}.
§ \textit{14}.}. Those who handed down the myths to us believed that the
accounts were objective, i.e.\ that God himself, outside human consciousness,
had taken on those finite forms. But this error is revealed and removed in the
course of time; for every subsequent historical stage always surpasses the
preceding one, and the revelation of the divine Idea becomes ever more and more
perfect\footnote{Wie aber Bildung
überhaupt Vermittlung ist: so wird die fortschreitende Bilder der Völker auch
der Vermittlungen immer deutlicher bewußt, welche die Idee zu ihrer
Verwirckligung bedarf; und so erscheint jene Differenz der neuen Bildung und
der alten Religionsurkunden \ldots\ Das Göttliche kann nicht so (theils
überhaupt unmittelbar, theils noch dazu roh) geschehen seyn; oder das so
Geschehene kann nichts Göttliches gewesen seyn --- so wird die Differentz sich
aussprechen. Strauß L.~J. § \textit{1}: Pag. \textit{2}.}. One now sees that
the things one has hitherto believed were false, and this becomes a source of
sorrow and grief; but one also sees, when cognition has reached its full
development, that it is the same activity of God which was formerly the source
of
% --- p. 133 ---
\opage{133} the myths that has now in turn brought about the cognition of their
untruth. It is so far, then, from being the case that God abandons us when we
have discovered the error, that he rather now draws nearer to us, and this is
the source of consolation and joy. For in this too panlogism diverges from
criticism, that the former has an ever-continued revelation of God, while the
latter, after denying the sacred records their validity, has at the same time
lost the last traces of an actual divine presence. --- For God does not exist
outside the human race: the spirit of God and the spirit of human beings always
preserve their identity, God goes out into human beings and human beings return
to God.

What thus holds of sacred phenomena in general can naturally also be applied to
the most excellent phenomenon of the absolute religion.

For the right understanding of the mythical Christ, i.e.\ of Christ as he is
depicted in Holy Scripture, the following three moments come into
consideration.

First one must consider the phenomenon as it could be apprehended through
everyday experience, --- imagine Christ as he would have appeared to the
enlightened people of our time, had he arisen among us. Considered from this
side, he must be assumed to have been a human being who was perfectly clearly
conscious of the idea of religion; otherwise he would not have brought about so
great a change in the religious sphere. --- But that he apprehended the absolute
idea of religion by no means entails that he was an eternal personality,
endowed with power to work miracles; indeed, such an assumption even conflicts
with the idea of religion itself. For insofar as Christ was a human being, he
could neither, as a single individual, contain the whole fullness of God in
himself, nor reveal it to his contemporaries. The presentation of Scripture can
% --- p. 134 ---
\opage{134} \setcounter{footnote}{0}%
thus give everyday experience no more than certain fragments and outlines of the
true image of the Messiah\footnote{Cf.\ Strauß: L.~J. Schlußabhandlung.}. ---
He appeared in the fullness of time, when the age already had a rich apparatus
of prefigurations and figurative designations; now since that religious
principle which stood clear before Christ's consciousness came in and, as it
were, gave these elements a new life, and since the particular phenomenon in
the service of the Idea became the prey of death, the religious idea conceived
from this had so great a power in the consciousness of those who had been the
great Master's disciples, that they restored his life and his personality anew,
and thus handed down to the believers the fullness of the Idea made intuitable
in a phenomenon which, as such, could become an object of worship. We must
therefore, although that Christ whom the apostles have described to us is
composed of such opposed moments, concede that the people of our time would not
have been able to grasp the religious idea in an inward way had it not first
been made intuitable for us in this outward way. One must therefore not despise
that Christ, for though he died as a single human being, as a unique phenomenon
(κατὰ σάρκα), he rose glorified as a symbol of the human race
(κατὰ πνεῦμα), and we thus owe it to him that we have come to know the absolute
concept of truth. --- Everyone who is content with the God of panlogism must
therefore also approve the mythical treatment of sacred history. Although, then,
no reality remains save with an altogether momentary existence, and although
subjective and phenomenal life is continually swallowed up by the objective and
essential: yet God remains the same in the realm of ideas, yet the human race,
despite the mortality of individual human beings, vindicates its stamp of
eternity, yet the true
% --- p. 135 ---
\opage{135} Messiah takes his seat on the mount of panlogistic blessedness and
teaches his people:

``Blessed are they who believe in the speculative immanence of heaven and earth,
for theirs is the kingdom of heaven; blessed are they who cleanse themselves of
all prejudices about a God existing above the world, for in the pure mirror of
the Idea they shall see God; blessed are they who scorn their own personal
immortality, for their reward shall be great in the realm of memory!''

\parag{21}{The Metaphysical Testing of Panlogism.}

Should anyone object that nothing but neglect of the study of Scripture can have
called forth this development of the Christian religion, we answer that it is
rather acute, learned men, well versed in Scripture, who recommend it. They
have neglected none of the critics' introductory investigations; equipped with
the whole apparatus of sacred philology, they distinguish with the greatest
subtlety among all of Scripture's ζεύγματα and other figures; indeed, they even
know all the attempts to reconcile Christ's genealogies down to the last jot,
--- but all this does not convert them. For the following reciprocal relation
obtains between life and philosophy: life posits the moments as positive,
separated, but at the same time preserves them all at once, lets them all
subsist side by side, connects the individual moments with one another without
sublating their singularity; thinking, on the other hand, negates and restores
the phenomena of life, conceives through an intuiting, intuits through a
conceiving, and in consequence insists now on the separation of the elements,
now on their identity. And precisely herein is grounded the whole legitimacy
and illegitimacy of critical deism and of speculative pantheism.
% --- p. 136 ---
\opage{136} Among the Orientals, where the dull eye of the human soul can
discover no other opposition than the opposition between the multiplicity of
phenomena and the abstract unity of essence, everything finite is swallowed up
by the infinite; \emph{there}, therefore, the reconciliation also takes place
in the most immediate way. --- Consciousness, which thinks, is subject to
sensuous nature; the natural moment thus preponderates in consciousness, and it
is therefore by a certain immediate instinct that consciousness grasps the ground
of its existence. It can then raise itself above nature only by voluntarily
giving nature what belongs to it; it therefore returns through asceticism to
its original nothing (essence in its pure abstraction). The consciousness of
nature is thus here limited by the nature of consciousness. --- With Christian
life the case is quite otherwise; it permits no sinking into dreams of nature;
only through manifold position and negation of the given facts, only after
spirit had overcome the manifold difficulties that stood in its way, only then
could panlogism appear as the final result.

We have already remarked earlier (cf.\ §~\textit{7}) that the restoration to
which the dialectical stimulus leads consciousness differs according to the
differing nature of the phenomena that are negated. Yet the restoration is not
always such a reinstatement of the negated phenomena that these come to occupy
the same dignity they had before; for the dialectical restoration is at the
same time a testing. We must therefore investigate whether panlogism will be
able to prevent those moments of the ancient essence of the Christian religion
which it has bound with the chains of dialectical negation and thrust down into
the shadow-world of subjective representations from stepping forth anew, --- or
whether it is the eternal ideas of truth on whose ruins speculative philosophy
has built its
% --- p. 137 ---
\opage{137} \setcounter{footnote}{0}%
palace, and which will therefore, after a longer or shorter time, burst their
prison once more and avenge the injury done to them. Fear not, they will say, we
are good Christians! The God we preach, is he not the absolute spirit, which
negates nature in all eternity precisely in order to glorify it? Is not the God
who through the metaphysical trilogy unceasingly goes out of himself and
returns to himself the Triune? Here there can be no talk of a profanation of
revelation; we do not profane God, as though he ever shut up heaven out of
envy; no, the loving God constantly makes human beings partakers in his life!
What else, after all, does the personal unity mean than precisely this inward
union? Yes, Christ is religion itself; he was a historical person in whose
consciousness the identity of both moments became visible; but where he now is,
what does that concern us? It is the human race that is that θεάνθρωπος; this
is really present in every human being, --- if only it comes to a revelation!
We gladly bid farewell, then, to the vanishing individual, Christ, in order to
come into possession of the eternal, present, and thus in truth historical
Christ\footnote{\ldots{}
durch die Belebung der Idee der Menschheit in sich, namentlich nach dem Momente,
daß die Negation der Natürlichkeit und Sinnlichkeit, welches selbst Negation des
Geistes ist, also die Negation der Negation, der einzige Weg zum wahren
geistigen Leben für den Menschen sei, wird auch der Einzelne des
gottmenschlichen Lebens der Gattung theilhaftig. (With these words Strauß
refutes the objection that Schaller raises against him in his work: der
historische Christus und die Philosophie. S.~\textit{64}.) Diess allein ist der
absolute Inhalt der Christologie: daß dieser an die Person und Geschichte eines
Einzelnen geknüpft erscheint, gehört nur zur geschichtlichen Form derselben.
Strauß, Schlußabhandlung L.~J. §~\textit{159} Pag.~\textit{768}.}. For what,
after all, is true
% --- p. 138 ---
\opage{138} history but the eternal presence of actuality? Why then should one
covet a subjective immortality? Whoever seeks his life shall lose it! Such
egoism was not found in the Christ whose praise you yourself so often have on
your lips. He sacrificed his subjective personality in order to rise in the
consciousness of others and ascend to the heaven of a grateful remembrance. ---
In Christ, then, we shall live, die, and rise!

We shall now test somewhat more closely that immanence of ideality and reality
on which the whole interpretive method of this view of religion is built. An
abstract ideality requires that all the moments which by the power of negation
mutually exclude one another must at the same time have this negation so
thoroughly reflected within their own limits that position and negation
everywhere coincide. But the truth of infinite determination and limitation
depends on the actuality of negation. Now abstract ideality contains only an
abstract identity with itself; its true opposite it keeps outside itself.
Therefore its moments everywhere lack, within its limits, a true opposition; the
ideal does seem to be limited by the ideal, negation by negation, but this
limitation does not actually take place. Since abstract ideality thus has its
true negative outside itself, which conflicts with the nature of ideality, it
vanishes. Only the realm of external negation, of empirical actuality, remains,
built on the ruins of ideality. --- Here it is the first principle that positive
things exclude one another; this mutual exclusion is indeed at first thought of
as such that each single thing rests undisturbed within its rightful limits;
but since the thing's limit is itself a thing, and the thing thus becomes its
own opposite: there arises an incessant opposition without rest, an incessant
% --- p. 139 ---
\opage{139} drive to alteration and annihilation pervades the marrow of the
thing, the realm of abstract reality annihilates itself. But as reality, through
this destroying, carries through the negation of itself, the relation of
ideality returns again. Since, then, neither pure ideality nor abstract reality
is anything in itself, it naturally follows that this restless oscillation
between ideality and reality is not only the beginning of, but also the
principle for, our existence in this life. Beginning and principle coincide.
Since nothing can emerge in the conclusion that was not already in the
beginning, it is no wonder that the fluctuation goes on to infinity. Ideality
does always seek to find its existence in the realm of reality; but the
individual moments of empirical knowledge are perishable; therefore all its
attempts are in vain; --- were ideality true and perfect in itself, it would
also have its subsistence in itself; it would then also be able to gather
together the scattered moments of reality and reveal itself in the empirical
sphere; but as it is, reality remains corruptible through the impotence of
ideality, and ideality remains impotent through the corruptibility of reality.

It is from this self-dissolution that that θεάνθρωπος suffers which is to take
Christ's place. In the God who dwells in the sphere of ideality the human race
sees its Father, but this God is not a subject established in itself,
unshakable, not Jehovah, who in his unshakable holding fast to himself has all
his opposites outside himself, but it is the Father who seeks his existence in
eternally begetting the Son, i.e.\ the human race. But what existence, now,
does this Son of God have? He is not pure ideality --- for this belongs to the
Father alone ---, nor abstract reality, which unfolds itself through individual
things and individual human beings; --- he is the very transition from the
Father's ideality to the world's reality,
% --- p. 140 ---
\opage{140} the Mediator (\textit{mediator}) between God and world. We must then
test this mediation more closely.

One must distinguish between the general consciousness of the human race and
the individual consciousnesses of individual human beings. Christ is this
general consciousness --- not that eternal individual I which once revealed
itself on earth and now reigns in heaven (for then one would surely have to take
refuge again in the Jesus born in Bethlehem and crucified on Golgotha, whom
speculative immanence has long since made antiquated!) ---, but the general
consciousness which procures itself existence through the individual
consciousnesses, so that, for the same reason, namely because it does not want
to be anything in and for itself but only to have everything through the
individuals, it demands that the individuals must not want to be anything in
themselves, but that they should give up their subjectivities and return into
it. The development of this general consciousness through the individual
consciousnesses, --- the historical operation by which the discord between the
human race and individual human beings is more and more sublated, is the Spirit,
--- yet not the Spirit that, as the third I, unites the life of the Father and
of the Son, leads God's negativity back to being an immanent negativity within
the limits of the divine life, and thereby can at the same time allow the
individual consciousnesses to have their own negative, true ideality, eternal
life, in themselves, --- but the spirit that imparts to the actuality of the
moment all the ideal treasures of past ages and sincerely declares that it has
its existence only in this immanence of ideality and reality. But the momentary
actuality that has unfolded itself through the individual consciousnesses
already shows itself to have become another (altered); the objective spirit is
thus limited by the sum of finite consciousnesses. For considered according to
its real existence, human consciousness has other
% --- p. 141 ---
\opage{141} consciousnesses posited outside itself, so that there is no ideality
present in their sum; considered ideally, on the other hand, the one
consciousness takes up the other's life into itself. The single real
self-consciousness thus takes up into itself the spiritual fullness of the other
consciousnesses. It is therefore in the individuals that one must seek the
immanence of ideality and reality. But this immanence continually dissolves into
scattered moments; for insofar as the really existing consciousness thinks, it
also has the immanence in itself, but insofar as it really exists, the
scattering sets in. For the ideality that is present in real consciousness is
only entrusted to it for a time; thinking consciousness does indeed recognize
that the fullness of this ideality is infinite, but it will never be able to
exhaust it completely. For since the real self-consciousness, which, as it is
said, has its negation outside itself but not within itself, must die, ideality
continually seeks out another real self-consciousness in order through this one
to preserve its identity with itself. But since every real self-consciousness is
subject to the same fate, contradiction and opposition penetrate right to the
innermost core of this spirit. For real self-consciousness remains the necessary
organ without which spirit cannot exist, but exposed as it is to infinite
alteration it will never be able to take spirit up into itself; this
incommensurable relation will then persist in all eternity.

Since the matter stands thus, one cannot confuse the panlogistic Trinity with
that of the Church. The God of panlogism is not our Father,
ὁ μόνος ἔχων ἀθανασίαν, but an intoxicated God who cannot collect his scattered
thoughts and reels drunkenly about among the scattered fragments of existence.
Its Son of God is not the eternal Son, who declares of himself that he had glory
with the Father before the world was, but the son whom the Father begot, spurred
on by the contradiction of ideality
% --- p. 142 ---
\opage{142} and reality, --- who brings to nothing the hope of the Father
grasping after existence, who mocks the nakedness of his abstract ideality and
is therefore followed by the Father's curse. We know this spirit too, which is
sent forth by such a Father and such a Son; it is the spirit of this world,
which everywhere betrays the σκάνδαλον of the Father and of the Son. Its
immanence is a contradiction, its reconciliation a horror. I, it says, wipe away
all tears and still all pains. For in order to be seized by pain one must have
the positive fullness of life in oneself, a fullness of life that suffers under
a contradiction penetrating to the innermost; but were life completely negated,
and had it only its negative outside itself, a sudden death would set in, and
then there would be no pain; the painful life, dissolved and scattered in its
innermost, thus remains equal to itself, and thus identity is preserved. Pain is
the innermost root of human existence; for as soon as the human being gathers
all his powers in order to intuit his own nature in complete clarity, ideality
and reality coincide: there is a true identity of them, for precisely this
constitutes life itself, but there is at the same time, as we have seen, a
fundamental contradiction between them, and herewith the dissolution of life,
and pain, are given. In this, then, consists sin, that the human being is
carried away by love of life; but love of life begets fear of death, and this
is precisely the punishment for the former; if, then, we are seized by both, we
cannot doubt that sin and punishment interpenetrate, and that reconciliation
consists precisely in this pain. Pain is the most perfect immanence of
contradiction and identity, it is precisely the eternal spur to striving after
eternal life, it is God's holiest mystery; pain is blessedness; --- but it is a
blessedness that recalls the Devil's roaring laughter! Human self-consciousness
% --- p. 143 ---
\opage{143}\setcounter{footnote}{0}%
hears it and shudders; for the more attentively it regards the elements of this
realm, the more sharply and clearly it sees the contradiction of pain, in a
thousand different shapes, stream through all its individual parts. But when in
this hell of darkness it hears the Church sing the praise of its Christ, who was
crucified but rose from the grave and sits at the Father's right hand, behold!
then he rises again, like the man who wakes from a dreadful dream and gladly
greets the light of day; what draws him is not a calculating egoism that knows
only how to seek its own, no, it is the objective power to which he gives
himself with all his heart, the true longing for the true God, and he is no
longer carried away by a blind passion, once he has thus experienced the bitter
pain of the dialectical contradiction. Get thee hence, Satan, thou who deniest
the holy! Unto us come the negation of the negation!

\parag{22}{Return to the Church. Speculative Theology.}

\begin{verse}
``Let hissing sound in serpents' yard\\
That I am now laid waste;\\
I crown my year no less for that\\
With fruitfulness and harvest;\\
Yes! up I rise like ear in field,\\
Like May in beechen woods,\\
And splendid, amid song of birds,\\
Like golden sun from wave!''\footnote{Grundtvig: Den christelige Tro.}
\end{verse}

These words, with which the poet sings the victory of the Christian faith amid
the various vicissitudes of the ages, resound not only in the life of the Church
but also in science. Human reason moves in two opposite directions, from God to
the world and to itself, from itself and the world
% --- p. 144 ---
\opage{144} \setcounter{footnote}{0}%
to God. It was not, indeed, to be expected that this transition should take
place calmly and without interruption in our sinful and finite world; but we
nevertheless always see that reason, although it often seems to turn altogether
away from the truth of Christ, repents of this apostasy again once it has fully
tasted its sorrowful consequences. --- One must undeniably recognize panlogism
as the highest stage of self-cognition to which abstract reason can attain, and
one cannot rightly reproach it with unwillingness to return to positive life;
but is it also the right way by which it returns to this? If there were no
world order that could be negated, the human I would never be able to arrive at
a development of the logical concepts either. But from this it follows that the
system of concepts that panlogism sets up is only a formal one, that it
everywhere lacks positive actuality and is at every point filled with a need, a
longing, for this. And panlogism itself fully recognizes this. But it is
precisely its πρῶτον ψεῦδος that it ascribes divine authority to this ideality
and takes its impotence for omnipotence. For this ideality is nothing other than
the shape of human reason itself. Yet the error lies not so much in its
assuming an identity of divine and human reason as in its declaring the
immanence of both to be immediately given in the latter. For an identity that
does not rest upon an opposition is, according to the testimony of speculative
logic, just as surely unstable and bad as a commingling of divine and human
reason is, according to the testimony of the Church\footnote{How forcefully and
acutely J.~H. Fichte has combated this rationalism is well enough known.}, the
true mark of rationalism.
% --- p. 145 ---
\opage{145} %
Panlogism has a false God; it is therefore illogical (ἄλογος) if it does not go
beyond itself; but it has a true logic and therefore needs only to take refuge
in Christ in order to find a God of a truly logical nature. It itself declares
that the way to immanence passes only through the greatest oppositions, but its
own oppositions are of an all too light nature; if, therefore, it comes to a
complete cognition of the logical principle that nothing can emerge in the
conclusion which did not in actuality lie in the principle, it is also compelled
to recognize that the opposition which shows itself in this world between the
ideal and the real has its presupposition and ground in an immanence of both
mediated from eternity, in other words, in a divine, uncreated reason,
\emph{qualitatively} different, as it is said, from our human reason. But once
God's absolute freedom has thus been recognized, the whole relation between the
positive moments and the logical categories is reversed; we then see that it is
a positive divine power that has founded the whole of real existence and
precisely thereby posited the ontological forms as immanent in things. We are
then referred to the whole fullness of actual life, in which all figures must be
considered as they meet us in life; and thus the moment of faith and of view
again asserts its right. We then consider the human community in which we find
ourselves, and through which the idea of eternal personality shines; we look to
the divine founder of this community, who came forth in the fullness of time,
when the law that had been given to the Jews had developed the consciousness of
sin; we then recognize that God, by this way of revelation, has come forth as
the Triune, and so on. It is contrary to reason to want to call some of these
elements into doubt, with the thought of arriving, by means of other elements,
at a dialectical necessity of that
% --- p. 146 ---
\opage{146} \setcounter{footnote}{0}%
which is the object of doubt; for if one really entertains doubt about one
moment or another, one must also necessarily think the remaining positive
moments with which it is connected in such a way that the doubt about the former
lies hidden in their concept. It therefore helps but little against the deist to
appeal to the revelation of Christ, since he has a one-sided concept of God. Nor,
therefore, is theology to prove that it follows from God's Trinity that Christ,
with dialectical necessity, had to reveal himself in this or that place; but
since he has once revealed himself, theology must be able to conceive how this
positive phenomenon of revelation hangs together with the other positive
phenomena; this revelation is, as it were, a real thought that must be capable of
being reconciled with the other real thoughts. The absolute presupposition of
religious conceiving is therefore absolute faith, and so far is science from
being injured by this inviolability of faith that it is rather only by its help
that dialectical necessity and negation can penetrate all the elements of the
system. As, therefore, abstract logic passes over into the positive elements, so
on the other hand must the positive elements necessarily pass over into a
concrete metaphysics and speculative empirical knowledge. When the objects are
put together in an external manner, it is all over with science, and it has
therefore almost become a common proverb that every philosophy ought to begin
with doubting, although one ought rather to say: ``with negating''; for
doubt\footnote{That an untruth lies hidden beneath this doubt, the philosophers
have already acknowledged. Thus J.~L. Heiberg: All true philosophy must unite both
the mottoes cited by Mr. Martensen: the fear of God it must have in its heart,
doubt in its eye. Perseus Nr. \textit{1}. P. \textit{41}.} stands in opposition
% --- p. 147 ---
\opage{147} \setcounter{footnote}{0}%
to faith, but that negation and abstraction do not injure it we have already
shown earlier. As, therefore, the positive moment of objectivity is absolutely
preponderant in faith, so through negation the absolute freedom of subjectivity
returns, for all moments must be negated and tested in such a way that not a
single phenomenon is left untouched\footnote{It is this testing that Sibbern
requires when he demands: ``an all-embracing debate of all against all.''}. What
is cited as a merit of Hegel's logic, namely the great acuteness with which,
through a thorough negation and contradiction, it sharpens the dialectical spur,
we shall find the same also in dogmatics, if only we take the right path.

The mutual opposition of the individual moments has far greater depth in
Christian dogmatics than in panlogism, and it is precisely thereby that the
former leads to a true identity and reconciliation. Precisely because panlogism,
in order to reach immanence, divided existence into two abstract parts, precisely
for that reason it could not reach identity; Christian dogmatics, on the other
hand, posits an absolute opposition between the truly triune God, complete in
himself, and the created world; but through this fundamental contradiction a true
reconciliation can come about between the peace of eternity and the struggles of
temporality, since the knots that are tangled here are already untied there, in
the principle itself.

In order to unite and reconcile all the ideas, panlogism sacrifices the
individual self-consciousnesses and precisely thereby moves in an eternal
opposition between reality and ideality; Christian dogmatics, on the other hand,
ascribes to the individual self-consciousnesses
% --- p. 148 ---
\opage{148} an inner infinite immanence in themselves, according to which they
die, to be sure, \textit{ad extra} (ὁ ἔξω ἄνθρωπος), but nonetheless \textit{ad
intra} (ὁ ἔσω ἄνθρωπος) preserve themselves for all eternity. But while it thus,
as it were atomistically, separates them all from one another, gives the nature
of sin the deepest development, and unshakably holds fast the ethical relation of
opposition between God and man, it also penetrates to a true immanence of all
these different moments, in that it develops the deepest concept of freedom, of
conversion, and of love.

When panlogism would transfer to the whole human race what the Church ascribes
to its Christ, without remembering on the one hand the inner, infinite fullness
that is the essence of human self-consciousness, and on the other the mutual
relation between Christ and the Holy Spirit, it mixes together in pitiable
confusion moments that are really opposed to one another. Christian dogmatics,
by contrast, ascribes to its Christ a perfect and absolute identity of God and
man, and so far are the equal right of the ideas and the identity of all from
being disturbed by this, that rather the incorruptibility of life becomes the
portion of all through the One, since the individual members of the body hang
together organically with the head.

When, finally, panlogism, without knowing the true and actual mediation of the
cosmic system, constantly has the word mediation in its mouth and assumes that
the whole dialectical course of development has in actuality always been the
same, as if the relation that obtains in our time between nature and subjective
spirit were eternal, it seems to contradict itself in a remarkable way when it
nevertheless promises a time when spirit and nature shall infinitely penetrate
each other. Christian dogmatics, by contrast, although it proceeds from a true
metaphysical identity
% --- p. 149 ---
\opage{149} of the divine will and the laws of nature, presents objective
mediation in such a way that God on the one hand has made nature the foundation
of subjective spirit and therefore also lets his ideas develop through the
objective forms of nature, and on the other hand also transfigures nature in a
peculiar manner by his presence, since sin had darkened human self-consciousness
and given the worldly and profane moment an isolated development, from which it
then follows that at times a sharp and actual opposition appears between the
general divine governance and supernatural revelation, --- a doctrine of
mediation that contains a true presentation of God's essential accommodation to
the peculiar character of different times and men.

Human reason does not shun these actual oppositions and fundamental
contradictions; for it recognizes that they all have their origin from, and
their existence in, itself. As, therefore, the I of the regenerate man remains
the same although an entirely new life is created in it, so too, far from the
oppositions that penetrate and, as it were, dissolve human reason abolishing its
identity with itself, it rather overcomes the contradiction of pain and attains
reconciliation in proportion as it more and more immerses itself in and
appropriates these oppositions. From this it becomes evident that created reason
is in actuality determined and conceived through the uncreated, and that the
mediation of the supernatural and natural moments takes place precisely through
the absolute opposition; indeed, one can truly say that every sacred history,
that is, every opposition between the pragmatic development of the world and the
miracles, has precisely the task of leading to a reconciliation of those opposed
moments.

% --- p. 150 ---
\opage{150} \parag{23}{Miracle and Sacred History.}

As it is with revelation, so it is also with sacred history; the truth of
revelation depends on the actuality of miracle, but this is developed through
the cosmic \emph{mediation} of things. For the inquiry into the actuality of
miracle does not aim at showing that God's will has influence on nature; for
there is an actual connection between God's activity and the forces of nature,
and every time we move hand or foot we ourselves give an example of how thought
has influence on the body; --- but its aim is, as we have already remarked
before, to demonstrate an immediate suspension of the opposition that exists
between the divine counsels and the laws of nature.

All miracles thus seem contrary to reason, because God always works everywhere
equally mediately and immediately; but such a presentation of the metaphysical
relation is just as unsatisfactory as it would be if someone thought he had
understood the concept of human freedom because he had observed that God works
everything in man and that man is at the same time self-active. All such
fundamental relations must necessarily unfold through the oppositions of
temporality. Thus it becomes impossible to present a thorough mediation of
divine and human relations without granting miracle a definite place.

Deism, which above all insists on a natural world order, is nevertheless itself
compelled to appeal to one miracle as the starting point of its whole pragmatism,
namely creation. The same holds of panlogism; for while it has discovered the
moment of illusion in the Christian
% --- p. 151 ---
\opage{151} \setcounter{footnote}{0}%
traditions, so that these no longer present God to us as immediately present, it
nevertheless also concedes that objective spirit has been able to reach the
standpoint of enlightenment that distinguishes our time solely and only through
this standpoint of illusion. What is presented mediately in philosophy was thus
immediately present in those narratives, so that the separation between ordinary
events and the phenomena of revelation returns here again. The development of
all things thus has the aim of leading to a mutual immanence of nature and spirit
at every point, so that the \emph{mediate} and the \emph{immediate} may
everywhere coincide; but if this harmony is to become actual, the elements must
first separate from one another, so that nature and spirit on the one hand each
have their own peculiar points where they work in isolation, in order on the
other hand, precisely by reason of this separation, to be able to penetrate each
other mutually at other points.\footnote{An example of this mediate and
immediate relation is offered by daily life itself: external nature is, namely,
transformed more and more by the human spirit, but in a \emph{mediate} manner.
Therefore Strauß rightly says: .\,.\,.\,.\ ``sie (die Menschheit) ist der
Wunderthäter: sofern im Verlauf der Menschengeschichte der Geist sich immer
vollständiger der Natur, im Menschen wie außer demselben, bemächtigt, diese ihm
gegenüber zum machtlosen Material seiner Thätigkeit heruntergesetzt wird. J.~L.
Schlußabhandl. § \textit{149} Pag. \textit{767}.

But this mediate activity is bound up with the immediate and real connection that
obtains between soul and body; and if one would now call this mythical, one would
also have to call all narratives of such works myths.}

If, therefore, one considers the miracles in isolation, they seem contrary to
reason; but the reason for this is precisely that the immediate is always
determined by the mediate, just as the latter in turn is determined only by the
former. If, on the other hand, they are considered in the general
% --- p. 152 ---
\opage{152} \setcounter{footnote}{0}%
light of all that goes before them and follows after them, they are of a logical
nature and can consequently also be conceived. Miracle can therefore be called
an anticipation of the relation which, according to the religious spirit, already
obtains, metaphysically taken, and shall one day be seen realized, between
created nature and the Creator. The God who appears in the miracle is therefore
the same who is present in religious consciousness. Miracle is thus the very root
of religion (we are speaking, namely, of the miracle that is the most important
in each religion); it contains the peculiar moment of view whose inner essence
the spirit of religion is more and more to unfold. Since miracle thus
constitutes the very content of the spirit, the religious spirit itself must also
waver as soon as doubt arises about the actuality of the miracle, and if it is
discovered that the miracle is untrue, the spirit itself must as a consequence
vanish and bow beneath a mightier spirit that is mediated through a truer
actuality. That the panlogistic spirit cannot reconcile itself to the historical
Christ\footnote{Cf. Strauß L.~J. Pag. \textit{766}.} therefore has its natural
ground in the fact that it is not mediated through Christ; --- if one considers
its three main objections, this will easily become evident. For it denies, first,
the possibility of an individual union of God and man, because it does not know
the inner fullness of the individual consciousnesses; but the reason for this
ignorance we have already stated above. It remarks, next, that it would conflict
with the nature of the Idea as such to enrich and exalt one particular organ at
the expense of the other organs; but here it does not take due account of the
real moment of mediation (cf. § \textit{22}). And when it finally objects that
Christ, as an individual phenomenon, is neither more nor
% --- p. 153 ---
\opage{153} \setcounter{footnote}{0}%
less than what presents itself to us in his phenomenal
existence\footnote{Allein, wie Schmid gründlich nachgewiesen hat, ein
geschichtliches Individuum ist eben nur das, was von ihm erscheint; sein inneres
Wesen wird in seinen Reden und Handlungen erkannt; zu seiner Eigenthümlichkeit
gehört die Bedingtheit durch Zeit= und Volksverhältnisse mit, und was hinter
dieser Erscheinung als An sich zurückliegt, ist nicht das Wesen dieses
Individuums, sondern die algemeine menschliche Natur überhaupt, welche in den
Einzelnen, durch Individnalitet, Zeit und Umstände beschränckt, zur Wirklichkeit
kommt. Strauß L.~J. Pag. \textit{740}}, we remark on this that existence cannot
be considered from its sensuous side alone, --- for from that one would merely
extract the contradiction that the sensuous is not the spiritual; if, on the
other hand, it is understood to mean that there lies nothing else in the
phenomenon than what can be derived by a negation of the phenomenon, then
precisely a spiritual interpretation becomes necessary; but what extent this is
to have, the Spirit alone can teach us. There thus arises an intimate connection
between Christ and the Holy Spirit: what lies \textit{implicite} in Christ, the
Spirit sets forth \textit{explicite}. The Spirit would thus not have presented
the idea of eternal personality to us in the Church, had not the infinite
fullness of the Godhead (πᾶν τό πλήρωμα τῆς θεότητος) dwelt in Christ. Therefore
Christ nowhere says that the Holy Spirit will one day come of its own prompting
and teach that he himself was only a vanishing individual; but he promises that
the Spirit shall come from the Father in his name, and adds: ``When the Spirit of
truth comes, he shall guide you into truth. For he shall not speak of his own,
but what he hears, that shall he speak, and the things to come he shall proclaim
to you. He shall glorify me; for he shall receive of mine and proclaim it to you.
All that the Father has is mine. Therefore I said that he shall take of mine and
% --- p. 154 ---
\opage{154} proclaim it to you.'' (John \textit{16, 13---15}). There are
therefore two fundamental miracles in the Christian religion, the
\emph{Incarnation} and the \emph{Inspiration}.

When the disciples at first gathered around Christ, already then it was in truth
the Spirit that unconsciously drew them, but only in that it shone forth
immediately from the Savior's own person; therefore Philip, with the words ``Come
and see!'', calls on Nathanael, who doubts that the Messiah can come from
Nazareth, to follow him (John \textit{1, 47}). Yet only when the apostles, after
having received the necessary religious preparatory formation in their life
together with Christ and through his oral teaching, had seen his death, --- when
he, glorified after the resurrection, had stood before their beholding gaze as
it were on the boundary between the heavenly and the earthly kingdom, presenting
to them in his essence the individual and the universal in true unity, --- when
he finally, lifted up to heaven after giving them the promise of the Comforter
who was to come, had left them for a time to their own reflections: only then was
it that there revealed itself immediately the sending of the Spirit in the
miracle of Pentecost, prepared by a true mediation within them. Christ's external
phenomenon was negated, and the interpreter who already lay hidden in their
consciousness came forth suddenly into the full light of consciousness. As,
therefore, Christ had previously been revealed in his glory in the external
world, so now for the first time he was present in perfect glory in the inner
kingdom of consciousness. This relation between Christ's immediate presence and
the immediate activity of the Spirit sent by him gives precisely the apostolic
age its peculiar character. If, therefore, one considers Christ's life in its
entirety, one will easily find that mutual relation again. For just as his
public life belongs most nearly to the phenomenal actuality of history, so that
the moment of view is here
% --- p. 155 ---
\opage{155} preponderant, but the true understanding and development of it is
first given by the Spirit, so his glory with the Father reveals itself most
nearly to the searching spirit, but the truth of the appropriation of this
revelation is in turn conditioned by the believing reception of the teaching
about his life here on earth. Christ's entry into the world and his exit from it
are thus forms of revelation that by the nature of the case hover between
phenomenal actuality and spiritual ideality: anyone who does not believe in the
Spirit mistreats history.

From this it is also evident that the reconciliation of the world, too, must be
referred to Christ's sacred history, rightly understood. For the reconciliation
has been accomplished in that Christ underwent the harshest oppositions of life
and came forth victorious from the struggle. Just as the natural man enters this
world in such a way that he has darkness behind him, so that he does not know
whence he comes, and just as he must pass through the hard trial of the struggle
against all the contradictions of life and death without knowing whither he
goes, so conversely Christ, who is the way and the truth and the life, goes forth
from the kingdom of eternal light, reveals God's presence in the world of pain,
while his divine spirit always remains like itself, and wanders through the
darkness of death back to the kingdom of glory. (``I came forth from the Father
and am come into the world; I leave the world again and go to the Father.''). He
thus not only shows individual men the way, but also reveals the true history of
the whole world.

But just as the miracle of inspiration is mediated through the individual Christ,
so the apostolic age also prepares the subsequent formation of the Church.

For it presents to us, on the one hand, a fixed and
% --- p. 156 ---
\opage{156} unchangeable image of Christ, so that Scripture utters its ``Come and
see!'' to anyone who longs to see Christ; but on the other hand it leads beyond
itself to another form of the Church. All the institutions of the Church aim,
namely, at developing the idea of freedom more and more, so that human
consciousness, reflected into itself, may through the cognition of its own nature
seek, find, and thoroughly appropriate God. That our time cannot show any miracle
is thus so far from grounding any suspicion against the truth of the miracles of
the apostolic age that we must rather turn the matter around and say that the
miracles of that age precisely explain to us the reason why no miracles take
place in our time.

The manifold fluctuation that takes place in the Church in the course of the ages
between real objectivity and abstract ideality, as they alternately succeed each
other and pass over into each other, serves to reconcile the concept and the view
with each other ever more truly and inwardly. If, therefore, one of these moments
becomes preponderant at one time or another, we always see the other emerge and
moderate it.

For the dialectical development always proceeds in such a way that on the one
hand spirit is formed through a negation of the phenomena of life, and on the
other these are developed and explained by it; as, therefore, the spirit is, so
also is the explication and interpretation of the phenomena. Therefore the
Apostle's words: ``Beloved, believe not every spirit, but try the spirits
whether they are of God!'' sound constantly through the Church. But that which
tries the spirit is precisely the spirit itself, for of the spirit it holds:
``\textit{verum est index sui et falsi}''. Therefore the spirit of abstract
freedom and personality will always come forth victorious from the trial. ---
It will then be clear from this that the principle in sacred history is God's
true, real accommodation. God, conscious from eternity
% --- p. 157 ---
\opage{157} %
of his absolute freedom, formed the counsel to develop a human freedom and
personality; therefore all such phenomena as men cannot but take for definite
manifestations of divine assistance, aimed at the realization of that counsel,
must also necessarily be taken as having in actuality belonged to the divine
counsels. If, therefore, the sacred phenomena are considered in the light of such
an accommodation, one must also ascribe actuality to them in the deepest sense of
this word. For if such phenomena are considered in and for themselves, they are
by their nature adequate expressions neither of the things that belong to the
temporal life nor of those that belong to the eternal. For insofar as the
supernatural things are revealed, they must also find expression in phenomena
that belong to our empirical knowledge, and thus put on forms that are
incongruous with their own nature; and although they descend into our sphere,
the phenomena in which they are revealed must nonetheless be glorified in a
peculiar way so as not to be confused with other, ordinary events, or --- the
descent of the essence into the lower sphere must at the same time be the
elevation of the phenomenon into the higher sphere. Therefore illusion, too, is
always a necessary moment in every temporal phenomenon of revelation. --- The
actuality of the sacred phenomena therefore always presupposes a peculiar
governance by God. When men see a peculiar and wondrous phenomenon (τὸ τέρας),
they conclude from it that God will make known to them his immediate presence
(τὸ σημεῖον); if this conclusion is false, even the most exalted phenomenon
remains without effect; if, on the other hand, it holds as true, the hovering,
double nature of the phenomenon puts no obstacle in the way of the working of
revelation. For it is God's supernatural power that gives the phenomenon its
peculiar stamp; and that it does not fully express the essence
% --- p. 158 ---
\opage{158}
itself stems from his accommodating himself to man's power of comprehension. It
is thus without significance for the truth of the miracle whether the angel
Gabriel revealed himself to Mary as it were at one stroke, or whether she,
through the cooperation of God's Spirit, only later saw presentiments that from
the beginning had merely stirred in the innermost depths of the virginal soul
rendered visible in complete clarity. For Gabriel, invisible by nature, becomes
visible at God's command, in order to give us a representation of Christ's
supernatural conception. Gabriel is only a serving instrument in God's hand. ---
The same holds also of the cloud that, according to the account of Scripture,
surrounded Christ when he ascended into heaven. For the cloud formed a visible
transition from the earthly to the heavenly kingdom. One can then say with as
much right that it existed as that it did not exist; it is therefore indifferent
with regard to the truth of the miracle whether it consisted in the
transformation taking place visibly in the external world, or whether God's
Spirit merely presented the phenomenon in this form to the thoughts of the
apostles, provided only that it stands firm that God has been a truthful and not
a misleading spirit. One can thus quite well hold different opinions about the
origin and formation of the phenomena and nonetheless have the most unshakable
faith in their full actuality. But this difference of opinion is not on that
account dependent on our subjective discretion; for we ought to take account of
vulgar empirical knowledge \textit{ad extra} just as much as of God's actual and
true accommodation \textit{ad intra}.

The older Church insisted one-sidedly on God's objective activity, but thus did
not give human freedom its full due. For the miracle of inspiration was developed
in a sensualistic manner, if we may say so, so that the activity of the Holy
Spirit was placed not so much in the interpretation of the
% --- p. 159 ---
\opage{159} %
events that had taken place as in the faithful recollection of them; as,
therefore, the Spirit had suppressed the self-activity of the apostles, so now
also ecclesiastical faith sought to bring reason into bondage. Criticism has
taught us that revelation did not take place in so slavish a manner. For just as
the religious relation between God and men first becomes true through the moral
exertion of the latter, so also inspiration has freed the consciousness of the
apostles in such a way that they could pass a freer judgment on the events that
had taken place. And so far is the truth of revelation from being weakened by the
discrepancy between the individual moments of sacred history that it is rather
strengthened by it. For had Christ not descended from his glory with the Father
into the realm of everyday experience, that is, into the frailty of our
existence, he would not have been a true man (υἱός τοῦ ἀνθρώπου); he would then
have remained standing in the clouds, like the Son of Man in the prophet Daniel's
vision (\emph{Bar enash}), and we should then have to call his earthly life an
apparent existence (δοκετίσμος). But by thus appearing in the realm of phenomena
he has given the spirit freedom to gather the scattered moments of experience
into the Idea's own clear, intuitable totality. Here, then, the old maxim also
holds, that disagreement in what is inessential confirms agreement in what is
essential, so that we have before us not an ingenious fiction but an actual
revelation. And just as the Spirit asserted its freedom to create the forms ---
so that the same divine honor belongs to the Holy Spirit as to the Son ---, so it
also granted the apostles and the apostolic age freedom to present the same truth
in different words. It regenerated the apostles, to be sure, but did not
sanctify them at one stroke, so that every manifestation of the frailty of the
human will
% --- p. 160 ---
\opage{160} %
suddenly became impossible for them; it enlightened their spirit, to be sure,
but not in such a way that it suddenly removed every subjective one-sidedness or
error from them. --- Scripture's depiction of the sacred phenomena therefore
contains at once the divine truth in its highest actuality and the freest,
purely subjective movements in the consciousness of the sacred writers; the
sacred phenomena are not either the one or the other, but are precisely only in
the intimate mutual transition of both these moments into each other, and cannot
be determined by any external boundary whatsoever.

But a revelation and tradition for whose whole economy freedom is the principle
naturally cannot be appropriated before the true method of freedom has been
found. Now since the older Church and criticism sought to appropriate the
supernatural idea by holding fast the phenomena by themselves, it followed as a
matter of course that they could only infer by analogy and careful analysis of
the individual moments. From this it followed that on the one hand the older
faith had to insist on the verbal agreement of the narratives, and on the other
criticism had more and more to extend its doubt to the miracles in general as
soon as it had discovered discrepancy in particulars; \emph{for if the lesser
miracles cannot stand the test of inquiry, how much less can the greater.} But
this positivism and probabilism is abolished at the root when, through the
negation of the phenomena, we penetrate to the Idea itself: we then also hold
unshakably fast to the fundamental miracles, in order to be able to consider the
more special miracles in their light. The restoration of the phenomenon of Christ
must then take place in such a way that the individual narratives of his life
each keep their place, in order thus each in its own way to reveal the idea of
the Incarnation.
% --- p. 161 ---
\opage{161} %
We thus return again to the ecclesiastical image of Christ, but in such a way
that, after having subjected its individual traits to a separate analysis, we
see that the same Christ whom the Church believed it could behold in the
external realm of everyday experience must not be considered outside the mirror
of the Spirit.

Inspiration as such ceased, to be sure, but the Holy Spirit continues its life in
the Church: amid all the vicissitudes of time it fights with might against the
spirit of this world; whenever sacred history is distorted, it is the Spirit that
anew refreshes, vindicates, and by its divine light transfigures the Church's
divine image of Christ; and it will continue to illuminate it until the moments
of the concept and of the view fully penetrate each other; --- until Christ is
revealed in his glory and his voice sounds to all: ``Come and see!''

\par\vspace{1.0\baselineskip}\centerline{\rule{0.25\textwidth}{0.5pt}}

\division{Conclusion.}

If one grants the truth of the development we have thus given, one will also
agree with us in taking the speculative treatment of sacred history to be the
only true one.

On the one hand it preserves the positive character, for in the recognition that
objective reason cannot be conceived except through a careful observation of the
given objects, it receives with joy the treasures that the learned have gathered,
and leaves none of the investigations of criticism unheeded. But on the other
hand speculative theology well knows how to distinguish between the apparatus of
science and science itself; therefore it follows with care
% --- p. 162 ---
\opage{162}
all philosophical investigations; not in blind adherence to this or that system,
but in order, as the teacher of abstract philosophy, after careful testing of
each of its systems, to lead created reason back to its true dependence on
uncreated reason, and through this reconciliation of both to interpret the
phenomena of religion. --- It is far from all craving for novelty; for its
endeavor aims precisely at bringing to light all the true moments that are found
in the older systems. But above all it demands a true transition from one sphere
of cognition to the other and an inner dialectical connection between the
individual parts; it therefore rejects every mixing together of opposed elements,
disdains all the artifices of eclecticism, and flees the mediocrity of the bad
middle way.
\end{document}
